% SOURCE docs/release/2026-10-29/source/docs/writing-repair/2026-10-29/src/67eb5baa2057b943.tex CHAR 0
\chapter{Language assistance, causal discovery and uncertainty}
\label{ch:la03-discovery}

This chapter follows a graph suggestion through the operations that could
make it useful for a decision. The opening comparison separates linguistic
knowledge from observations and interventions. The subsequent studies develop
information bounds, numerical alternatives, search failures, expert evidence
and executable graph contracts. Temporal and latent-variable examples then
show what those operations cannot determine from the supplied information.
The final identification study asks whether the resulting answer still
addresses the original population and action. These are successive obligations,
not interchangeable measures of discovery quality.

\section{The scientific object behind a graph proposal}

A banking analyst may begin with a list of columns and a request to explain why
balances move after a change in deposit pricing. A language model can recognize
the vocabulary, propose a market-rate variable, distinguish a customer offer
from a realized rate, and suggest that customer characteristics affect both
assignment and retention. These suggestions are useful because a statistical
algorithm does not recover the meaning of a column from its numerical values.
The scientific question nevertheless remains open. The proposed relation may
describe a mechanism in another population, a historical association, a direct
effect, an indirect path, or simply a familiar phrase. The role given to the
suggestion determines the claim the resulting graph can support.

The inspected literature offers several roles rather than one method called
``LLM causal discovery.'' Pairwise language judgments, breadth-first graph
construction, retrieval-assisted discovery, response-based independence
queries, statistical prompting, soft graph priors, multi-agent debate and
numerical experiment design alter different parts of an analysis
\cite{v12_la_2305_00050,v12_la_2402_01207,v12_la_2402_15301,
v12_la_2406_07378,v12_la_2402_01454,v12_la_2405_13551,
v12_la_2407_15073,v12_la_2608_09696}.
They should be compared by what information they consume, what object they
change, and what uncertainty survives that change.

Let $D$ denote numerical observations, $T$ the retrieved or supplied technical
text, and $M$ the language-model mechanism, including its version, prompts and
decoding policy. A proposal is then a function or random output of
$(D,T,M)$. It is not an additional independent sample from the banking process.
If the model proposes a graph $G$, the relevant questions are whether the
variables have the intended meaning, whether the graph restrictions are
scientifically warranted, whether the data support the inference procedure,
and whether the target effect is identified under those restrictions. These
questions remain distinct even when one conversational interface presents all
of them together.

The following analysis uses inspected primary methods and evaluations. Reported
benchmark gains are the authors' results; no external discovery package or
language-model endpoint is executed here. The finite examples and probability
calculations below are project constructions. Their role is to test the logic
of proposed interfaces and interpretations, not to claim bank validation.

\section{What language-only discovery measures}

The broad proof of concept of K{\i}c{\i}man and colleagues motivates language
support for several causal tasks, including the elicitation of background
knowledge and reasoning about causal relationships
\cite{v12_la_2305_00050}. A model can answer a question such as whether an
increase in a borrowing rate affects demand without receiving a numerical
sample. Such an answer may encode useful inherited knowledge. It cannot reveal
which of several observationally compatible assignment mechanisms generated a
particular bank's records unless the prompt or the model's knowledge provides
additional applicable information.

Feng and colleagues investigate relation occurrences in accessible OLMo and
BLOOM pretraining corpora and report that relation frequency is associated with
prediction success \cite{v12_la_2407_19638}. Their controlled additions of
synthetic relations and the effects of reversed or negated statements help
separate recall from a claim of general discovery. The scope matters: their
pretraining investigation concerns accessible models and predominantly
commonsense relationships. It is evidence for a contamination concern and a
test design, not proof that every model operates only by memorization.

A sensible banking evaluation therefore crosses two kinds of novelty. One can
change the variable names while holding the mechanism fixed, or change the
mechanism while retaining familiar labels. The first tests reliance on
surface vocabulary. The second asks whether a familiar story prevents the
system from responding to contrary evidence. Neither test should declare a
new label intrinsically uninformative: operational definitions often carry
real information about intervention feasibility and time order. Removing all
meaning from a variable name changes the task as well as its contamination risk.

Context must also be handled carefully. In the inspected reliability study,
positive and negative contexts alter affirmative response rates, with contexts
generated by a model and assessed by human annotators
\cite{v12_la_2407_19638}. Some contextual changes legitimately change a causal
claim. A policy that affects demand in one regulatory regime may have no effect
in another. A robustness test should either hold the target world fixed or
declare the changed world and its new truth. Scoring every changed answer as
instability would conflate semantic responsiveness with error.

Jiralerspong and colleagues use breadth-first organization to reduce graph
construction calls \cite{v12_la_2402_01207}. This offers an engineering
alternative to asking about every pair separately. Call complexity is not total
resource cost. A call that lists many nodes may have a growing prompt and
response, require retries, or depend on a long context. A complete comparison
must record input and output tokens, wall time, cache policy, model cost,
failure frequency and the review time needed to adjudicate proposals. Rejecting
a cycle is useful structural validation, but does not make the remaining
acyclic graph true.

\section{Response tests are not sampling-law independence tests}

Cohrs and colleagues investigate language-model judgments in a
constraint-based discovery workflow \cite{v12_la_2406_07378}. This is a
different mechanism from applying a conditional-independence test to observed
bank records. For example, the statement
$X\mathrel{\perp\!\!\!\perp}Y\mid Z$ is a claim about a distribution. A
model's frequency of saying ``YES'' to that statement is a property of a
prompted response process. Even a correctly estimated response probability
does not establish the distributional statement.

Static inspection of the pinned \texttt{causal\_GPT} source shows a further
issue that can be examined without making any model calls. Its response
aggregation counts YES and NO within the same response collection. The
inspected test forms proportions $\widehat p_N=N_N/n$ and
$\widehat p_Y=N_Y/n$, pools them as
$\widehat p=(N_N+N_Y)/(2n)$, and uses the denominator
\begin{equation}
 \widehat s_{\mathrm{ind}}=
 \sqrt{\frac{2\widehat p(1-\widehat p)}{n}}.
 \label{eq:la03-independent-denominator}
\end{equation}
The formula resembles an independent two-proportion standard error. Its
interpretation needs the actual sampling relationship between the counts.

Consider the deliberately simple project response model in which independent
draws take values NO, YES or OTHER, with probabilities $p_N,p_Y,p_O$. Let
$I_N$ and $I_Y$ be indicators for the first two categories in one draw. The
categories are mutually exclusive, so
\[
 \operatorname{Cov}(I_N,I_Y)=-p_Np_Y.
\]
Consequently the exact variance of the difference of the sample proportions is
\begin{equation}
 \operatorname{Var}(\widehat p_N-\widehat p_Y)
 =\frac{p_N+p_Y-(p_N-p_Y)^2}{n}.
 \label{eq:la03-multinomial-variance}
\end{equation}
This elementary calculation is a project derivation under the stated
independent categorical law. It is not a claim that actual model calls are
independent or stationary.

Under the binary symmetric response null $p_N=p_Y=1/2$, the true variance is
$1/n$, whereas the square of
Equation~\eqref{eq:la03-independent-denominator} is $1/(2n)$. The denominator
is too small by a factor $\sqrt{2}$. For $n=100$ and a nominal one-sided
level of $0.01$, summing the exact binomial probabilities over the inspected
formula's rejection region gives approximately $0.044313$. The calculation
enumerates all counts with their rational binomial probabilities; it does not
simulate responses or execute the downloaded implementation.

Correcting this variance would address one response-model issue. It would not
convert language judgments into a bank conditional-independence test. Real
calls may share context, cached outputs, decoding artifacts or a changing
backend. The source also converts the response decision to a numerical zero
or one at the independence-test interface. Those values need a type that
exposes their meaning. Passing them through an API named ``pvalue'' cannot
confer sampling-law calibration.

In a funding example, repeated model agreement that customer risk is unrelated
to pricing conditional on branch does not validate the bank's assignment law.
The remedy is to separate a language-evidence object from a numerical test
object. The former can carry source anchors, answer frequencies and unresolved
scientific questions. The latter must identify its observations, null,
sampling assumptions and calibration. Neither object should silently replace
the other.

\section{Retrieval, constraints and soft priors}

Retrieval changes the information available to a model. It may supply an
operational definition or an expert restriction that the numerical data cannot
identify. It can also supply irrelevant or contradictory evidence. Zhang and
colleagues' retrieval-assisted method reports mixed results: in the inspected
ASIA setting, documents reduce skeleton F1 from 1 to 0.73, while other settings
benefit \cite{v12_la_2402_15301}. The result argues for keeping retrieval
conditions visible in a comparison, rather than labeling retrieval as a
uniform improvement.

The reported orientation accuracy in that work is evaluated on recovered
true-positive edges. Its denominator excludes missing or false edges. For a
reader deciding whether a graph is useful for intervention, skeleton error and
orientation error both matter. A system must report the metric and its
denominator before drawing a conclusion about the full graph. Downstream
effect behavior remains a separate outcome.

A retrieved statement may concern an ancestral path rather than direct
adjacency. Ban and colleagues elicit ancestral constraints and use soft
penalties alongside numerical scores \cite{v12_la_2306_16902}. In a chain
$A\to M\to Y$, $A$ is an ancestor of $Y$ even when no direct $A\to Y$
edge exists. If an adapter turns ``can affect'' into ``directly causes,'' it
changes the scientific input. Path, adjacency, temporal-order and exclusion
constraints therefore need different representations and conversion tests.

Hard constraints delete graphs. If a wrong constraint removes the true graph,
increasing the numerical sample does not restore it within the constrained
search. Darvariu and colleagues investigate soft language-derived priors that
retain support for possible graphs \cite{v12_la_2405_13551}. In schematic
form, one can write a search objective as
\begin{equation}
 \widehat G_\lambda\in\arg\max_G
 \{S(D,G)+\lambda L(T,M,G)\},
 \label{eq:la03-soft-score}
\end{equation}
where $S$ is a numerical score, $L$ a language-derived preference and
$\lambda$ a declared strength. This notation is a common project comparison
device; the actual score and statistical assumptions must be taken from the
chosen method.

Positive prior support prevents one particular failure: a graph is not
logically excluded solely because a model dislikes it. A large wrong
preference can still dominate weak numerical evidence. Subsequent pruning can
again remove the relevant graph. The inspected experiments include selected
combined gains and cases where standalone priors harm performance. A useful
evaluation fixes the prior rule, compares $\lambda=0$, varies evidence quality
and preserves every failed setting. Tuning $\lambda$ after viewing the final
test graphs invalidates a simple interpretation of the reported comparison.

Takayama and colleagues instead feed statistical discovery information into
language prompts, then convert response probabilities into forced, forbidden
or unknown constraints for particular algorithms
\cite{v12_la_2402_01454}. The conversion differs across PC, exact search and
DirectLiNGAM. Their inspected examples contain both gains and failures and
caution against treating the procedure as an opaque automatic solution.
A bootstrap frequency, a response frequency and a posterior model probability
are three different objects. A user interface should not place them under one
unqualified confidence label.

Statistical prompting also raises a data-reuse issue. If $D$ determines the
initial graph summary and that summary shapes the language restriction, the
restriction is a function of $D$. A later fit cannot casually treat it as
independent external evidence. This does not make every such procedure
invalid; it changes the procedure that must be justified or evaluated.
Tracking data hashes and the full adaptive rule makes the reuse visible.
Sample splitting may help a particular construction, but it neither proves the
restriction true nor automatically supplies valid downstream uncertainty.

\section{Debate, disagreement and scientific evidence}

The multi-agent discovery approach of Le and colleagues assigns different
roles to language agents and uses their interaction to construct a graph
\cite{v12_la_2407_15073}. This can organize critique and expose alternative
stories. Its inspected assumptions exclude hidden confounding and cycles, and
the reported canonical benchmark scale is limited. Those conditions should
remain attached to a banking interpretation.

Several agents need not provide several independent pieces of evidence. They
may use the same model family, encounter overlapping training material, or
retrieve the same article. A majority can therefore amplify a common error.
Disagreement is useful as a review signal, but agreement cannot be converted
to a causal confidence interval without a model of the information and
dependence involved.

For servicing, one agent may propose that a contact causes distress, another
that distress causes contact, and a critic may resolve the conflict by choosing
the more plausible story. A more defensible workflow keeps both candidates,
asks which timestamp and assignment evidence would distinguish them, and
identifies the effects under each. A later scientific owner may supply an
applicable restriction. The agent's role is to formulate the unresolved
question and retrieve the evidence, not to manufacture the answer.

Candidate remedies include comparing one-agent and multi-agent conditions,
using the same numerical evidence and total budget, and clustering evaluation
results by shared mechanism rather than counting paraphrased prompts as
independent cases. These are proposed tests. No multi-agent trial has been
executed here, and diversity of models alone does not verify assumptions.

\section{Graph ambiguity and the target effect}

A graph matters because it changes the interpretation of an intervention.
Suppose an observed distribution $P$ is compatible with a class
$\mathcal G$. For a target $Q$, let $\Theta_G(P,Q)$ contain the effects
allowed by the distribution and assumptions under $G$. The scientifically
relevant target set is
\begin{equation}
 \Theta_{\mathcal G}(P,Q)=\bigcup_{G\in\mathcal G}\Theta_G(P,Q).
 \label{eq:la03-effect-union}
\end{equation}
This is a definition, not a claim that the set has been computed or that a
learned class covers the truth. An unrestricted member can make the union
unrestricted. A point estimate from the most fluent graph narrative would
answer a different, more strongly conditioned question.

The project's two-variable example makes this concrete. Under model F,
$A\sim\mathrm{Bernoulli}(1/2)$ and $Y=A$ with probability $0.9$, with an
independent flip otherwise. Under model R the roles are reversed:
$Y\sim\mathrm{Bernoulli}(1/2)$ and $A=Y$ with probability $0.9$. Both
models yield
\[
 P(0,0)=P(1,1)=0.45,\qquad P(0,1)=P(1,0)=0.05.
\]
Replacing $A$'s equation by an intervention gives an average effect of $0.8$
under F and zero under R. In this deliberately finite class the target set is
$\{0,0.8\}$. Neither acyclicity nor observational fit selects the effect.
A sufficiently informative scientific time order or intervention could remove
an alternative, but the example intentionally assumes only a coarse shared
recording time.

The distinction between set identification and confidence in a learned set is
equally important. Let $G_\star$ be the true graph, $\widehat{\mathcal G}$ a
random retained class, and $C_G$ a random set for the target under $G$. Define
$C=\bigcup_{G\in\widehat{\mathcal G}}C_G$. If the actual joint procedure
satisfies
\[
 \Pr(G_\star\notin\widehat{\mathcal G})\le\alpha_G,
 \qquad
 \Pr(\theta_\star\notin C_{G_\star})\le\alpha_\theta,
\]
then the elementary union bound gives
\begin{equation}
 \Pr(\theta_\star\notin C)\le\alpha_G+\alpha_\theta.
 \label{eq:la03-coverage-decomposition}
\end{equation}
The second premise must refer to the actual procedure; it cannot simply be
borrowed from a fixed-model interval after unaccounted model selection. No
independence between the two events is needed for the bound. Neither premise
has been established for arbitrary language-assisted temporal discovery here.
The formula exposes what would have to be supplied, rather than supplying it.

CPDAG and PAG representations encode ambiguity that should survive
serialization and explanation. Circle marks, bidirected relations and
unoriented edges cannot be replaced with arrows merely to produce a cleaner
diagram. A PAG is not an assertion that every latent variable has been named.
Similarly a finite list of proposed DAGs is not automatically a complete
equivalence class. The declaration validator preserves the marks and the
declared alternatives while delegating full class recognition and
identification to a formal engine.

\section{Temporal discovery is not longitudinal policy inference}

The word ``sequential'' refers to several distinct problems. Verma and
colleagues process successive batches of contemporaneous observations while
querying a language model as a noisy expert and retaining PAG information
\cite{v12_la_2506_16234}. The query budget and answer histograms concern
the information supplied to a discovery procedure. They do not identify the
value of a sequence of treatment decisions. Semantic entropy of model answers
does not by itself establish simultaneous graph coverage.

RealTCD combines language-derived initialization with numerical optimization
over temporal graphs and intervention targets
\cite{v12_la_2404_14786}. Its setup assumes causal sufficiency, a known
observational regime and an unrolled finite-lag DAG. The inspected theorem
conditions its population optimization result on model capacity, intervention
faithfulness, positivity and finite entropy. A stochastic optimizer with a
thresholded adjacency matrix requires additional analysis before it can inherit
any finite-sample guarantee.

The real cooling-system study does not have a complete ground-truth DAG. It
uses selected expert physical-location rules to assess edges. That is useful
domain evidence but a different evaluation from full known-truth graph
recovery. The method also models abnormal windows as intervention regimes.
In a bank, a sudden spread movement could reflect a common market shock, a
changed selection mechanism or measurement failure. Treating it as an
exogenous intervention requires substantive evidence, not a threshold crossing.

Kougioulis and colleagues' Large Causal Model is an adjacent numerical
transformer rather than a language-model wrapper
\cite{v12_la_2602_18662}. Its inspected scope includes additive-noise,
stationary, causally sufficient, lagged-only systems. These conditions can be
useful for a declared class; they exclude several concerns in private banking
states and common shocks. The external pretrained weights have not been
audited or executed locally. A language interface should not rename this
restricted temporal predictor as a general latent-graph solution.

For longitudinal servicing, the analyst instead needs a sequence of decisions,
available histories, treatment-confounder feedback, censoring and a value
target. A graph of temporal dependence is only one ingredient. The
intermediate account state can be affected by an earlier action while
confounding a later action's effect. The correct question and inferential law
must be specified before an agent chooses an estimator. None of the inspected
temporal discovery results closes the project's general latent temporal
effect-set problem.

\section{New interventions, mechanism proposals and information value}

Lampinen and colleagues show how passive learning from expert demonstrations
can support active causal strategies \cite{v12_la_2305_16183}. Their
constructive argument first acquires causal information through experiments,
then exploits it. The evaluated learner can intervene in its test environment.
This is a valuable qualification: passive training is compatible with active
information acquisition, but the new interventions supply information that an
analyst restricted to historical bank records may lack.

The inspected empirical settings use small finite acyclic systems and
constrained observations and actions. A banking translation must specify which
interventions are feasible, who can authorize them, their costs and their
effect on future information. Merely asking a model another question is not
an intervention on the customer or the funding market. The two actions have
different likelihoods, costs and ethical or operational constraints.

Murphy's Model Discovery Agent proposes executable mechanisms with an LLM and
fits and compares them numerically \cite{v12_la_2608_09696}. The inspected
studies include algebraic mechanisms, latent ODEs and stochastic latent
dynamics. The LLM expands the candidate library; BIC or SMC approximations
provide model evidence, and a model-disagreement surrogate approximates
expected information gain for experiment selection.

This architecture gives each component a clearer role than requesting a final
effect directly from the model. It also exposes new obligations. An adaptive
library may omit the true mechanism. Pruning can discard a mechanism that
would become useful later. BIC weights within a retained pool are conditional
model weights, not uniform graph-confidence statements. Residual checks can
reveal some misspecification but cannot certify that no policy-relevant
mechanism is absent.

The paper explicitly states that its finite held-out prediction tests do not
establish causal identification. The experiments use open-loop policies, and
the reported comparisons focus on sample efficiency rather than matched total
compute. The chemistry comparison retains matched evaluable cells from a
larger attempted collection. Those denominators must accompany the gains.
For a bank, the value of information is the expected improvement in an actual
decision after paying acquisition and implementation costs. Information gain
within a model pool can be a design surrogate, but it is not that monetary
quantity without an explicit connection to the bank's loss.

\section{An auditable discovery interface}

Zheng and colleagues describe causal-learn+ with agents assisting data
inspection, assumption elicitation, tool choice and interpretation around a
formal discovery core \cite{v12_la_2606_23608}. They emphasize separate
objects for suggestions, user decisions, algorithm settings and graph outputs.
The design is a useful default for a banking library because it makes changes
to scientific input visible. It is not a universal impossibility result for
hybrid discovery; a hybrid can be evaluated as its own fully specified method.

CausalAgent supplies another conversational architecture, including MCP tool
interfaces \cite{v12_la_2602_11527}. Static source review finds that an
inspected OLC tool calls a function not imported or defined in that module,
and a parser can return an empty dictionary after malformed JSON. These are
integration findings about pinned source, not measurements of a running
service. They show why transport success, parse success and scientific success
need separate dispositions. A successful model message is not a valid graph.

For each request, an auditable interface should retain the business target,
variable definitions and clocks, numerical-data identity, retrieved-source
identity, model and prompt version, proposal, applied restrictions, selected
algorithm and returned graph class. A reviewer should be able to distinguish
an expert's asserted exclusion from an algorithmic orientation and from a
model's suggestion. Changed restrictions create a new analysis identity;
they should not overwrite the prior result invisibly.

Failure recovery must preserve scientific meaning as well as computation. If a
job stops after producing a partial graph, replay must not publish it as a
complete answer. If a source or data hash changes, an earlier identification
claim cannot simply be reused. A report should expose a refused or unknown
effect when that is the formal outcome. The completed local operation contract
provides a bounded implementation surface; new discovery adapters still need
their own source, semantic and execution qualification.

\section{What would establish useful progress}

A meaningful comparison begins with a fixed question and information set.
The numerical baseline, language-only baseline, retrieval condition, hybrid
and agent-assisted workflow should receive explicitly matched information or
have their information advantage declared. Accuracy on familiar graphs is
only one outcome. The study also needs identification correctness, retention
of compatible alternatives, valid abstention, effect-set coverage where a
guarantee is claimed, interval or set width, decision loss and full cost.

Every request belongs in the denominator. Failed parsing, unavailable models,
timeouts, nonidentification and legitimate refusals have different meanings
and should be recorded separately. A graph metric conditional on successful
requests can be informative, but it must accompany the success probability
and the selection rule. A narrow accurate set emitted only in easy cases is
not the same product as one that delivers useful valid output reliably.

The proposed remedies are correspondingly bounded. Novel labels test one
shortcut; counterevidence tests retrieval dependence; correctly typed
constraints prevent a specific semantic conversion; fixed prompts and budgets
make an experiment reproducible; and preserved graph classes prevent one form
of overstatement. None of these repairs creates missing bank assignment
truth, resolves observational equivalence by itself, or proves arbitrary
dependent-sample calibration.

This opening comparison supports language assistance for explicit proposals
and reviewable workflow support, with each hybrid method qualified separately.
The following studies develop that position through information limits,
numerical alternatives, source-specific failures and decision consequences.
They distinguish the restricted results already available from the general
temporal graph and useful-precision questions that remain open. New scientific
information, justified assumptions or additional experiments may change the
answer; more fluent explanations of the same information cannot.


% SOURCE docs/release/2026-10-29/source/docs/writing-repair/2026-10-29/src/3fac68a68022ec2f.tex CHAR 0
\hypertarget{temporal-predictions-and-interventions-on-conversations}{%
\section{Temporal predictions and interventions on
conversations}\label{temporal-predictions-and-interventions-on-conversations}}

\label{exec13:census-temporal}

A time series can support a useful prediction without identifying a
useful intervention. The distinction remains necessary when a language
model produces a graph and an accompanying explanation. Augur offers a
concrete example: a teacher model proposes relations between covariates,
a student learns from selected explanations, and the downstream task
predicts a future trend. The inspected preprint evaluates air, power,
traffic and finance settings, using classification F1 and AUROC for
prediction. These are relevant forecasting outcomes; they do not
estimate the value of a deposit rate change or a sequence of servicing
actions. \cite{exec13_census_augur_arxiv}

\hypertarget{what-a-temporal-graph-proposal-establishes}{%
\subsection{What a temporal graph proposal
establishes}\label{what-a-temporal-graph-proposal-establishes}}

Augur first retains highly correlated pairs using Spearman ranks, asks a
language model to choose among directed, confounded and spurious
relationships, and repairs graph cycles. Its stability score measures
agreement across candidate graphs. Its groundedness score checks that a
summary agrees with the graph from which it was generated. Each
component has a clear engineering purpose, but the latter two can be
high when all candidates share the same wrong causal premise. Acyclicity
is a structural property; agreement with one's own graph is an internal
consistency property. \cite{exec13_census_augur_arxiv}

The authors explicitly acknowledge causal sufficiency and the
possibility that correlation pruning misses nonlinear or lagged
dependence. Asking the model to output ``confounded'\,' does not
independently demonstrate hidden-state identification. Nor does a
chronological train/test split create exchangeability of a proposed
action. In a deposit application, funding cost and deposit outflow may
respond to an unobserved liquidity shock. A predictor can exploit that
relationship while remaining unable to answer what outflow would have
been under a different offered rate.

The proposal should retain this architecture as a forecasting and
hypothesis generation comparator. A causal evaluation needs a separate
target and separate truth: for example, generated lagged mechanisms with
known interventions, omitted common causes and nonlinear dependence
deliberately missed by marginal screening. Report predictive performance
and intervention performance separately. A failure of intervention
recovery can coexist with a useful forecast and should lead to a
restricted product role, rather than erasing the forecast result or
relabeling it as policy evidence.

\hypertarget{a-directly-relevant-longitudinal-method}{%
\subsection{A directly relevant longitudinal
method}\label{a-directly-relevant-longitudinal-method}}

Zhang, Wang and Dhillon study sequences of human edits and model
replies. Their CausalCollab method proposes an Incremental Stylistic
Effect (ISE), uses a conditional variational autoencoder to represent
editing styles, and combines a fitted outcome model with Monte Carlo
evaluation of a longitudinal functional. Their identification theorem
explicitly assumes sequential exchangeability and positivity for the
style interventions. Thus there is a concrete literature route for
studying longitudinal human--LM collaboration; its existence should not
be obscured by a search that returns mostly generic causal-inference
papers. \cite{exec13_census_collaboration}

For clarity, let \(L_t\) denote information presented before editing
step \(t\), \(A_t\) the human edit, and \(H_t=(\bar L_t,\bar A_{t-1})\)
the recorded history. A potential terminal outcome is \(Y(\bar a_T)\).
Exchangeability requires the recorded history to remove the common
causes of the edit and its downstream outcome. An analyst's private
customer knowledge, skill or intention can violate this requirement even
when every message has been logged. The low dimensionality of a style
score also does not establish overlap: legally required text may allow
little variation in some histories, and certain kinds of edit may never
occur for distressed customers.

The source defines the outcome associated with a style functional
through a conditional mixture over text versions having that style. This
matters for implementation. A learned coordinate and an instruction such
as ``make this more reassuring'\,' need not describe the same
intervention. A model can make a response more reassuring by changing
tone, by promising compensation, or by omitting a warning. Those changes
can have different effects and costs. Before adopting the estimand, the
proposal must specify which distribution of complete edits implements
the change and which content must be held fixed.

Consider a second project-generated reference example. Two proposed
messages receive the same ``polite'\,' label. Their stipulated response
probabilities are \(1/5\) and \(4/5\), while the reference message has
response probability \(1/10\). Suppose an editing policy selects the
second message with probability \(q\). Its effect against the reference
is

\[
 (1-q)\frac15+q\frac45-\frac1{10}
 =\frac1{10}+\frac35q.
\]

The effects are \(1/10,2/5,7/10\) at \(q=0,1/2,1\), respectively,
despite an identical style label. These are exact values under an
invented law, not observed customer effects. They show why the mixture
is part of the intervention definition. A randomized choice among fixed
text variants is a substantive alternative to inferring a latent
stylistic effect. It has a narrower target but makes treatment versions
directly inspectable.

\hypertarget{what-the-available-evidence-leaves-open}{%
\subsection{What the available evidence leaves
open}\label{what-the-available-evidence-leaves-open}}

CausalCollab constructs counterfactual edits with ChatGPT, uses
model-assigned outcome labels, and changes the association between a
confounding signal and the outcome in its evaluation data. Its three
sources include human--model collaborative writing, model self-chat, and
a hybrid dialogue setting. They should not all be described as
independently observed human intervention outcomes. Performance under
this constructed design is useful evidence about that design, but is not
a direct measurement of how real users would have responded to unchosen
edits. \cite{exec13_census_collaboration}

This supports an affordable proposal-stage evaluation. Begin with public
or generated interaction histories, fix a small set of
content-preserving edit policies, and expose the evaluator to known
latent intentions and missing history in deliberately adverse cases.
Compare a fixed editing policy, an observed-association rule, a
correctly specified numerical longitudinal method, and the proposed
representation-based method. Keep the true generator inaccessible to the
method. Assess policy-value error, support, refusal and total review
effort, in addition to language quality. Generated labels should be
identified as such and their sensitivity to a second label rule should
be reported.

A live human or bank experiment becomes valuable when this work
identifies a consequential uncertainty that the generated law cannot
answer. It is not a prerequisite for specifying the architecture or
exposing its failure mechanisms. Conversely, neither the inspected
theorem nor a generated experiment entitles the proposal to claim
identification under unrecorded human intentions. Direct replication
still requires review of the exact representation training, integration
scheme, outcome labels and source-code versions.

\label{exec13:census-temporal-end}


% SOURCE docs/release/2026-10-29/source/docs/writing-repair/2026-10-29/src/a651df2c07ecf09e.tex CHAR 0
\section{What additional information can language actually supply?}
\label{exec14:l4-development}

The preceding comparisons leave a practical design question. If an assistant
can explain a sensible causal story but the data identify only an equivalence
class, which additional information is the explanation contributing? This can
be answered precisely before choosing a model provider. A restriction extracted
from a dated policy document may eliminate a graph because the policy fixes the
order of decisions. A recollection about lending in general may merely express
a prior preference. A paraphrase of the same empirical association contributes
no new identifying information. The interface should make these three roles
visible to the analyst.

Write $P$ for the observed-data law and $\mathcal G(P;\mathcal A)$ for the set
of graphs compatible with that law under assumptions $\mathcal A$. A numerical
discovery procedure estimates features of this set. A text-derived restriction
$R$ instead proposes the smaller set
\[
 \mathcal G_R(P;\mathcal A)=\{G\in\mathcal G(P;\mathcal A):R(G)=1\}.
\]
This is a useful operation if the restriction is applicable and justified. It
is also a way to manufacture certainty when an erroneous restriction deletes
the true graph. The source, date, population and interpretation of $R$ are
therefore part of the scientific object. A statement that an account review
occurs before an offer concerns a time order; it does not by itself assert that
review has no other causes or that all such causes were measured.

For a deposit-pricing question, a policy may establish that a national rate is
announced before branch offers are made. That justifies a temporal restriction
for the relevant dates. It does not establish that managers did not anticipate
the announcement, that the national rate has no other effect on deposits, or
that branch offers were assigned independently of unrecorded local risk. A
language assistant should return the restriction actually supported by the
passage and list the stronger restrictions that would be needed for the desired
intervention calculation. This follows the distinction between an assumption
and its observable implications in graphical model construction
\citep{exec13_census_credible}.

An empty $\mathcal G_R$ is an informative conflict. It can result from an
incorrect text interpretation, an unsuitable statistical model, sampling error,
a missing variable or inconsistent restrictions supplied by different owners.
The proposed tool should return the smallest diagnosed conflicting subset when
it can compute one, and otherwise retain the full conflict. It should not
silently discard low-ranked edges until a DAG appears. The original evidence
and all relaxations remain visible. Producing a graph after changing the
restrictions is a new analysis whose identifying content has changed.

A soft restriction makes a different tradeoff. A score of the form
$S(G;D)+\lambda R_M(G;T)$ combines statistical fit with a model-derived
preference. The regularization weight $\lambda$ determines how much the
language preference can change the selected structure. Neither its numerical
scale nor an edge frequency from repeated prompts supplies a posterior
probability that the banking mechanism contains an edge. A meaningful
sensitivity analysis varies the strength of the preference, keeps the data
term and feasible graph family fixed, and traces which decisions change.
Soft-constraint discovery is a legitimate alternative to a hard restriction,
but it must be compared under its stated meaning
\citep{v12_la_2405_13551,v12_la_2407_15073}.

\subsection{An observational equivalence bound for an assistant}

A short project derivation explains a limit that a more fluent assistant cannot
remove without new information. Suppose two structural models $\mathcal M_0$
and $\mathcal M_1$ generate the same observed law for every available record and
the same admissible documents, while their target effects are $\theta_0$ and
$\theta_1$, with $\Delta=|\theta_1-\theta_0|>0$. Let the assistant, numerical
code and critic together be any randomized procedure using only those records
and documents. Its internal randomization has the same distribution in both
worlds and contains no world-specific oracle. Then its reported interval $C$
also has the same distribution in both worlds.

If the procedure claims coverage at least $1-\alpha$ under each model, the two
events $\{\theta_0\in C\}$ and $\{\theta_1\in C\}$ have probabilities at least
$1-\alpha$ under that common report law. The elementary union bound gives
\begin{equation}
 \Pr\{\theta_0,\theta_1\in C\}\geq 1-2\alpha,
 \qquad
 \Pr\{\operatorname{diam}(C)\geq\Delta\}\geq1-2\alpha.
 \label{exec14:language-information-bound}
\end{equation}
The diameter statement also holds for a set-valued report. It does not require
normality, a particular discovery algorithm or a claim that an LLM is
statistically calibrated. It says what uniform coverage would require in this
two-world experiment. For nominal 95 percent coverage and an effect separation
of $2/5$, diameter at least $2/5$ is necessary with probability at least $9/10$.
An algorithm that always returns a narrower interval cannot satisfy both
coverage claims. A refusal is an allowed scientific result; counting it as a
successful narrow interval would change the claim.

The same argument has a decision form. If the correct binary action differs
between the two worlds, let $q$ be the probability that the workflow selects
action one. With unit loss for a wrong action, its risks in the two worlds are
$q$ and $1-q$, in one order. Their maximum is at least $1/2$. More independent
prompt draws can estimate $q$ more precisely; they cannot make both risks small
when the observations and applicable text are the same. The way to improve the
decision is to obtain a discriminating restriction or intervention, accept a
robust action, or state the unresolved risk. The source of a useful restriction
may be a person or a document, but its validity is additional information.

This limitation is directly relevant to the funding equivalence example in the
Java adapter discussion, where several Gaussian models share an observed
covariance yet differ in the intervention response. It also sharpens a fair
language evaluation. If the supplied narrative explicitly tells the model which
world generated the data, the task assesses use of that information. It must
not be described as identification from the numerical observations alone. If
model pretraining contains relevant proprietary policy information, that
information is also part of the effective input and raises a provenance and
transport question; it is not an invisible extension of a statistical theorem.

\section{From edge accuracy to the bank's decision loss}

The discovery target must be defined before a leaderboard is useful. A DAG,
a CPDAG, a PAG and an RFCI-PAG make different assertions. A DAG adds directions
that may be unresolved in a CPDAG. A PAG can retain circles because ancestral
relations differ across compatible latent-variable models. An RFCI adjacency
can have a weaker meaning than an FCI adjacency. Comparing all four as lists of
directed edges against a single chosen DAG can reward unwarranted certainty
and punish scientifically correct ambiguity
\citep{exec13_latentr_colombo2012,exec13_latentr_primary}.

A proposed evaluation should therefore publish separate losses for the
skeleton, endpoint marks, invariant ancestral claims, extra assumptions and the
downstream estimand. If $E^\star$ is the true adjacency set under a declared
reference and $\widehat E$ the reported adjacency set, precision and recall
answer whether proposed links are right and whether reference links are found.
An orientation score computed only on $E^\star\cap\widehat E$ answers a
conditional question; it can look excellent when most links were missed.
Unknown endpoints should have their own accounting, with a penalty only for
claims the task actually requires. The interpretation of retrieval-assisted
orientation rates depends on precisely this denominator
\citep{v12_la_2402_15301}.

For an intervention, graph error matters through a loss function. Let
$\Theta(\widehat{\mathcal G})$ collect effects compatible with a reported graph
class and other stated restrictions. A proposed action can maximize worst-case
utility, maximize utility under an explicitly elicited model distribution, or
request more information. These are different decision rules. The assistant
must not silently replace the first with the second by selecting the graph it
finds most plausible. The nominal effect interval and the range over graph
models are also different objects; combining them requires a stated procedure.

Consider an exact planning example in which an offer costs 3 units per account
and its incremental gross return is either 0 or 8 under two observationally
compatible models. The net value of offering is either $-3$ or $5$, whereas
not offering has value zero. The worst-case expected-value rule selects no
offer. If an elicited probability $q$ is placed on the high-return model, the
expected net value is $8q-3$ and offering is chosen only for $q>3/8$. This
threshold makes the consequence of the extra belief visible. A language-model
vote frequency cannot be substituted for $q$ without a model relating votes to
which world is true. The costs and returns are stipulated units, not bank
prices or measured ROI.

An ideal, free observation that perfectly reveals the model would yield value
$5q$, since the offer would be made only in the high-return world. The value of
that information before paying for it is
\[
 5q-\max(0,8q-3).
\]
At $q=1/2$ it equals $3/2$ units per account. A proposed experiment costing more
than that amount per account could not be justified by this particular perfect
information calculation, and a noisy experiment cannot exceed its value under
the stipulated model. This is a planning bound, not a recommendation to conduct
a bank trial now. It identifies what an inexpensive public or generated
prototype should measure: whether the workflow preserves the ambiguity and
correctly reports the threshold for a decision.

The criterion changes when loss includes customer harm, regulatory constraints
or uneven effects across groups. Those obligations must enter the feasible
action set and utility definition explicitly. A large mean return cannot by
itself authorize an infeasible offer. Similarly, a loss expressed per account
requires care when one customer holds several accounts; the intervention and
interference units may differ. The graph workflow should surface these choices
before performing a cost comparison, rather than use the presence of a currency
symbol as evidence that the causal problem has been economically specified.

\section{Two channels of uncertainty and a calibrated refusal}

Graphical ambiguity and statistical uncertainty should survive into the same
report without being conflated. The first arises when the same population law
is compatible with more than one target effect under the stated information.
The second arises because that law is estimated from a finite sample. A third
source, uncertainty about whether the variables and assumptions are meaningful,
is not automatically represented by either a confidence interval or a PAG.
A language assistant is often most useful on this third source: it can find an
ambiguous definition or an unsupported restriction. That usefulness does not
make its confidence score a measure of the first two.

One proposed report therefore has a target object, an identification object,
a sampling object and an evidence-review object. The identification object
states ``identified under these assumptions'', ``nonidentified in this model
class'', or ``unknown because the available procedure or input is inadequate''.
Failure to find a back-door adjustment set with a restricted implementation is
not a proof of nonidentification. The sampling object states the estimator,
method of uncertainty calculation, population and any trimming or weighting.
The evidence object links assumptions to inspected sources and distinguishes
missing evidence from a retrieval failure. A failed numerical test is neither
independence nor evidence of an unobserved common cause.

CausalDS is a useful comparator because it makes abstention on specified
nonidentified tasks part of the evaluation, and separates conceptual variables
from their noisy observation layer. Its finite benchmark findings remain
reported evidence for those tasks. The source-level adaptations needed for our
prototype are explicit three-state identification metadata, finite output
checks and denominators that distinguish principled abstention from missing
files, parsing errors and numerical failures
\citep{v12_la_2607_08093,exec14_estimation_sources}.

A refusal can be correct without being useful enough for a product. The
comparison should score whether it names the missing assumption, unsupported
operation or conflicting clock and whether a feasible next action is supplied.
A rule that refuses every request avoids some false causal claims while failing
all identifiable tasks. A rule that returns a number for every request has high
completion while failing nonidentified ones. Plotting error and returned-answer
rate over a frozen threshold is more informative than a single completion
percentage, provided the threshold was selected separately from the held-out
assessment. All attempted requests, including refusals and malformed outputs,
remain in the operating-cost denominator.

A critic can select which analyses reach the report, but acceptance by a critic
is not automatically conditional statistical coverage. As an exact counting
example, six covered intervals among six accepted requests, with two refusals,
give selected coverage one and an all-request return-and-cover rate of $3/4$.
Neither denominator alone shows a nominal conditional guarantee. To investigate
that claim, a prospective experiment must define the selection rule, tasks and
repeated sampling law, and report uncertainty for its empirical rates. The
critic's access to a reference answer must be separated from the estimator's
information. Reviewing one's own answer with the same model adds an interaction,
not independent scientific evidence.

\section{Recent agent systems as competing scientific workflows}

Recent systems should be compared by the boundary at which the agent's output
becomes authoritative. In one architecture the agent directly proposes edges.
In another it chooses settings for a numerical engine. In a third it edits a
shared data model or proof statement. A fourth retrieves evidence for a human
reviewer. All can be described as causal agents, but they place different
obligations on the scientific core and on the person who accepts the result.
The banking prototype should choose a boundary that can be tested with the
information it actually has.

Zhu and colleagues' CausalAgent combines data processing, structure learning and
reporting with MCP tooling and a Sachs demonstration. The inspected paper places
quantitative effect estimation and counterfactual modules among its future
directions. Its strongest immediate contribution to this proposal is a
workflow pattern for moving from data and a question to a structured report.
The proposal cannot use the presence of MCP tools as evidence that an early
arrears-review effect has been identified, or that the future quantitative
module is already available. A direct integration decision must inspect the
called implementation and what it does when parsing or a backend call fails
\citep{v12_la_2602_11527,exec13_data_causalagent_code}.

Han and colleagues' separately authored Causal Agent is a different source.
It illustrates named graph objects stored in tool memory and a generated
question benchmark. This is useful for persistent state: an analyst should be
able to distinguish the graph before a restriction from the graph after it,
and the target before trimming from the target after it. Yet storing a graph
under a durable name does not validate its construction. The proposed service
binds a graph's content, variable definitions, information clock and added
assumptions together. Two graph objects with the same edge list can still be
different scientific inputs if their variables or populations differ
\citep{exec13_census_causal_agent}.

ORCA places more emphasis on relational state, including joins, variable roles,
graphs and estimation configuration. This addresses a real banking problem:
joining an account table to a customer table can multiply rows, change the unit
and create a misleading apparent sample size before discovery begins. Its
reported PAG-to-DAG step, however, makes additional orientation or editing
choices. Those choices must be visible as assumptions. The inspected report
also contains a discrepancy between the abstract's 42-percentage-point
completion improvement and a table with completion rates $.83$ and $.75$.
Completion is defined there by production of an ATE without requiring accuracy.
The proposal retains the discrepancy and avoids treating either a completion
figure or the existence of a final estimate as evidence of correct causal
inference. No corrected experiment is invented
\citep{exec13_census_orca}.

The causal-learn+ position puts assistance around a protected discovery core.
Agents can help inspect data, explain algorithms and interpret results while
formal algorithms retain responsibility for the numerical discovery operation.
That division is the preferred initial design here because it makes the
scientific method and the agent's added assumptions separately inspectable.
It still requires a correct core: a well-typed call to a flawed algorithm is
not a sound discovery result. It also leaves measurement, latent structure and
sampling assumptions open to review. The paper's call for future evaluation
supports a candidate architecture, not a measured banking advantage
\citep{v12_la_2606_23608}.

CausalSmith adds a different kind of protection by making a formal statement,
its proof dependencies and compiler outcome explicit. A checked proof answers
whether the formal statement follows within its logical environment. It does
not settle whether ``recovery'', ``offer'' and ``available before treatment''
were formalized faithfully, or whether the assumptions describe the bank.
The inspected limitations retain unvalidated model judgments about statement
faithfulness and research quality. The appropriate use is to expose a proof
obligation and its assumptions, then review the mapping to the business problem.
It is not to treat successful compilation as a substitute for that mapping
\citep{exec13_census_causalsmith}.

These choices lead to a concrete reuse recommendation. Adapt stateful
specification and evidence-reporting ideas, retain a fixed independently
reviewed numerical core, and require typed proposals for changes to the target,
variables or graph restrictions. Defer unrestricted autonomous execution until
the full package and operational path is audited. Reject automatic conversion
of an uncertain graph into an asserted DAG merely to enable a downstream
estimator. Keep proof assistance as a separate declared capability whose
formal assumptions remain inspectable. The proposal claims that these boundaries
make errors easier to observe; an empirical claim that they reduce review cost
or improve accuracy awaits a matched study.

This decision is revisable. A more autonomous system could be preferable if a
fixed public or generated evaluation shows that it preserves targets and graph
uncertainty with less total effort. A protected core could be ineffective if
most errors occur in the semantic mapping before the core is called. Both
outcomes are plausible hypotheses. The next experiment should therefore measure
where corrections occur and what information was needed to make them, rather
than count the number of agent roles or reward the longest explanation.

\section{Following the full package path rather than the package name}

A research package is not interchangeable with an executable operation of the
proposed library. The first path is a selected PC call in causal-learn, its
returned graph representation, a conversion into the shared graph schema,
a context-aware diagnostic report and a bounded CLI or MCP response. The
inspected MIT-licensed revision offers a plausible direct integration route,
but its twelve direct dependency declarations do not constitute a full pinned
environment. Porting the algorithm now would add maintenance responsibility
without a measured embedding constraint that justifies it. Executable admission
is therefore deferred while this path is completed; the reuse choice itself is
concrete and reviewable \citep{exec14_causallearn_reuse}.

The PC and FCI profiles must remain distinct. A PC bidirected conflict under an
orientation rule cannot silently become the latent-common-cause assertion of a
PAG. The conversion needs the two proximal endpoint marks, a bijection from
input variables to nodes, and isolated nodes. Unsupported composite endpoint
encodings remain unsupported. A Fisher-Z failure due to a singular inverse,
an invalid $n-|S|-3$ term or a nonfinite numerical value cannot become a valid
independence decision. An $\alpha$ threshold for one query does not supply
whole-graph or effect-set coverage. These are requirements for the first scoped
adapter, not benefits already delivered by listing causal-learn as a dependency
\citep{exec14_causallearn_reuse}.

Tetrad provides a materially different Java path. At the inspected revision,
FCI initializes circle endpoints and invokes orientation code with a flag for
excluding selection bias. The class-level reference to a CPDAG is inconsistent
with the PAG-oriented search path and cannot define the interchange contract.
RFCI is a separate procedure with its own graph meaning. The Edge constructor
can reverse the stored order of nodes for a left-pointing edge, so conversion
must bind each mark to its proximal node, not assume that construction order
is stable. Star and null endpoint values exist in the Java enum but are outside
the proposed initial PAG interchange profile
\citep{exec14_java_latent_sources}.

Static inspection also identifies a specific reason to withhold full FCI
behavioral qualification at that revision. In the inspected possible-D-SEP
block, the code builds one candidate set from each endpoint, with a path-length
argument of three, and tests that complete set. That block does not enumerate
all of its subsets. Conditional independence is not monotone in the conditioning
set: adding a collider can turn independence into dependence. The existing
exact Gaussian counterexample demonstrates this logical point; it is not an
executed counterexample to the entire Tetrad algorithm, whose earlier skeleton
phase and later rules can affect the actual candidate graph. The required next
check follows the whole search on an oracle fixture that reaches this block
\citep{exec13_java_examples,exec14_java_latent_sources}.

A second branch matters when configuration is changed. The ordinary
\texttt{finalOrientation} call is inside the possible-D-SEP branch in the
inspected search method. Setting that option false skips that call; a separate
optional guarantee-PAG branch is not a substitute for an unexamined claim of
algorithmic equivalence. The conclusion is bounded: the source path needs
qualification, not that every Tetrad release or every returned graph is wrong.
The available local Java 11 runtime and absent compiler also do not meet the
Java 21 build profile previously reviewed. No Java FCI, RFCI or endpoint
roundtrip has been executed in this review
\citep{exec14_java_latent_sources,exec13_tetrad_reuse_primary}.

The proper alternatives are an isolated, pinned Tetrad environment with
method-specific oracle and numerical tests; the separately reviewed R pcalg
FCI/RFCI routes; and a separately reviewed Python FCI profile. Their settings,
selection assumptions, graph classes and conversion conventions have to be
matched before any speed comparison. A benchmark that times a truncated search
against a fuller search without labeling the difference measures a change in
scientific work as well as an implementation difference. A process bridge adds
startup, serialization, memory and recovery costs that belong in the complete
workflow comparison.

\section{A hidden-risk example with identical observable distributions}

The information bound becomes tangible in a nondegenerate Gaussian example.
Here the banking action is a standardized concession intensity, rather than the
binary offer used in the earlier motivating brief. Consider two centered models.
In the first, independent disturbances have variances
$\operatorname{Var}(U)=1/2$ and
$\operatorname{Var}(\epsilon_A)=\operatorname{Var}(\epsilon_Y)=1/2$, and
\[
 A=U+\epsilon_A,\qquad Y=U+\epsilon_Y.
\]
In the second, $A$ has variance one and is independent of $\eta$, with
$\operatorname{Var}(\eta)=3/4$, and
\[
 Y=\tfrac12 A+\eta.
\]
Both produce the bivariate normal law with variances one and covariance $1/2$.
Both have a positive observed density on the plane. Their observed regression
of $Y$ on $A$ has slope $1/2$. In the first model, intervening on $A$ leaves
$Y$ unchanged, so the derivative of the intervention mean is zero. In the
second it is $1/2$. Every statistic of the observed pair has the same sampling
distribution in the two models for every sample size. Recovering a precise
association does not decide which intervention response is right.

The example isolates the role of private borrower information. The first model
contains a common cause that is unmeasured in the supplied table; the second
attributes the association to the concession. An agent given only the same
table and an ambiguous description has no model-specific evidence that selects
one. Saying that the concession ``obviously affects recovery'' introduces an
assumption. It may be reasonable as part of a scenario, but it must be exposed
as such. A PAG over two observed variables can retain unresolved endpoint marks;
returning a single oriented edge removes an uncertainty that the data did not
remove.

Suppose a new predecision score $W=U+\epsilon_W$ is offered as a remedy, with
independent noise variance $s$. In the common-cause model the residual variance
of $U$ given $W$ is
\[
 v=\frac{(1/2)s}{1/2+s}=\frac{s}{1+2s}.
\]
After residualizing $A$ and $Y$ on $W$, their covariance is $v$ and the residual
variance of $A$ is $1/2+v$. The adjusted coefficient is therefore
\[
 \beta_{A\mid W}=\frac{v}{1/2+v}=\frac{2s}{1+4s}.
\]
For $s=1/2$ it is $1/3$ although the true causal response is zero. With an exact
score, $s=0$, it becomes zero; as the score becomes uninformative it approaches
$1/2$. This is a separate model from the unit-variance text-measurement example
in the earlier chapter. The variance choices and action definition are stated
so that the two examples are not mistakenly pooled. A document extractor can
help create $W$, but the identifying question is what information its error
model and validation sample supply.

The remedy is specific. Obtain a defensible measurement model with information
sufficient for the proposed correction, a justified assignment restriction, or
an intervention that distinguishes the worlds. A smaller standard error from
more observations of $(A,Y,W)$ does not remove the residual coefficient when
$s$ stays positive. A language agent can explain the alternatives, request the
necessary metadata and prevent a false report of identification. It cannot
invent an independent validation sample by generating plausible labels.
The exact covariance and residual-variance calculations in this section are
project derivations, reproduced in rational arithmetic
\citep{exec14_discovery_oracles}.

\section{Why a CI threshold is not graph or effect coverage}

A discovery interface often exposes one number called alpha. Its scope must be
made explicit to an agent and to a business reviewer. For a particular fixed
conditional-independence test, a level-$\alpha$ statement concerns the
probability of rejecting that null when it is true under the declared sampling
law. It does not bound the probability that a whole graph is correct, that all
orientations are right, or that a derived effect interval covers its target.
A graph can fail because a false independence is accepted, an assumption is
wrong, a query is numerically invalid, or an orientation is misinterpreted.
Those errors are not all controlled by a single null-rejection level.

Even the limited false-rejection claim changes when many queries are possible.
For $p$ observed variables, all unordered variable pairs and all conditioning
subsets of the remaining variables form a family of size
\[
 M=\binom p2 2^{p-2}.
\]
If each fixed query has unconditional type-I error at most $\alpha$, then a
union bound controls at least one false rejection among that fixed finite
family by $M\alpha$, truncated at one. No independence between query decisions
is needed for this inequality. It applies to any adaptively selected subset of
the family only when the procedure uses the same prespecified tests and the
per-query calibration holds for every fixed query; an adaptive prompt that
changes the data or test definition is not covered by that argument.

With six variables there are 240 possible pair/subset queries. At $\alpha=.05$
the bound is uninformative. At $\alpha=.0002$ it is $.048$ for the specified
false-rejection event. This second number is still not a 95.2 percent guarantee
for the graph: it does not control missed alternatives, faithfulness violations,
latent-variable misspecification or endpoint conversion. The example is useful
because it gives the assistant a precise explanation of what alpha can and
cannot mean. It is not a recommendation to choose a universal threshold of
$.0002$ for bank discovery.

The same distinction applies to language-response tests. Their sampling law is
a law of model answers under a prompt and decoding configuration. A response
frequency can be estimated under assumptions about repeatability, but it does
not inherit the error guarantee of a numerical CI test. An agent that reads a
column named ``pvalue'' must check the generating procedure. The response-based
variance calculation in the preceding material demonstrates one implementation
hazard; the more fundamental issue is that the response null and the banking
distributional null are different statements.

For the prototype, every CI result should therefore include the query, sampling
unit, conditioning variables, test method, usable observations and computation
status. A singular matrix produces an unavailable test, not a very strong
scientific dependence claim. The search result should state whether all
required queries completed and whether any search bounds or time limits
changed the explored family. The report should retain sensitivity to thresholds
and to permitted restrictions without presenting that sensitivity exercise as
a calibrated confidence set. A qualified effect procedure needs its own
identification and uncertainty argument after the graph step.

\section{A public evaluation design that can change the proposal}

The immediate evaluation should test whether language assistance improves the
causal workflow within a fixed information budget. It can use public documents
for realistic vocabulary and project-generated mechanisms for known truth,
with the two origins clearly separated. Generation does not create authentic
bank assignment evidence, but it makes intervention effects, measurement error,
latent variables and timing failures observable to a held-out scorer. Private
truth must not enter the assistant's prompt, retrieval corpus, tool output or
file tree. The predeclared deterministic baseline receives the same permitted
brief and data as the assisted workflow.

A useful task portfolio crosses four dimensions. One varies whether the effect
is identified from the supplied graph and observations. Another varies whether
variable names are informative, neutral or misleading. A third changes whether
the document's reported time and its availability time agree. A fourth varies
whether text contains a valid additional restriction, a contradiction or no
relevant information. This is a design proposal; choosing the generator, sample
size, splits and replication count requires a prospective protocol before any
performance study. A successful run on a few authored examples is a conformance
check, not a retrospective preregistration.

The banking tasks should include an offer distinct from uptake, a branch-level
intervention with account-level observations, an unrecorded borrower-risk
variable, and a postdecision note that looks like a baseline covariate. They
should include a graph with correct but unresolved endpoints and a case where
an apparent numerical failure is caused by exact accounting identities. A
language model that recommends a valid operation but uses the wrong treatment
variable has not answered the task. Conversely, a workflow that retains the
right target and returns a justified nonidentification result can succeed on a
case for which no point estimate is licensed.

Names and mechanisms should be perturbed separately. Replacing all banking
terms with anonymous labels can remove legitimate intervention and timing
information, so that perturbation is scored as a different information regime.
Holding the definitions fixed while changing irrelevant wording probes another
kind of robustness. Reversing a generated mechanism under familiar names tests
whether the model follows data or a memorized story. The inspected work on
pretraining relation frequency motivates this distinction, but its findings
for accessible models and commonsense relationships are not a theorem about
all language models or all banking questions
\citep{v12_la_2407_19638}.

Document-level leakage also needs a precise split. If the same policy is
paraphrased into both development and evaluation sets, a row-level random split
does not produce an unseen policy. If several accounts share one assignment
rule or one customer, independent-row formulas may overstate the effective
sample. The proposed protocol should group related documents, mechanisms and
customers according to the claim it seeks to assess, and say which kind of
transfer remains untested. No generic instruction to ``hold out 20 percent''
settles these scientific dependencies.

For each request the result should contain the target identity, permitted
information, proposed graph class, assumptions added by the agent, formal
identification state, numerical method and final disposition. Metrics should
include false causal assertions, missed identifiable opportunities, unsupported
orientations, effect error where meaningful, interval loss and all-request
resource cost. Model-only latency omits retrieval, serialization, review,
retries, corrections and refused jobs. A cheap request that repeatedly needs an
analyst to repair a target can cost more than a slower fixed workflow.

The proposal changes if the assisted path reliably reduces a specified review
burden while maintaining these scientific boundaries on a frozen task family.
It also changes if assistance adds unsupported restrictions or hides failures,
even when the prose improves. A later live comparison must bind provider and
model versions, prompts, decoding, tool identities, custody and actual prices;
none of these results is inferred from this design. An expensive bank experiment
becomes relevant only after the prototype isolates an uncertainty that public
or generated evidence cannot answer and shows why resolving it can alter the
investment decision.
\label{exec14:l4-development-end}


% SOURCE docs/release/2026-10-29/source/docs/writing-repair/2026-10-29/src/9235cdc16df4fd1e.tex CHAR 0
\section{Choosing the numerical core before asking an agent to choose}
\label{exec14:l4-alternatives}

The proposed protected core still leaves a substantive choice. Constraint
search, score optimization, functional equations and changes of environment
can all return a graph, but they use different information to distinguish
directions. An assistant should explain the reason for choosing among them.
Selecting the package that happens to accept the input table is insufficient.
For a branch-pricing problem, the first questions are whether the table records
independent units or a dependent panel, whether omitted common causes are
plausible, whether the action and outcome have a reliable order, and whether
the proposed structural restrictions are defensible for those variables.
Only then is it useful to compare runtime and interface convenience.

PC uses conditional-independence information with a causally sufficient DAG
interpretation under its assumptions. Its partially directed output preserves
orientations that the available independence information does not determine.
FCI permits latent and selection variables under a different model class and
returns PAG information. RFCI trades some adjacency information for a less
expensive search. The differences matter when an analyst asks whether a
borrower's private assessment jointly influences an offer and subsequent
recovery. Selecting PC after simply deleting that unrecorded assessment from
the variable list does not remove the common cause. Selecting FCI does not
make the effect point identified. Its output may correctly retain precisely
the ambiguity that prevents a unique intervention calculation
\citep{exec13_latentr_colombo2012,exec14_java_latent_sources}.

A score search instead compares how well candidate structures represent the
data after accounting for complexity or prior preferences. A selected DAG
can be an effective computational summary without being the only causal
explanation. The score, its statistical interpretation, the search space and
the optimization outcome are separate choices. If two structures have the
same observable implications under the supplied information, an arbitrary
tie break does not become a discovered direction. A language prior can break
the tie, but the report must then attribute the direction to that added
restriction or preference. For the proposed bank tool, a ranked candidate
list is useful if its conditional meaning remains visible; an unsupported
claim that the first entry is the true mechanism is not.

NOTEARS illustrates both the appeal and the limit of a continuous formulation.
For data matrix $X$ with $n$ rows and $d$ columns, the inspected linear
formulation uses a coefficient matrix $W$ and the objective
\[
 F(W)=\frac{1}{2n}\lVert X-XW\rVert_F^2+\lambda\lVert W\rVert_1,
 \qquad h(W)=\operatorname{tr}\exp(W\circ W)-d=0.
\]
Here $W_{ij}$ represents a coefficient from variable $i$ to variable $j$,
$\circ$ denotes entrywise multiplication, and the exponential is a matrix
exponential. This convention differs from a target-row coefficient matrix.
The acyclicity characterization is an inherited result, not an identifying
assumption that language can supply. If $S=W\circ W$, every entry of $S$
is nonnegative and each $\operatorname{tr}(S^k)$ sums weighted closed walks.
The expansion $\operatorname{tr}\exp(S)=d+\sum_{k\geq1}
\operatorname{tr}(S^k)/k!$ is equal to $d$ exactly when the graph has no
directed cycle. Squaring the entries prevents positive and negative edge
weights from cancelling a cycle in this argument
\citep{notears2018}.

That theorem separates a structural requirement from the inferential question.
A small computed $h(W)$ is a numerical diagnostic of one constraint. It is
not a probability that the graph is correct. The original paper expressly
distinguishes stationary-point optimization from global minimization, and
its statistical interpretation refers to assumptions studied separately.
The inspected reference implementation also has an $\ell_1$ penalty,
augmented-Lagrangian stopping rules and a final weight threshold. Each can
change the output. An agent that raises the threshold until a graph looks
sparse has changed the analysis; it has not improved a calibrated confidence
level. The exact returned graph needs a structural check after thresholding,
and the fit needs an appropriate scientific interpretation
\citep{notears2018,exec13_pydisc_notears}.

An apparent advantage of one discovery method can also depend on units.
Multiplying a deposit balance column by a thousand changes an unstandardized
squared-error objective and the relative meaning of a common coefficient
penalty. A nonlinear transformation changes the functional model still more.
The transformation, back-transformation and intervention scale should be
part of the proposed specification. A rank-based numerical method, a linear
equation method and a language judgment need not have the same invariances.
The comparison should expose this difference instead of assuming that an
identical CSV file makes all models scientifically matched.

LiNGAM obtains directional information from a different restriction: a linear
acyclic system with mutually independent, non-Gaussian disturbances and no
unobserved confounders. Its original model assumes repeated independent
vectors generated with the same equations and disturbance distributions.
The inspected original construction uses independent component analysis;
the currently inspected DirectLiNGAM code is a separate algorithmic route.
The source explicitly notes that this model need not impose faithfulness.
Consequently, a failure of a faithfulness-based argument does not by itself
refute LiNGAM. Conversely, apparent skewness of observed loan balances does
not establish the non-Gaussianity and independence of the structural
disturbances in the chosen equations. These are requirements on the model,
not labels that can be inferred from column names
\citep{exec13_pydisc_shimizu2006,exec13_pydisc_lingam}.

The Gaussian equivalence example earlier in this chapter intentionally lies
outside that non-Gaussian restriction. It cannot be used as an impossibility
proof for every functional discovery model. This distinction is productive:
one can escape an observational equivalence by supplying additional
information or by restricting the admissible structural models. The proposal
should state which of these it is doing and what evidence supports the
restriction. The language component can retrieve a scientific argument for
linear mechanisms, question an independence assumption or compare an
alternative model. It cannot verify disturbance independence merely by
producing a coherent explanation.

For time series, the underlying sampling and causal model changes again.
The inspected PCMCI+ result requires its stated sufficiency, stationarity,
faithfulness, lag coverage and test consistency conditions; selection and
measurement error are outside that theorem. Contemporaneous directions can
remain unresolved. A calendar panel containing common shocks and repeated
accounts cannot inherit the result solely because the records have a date
column. The assistant should identify the sampling unit and candidate lag
family, explain which common causes are represented, and retain uncertainty
about contemporaneous effects. Declaring a lag of one because the prompt is
shorter is a substantive restriction, not a harmless interface choice
\citep{runge2020}.

These comparisons imply a staged software decision. Retain method-specific
profiles for PC, FCI/RFCI, functional equations, continuous score search and
temporal discovery. Each profile binds the statistical assumptions, graph
meaning, matrix convention, search bounds and failure behavior. Reuse a
reviewed implementation when its complete path can be made reproducible;
consider porting only when a measured deployment or numerical requirement
justifies the maintenance burden. A common result schema can preserve their
differences. It should not force every method to emit a fully oriented DAG or
an effect estimate when that is stronger than the method supports.

\section{Faithfulness can fail even with complete measurements}
\label{exec14:l4-faithfulness}

A second exact project example separates missing variables from cancellation.
Let $A,E_M,E_Y$ be independent centered Gaussian variables with unit variance,
and define
\[
 M=A+E_M,\qquad Y=-A+M+E_Y=E_M+E_Y.
\]
The structural graph contains $A\to M$, $A\to Y$ and $M\to Y$. All three
variables are observed, all disturbance variances are positive, and the
joint distribution has a nonsingular covariance matrix. Nevertheless
$A\perp Y$: the direct and mediated paths cancel in the marginal law.
The covariance matrix, in the order $(A,M,Y)$, is
\[
 \Sigma=\begin{pmatrix}1&1&0\\1&2&1\\0&1&2\end{pmatrix},
 \qquad
 \operatorname{Cov}(A,Y\mid M)=0-\frac{1\cdot1}{2}=-\frac12.
\]
Thus independence at the empty conditioning set and dependence after
conditioning can arise here without the structural collider explanation
that a faithful model would use. This is not a numerical roundoff example.
The zero is an exact property of the stipulated population law.

There is indeed a different structural model with that faithful collider
explanation. Take mutually independent $A\sim N(0,1)$, $Y\sim N(0,2)$
and $\eta\sim N(0,1/2)$, and set
\[
 M=A+\tfrac12Y+\eta.
\]
It produces exactly the same $\Sigma$ and therefore the same centered
Gaussian distribution. In the triangle model, changing $M$ while leaving
$A$'s equation intact changes the mean of $Y$ with derivative one. In the
collider model, intervening on $M$ leaves $Y$ unchanged. The target effects
are different despite complete observed-variable measurement in both models.
Invoking faithfulness selects a model class that excludes the triangle; it
does not make the observable covariance itself reveal which structural world
generated the records \citep{exec14_discovery_design_oracles}.

For a banking interpretation, $A$ may represent standardized offer intensity,
$M$ a subsequent account-balance response and $Y$ net revenue. A direct cost
and an indirect revenue channel can offset each other. The equations are
stipulated and should not be taken as an empirical model of that mechanism.
They explain why an observed zero association can have several meanings.
A reliable chronology establishing $M$ before $Y$ would rule out the second
model as an ordinary acyclic causal explanation. That chronology is additional
information. With only a coarse accounting period, an agent must not assume
that the order of columns supplies it.

Cancellation is also a finite-precision concern. Replace the direct
coefficient by $-1+\delta$. Then
\[
 Y=\delta A+E_M+E_Y,\qquad
 \operatorname{Cov}(A,Y)=\delta,
 \qquad \operatorname{Cov}(A,Y\mid M)=\frac{\delta-1}{2}.
\]
The marginal signal can be arbitrarily small while the conditional signal
remains appreciable. A consistency statement for a fixed nonzero signal does
not guarantee useful uniform recovery as $\delta$ approaches zero. This
explains why the prototype should vary cancellation strength and retain
weak-signal failures. Increasing prompt detail does not increase the numerical
separation between these distributions.

The appropriate remedy depends on the claim. A verified time order can exclude
some candidate graphs. An intervention on the consequential variable can
supply new information. A functional model can use restrictions not present
in conditional-independence information, provided those restrictions are
defensible. A conservative report can retain multiple mechanisms and compare
their decision implications. None of these remedies is equivalent to choosing
an orientation because it sounds natural in a bank narrative. The exact
example supports a test of this distinction without paying for live calls or
obtaining bank data.

\section{A search budget changes which separations can be found}
\label{exec14:l4-depth-example}

The conditioning-set budget has a similarly concrete scientific meaning.
Consider mutually independent unit-variance centered Gaussian disturbances
and the structural equations
\[
 A=E_A,\qquad M_1=A+E_1,\qquad M_2=A+E_2,
 \qquad Y=M_1+M_2+E_Y.
\]
There are two parallel paths from $A$ to $Y$ and no direct edge. Their total
intervention derivative is two. The observed covariance matrix is
\[
 \Sigma_{A,M_1,M_2,Y}=
 \begin{pmatrix}
 1&1&1&2\\1&2&1&3\\1&1&2&3\\2&3&3&7
 \end{pmatrix}.
\]
The Gaussian conditional covariance calculation gives
\[
 \operatorname{Cov}(A,Y)=2,\quad
 \operatorname{Cov}(A,Y\mid M_1)
 =\operatorname{Cov}(A,Y\mid M_2)=\tfrac12,
 \quad \operatorname{Cov}(A,Y\mid M_1,M_2)=0.
\]
For the last equality, the mediator covariance inverse is
$\frac13\left(\begin{smallmatrix}2&-1\\-1&2\end{smallmatrix}\right)$;
multiplying the cross-covariances subtracts exactly two. A search limited to
conditioning sets of size at most one cannot find the separator for this
pair, even with an exact oracle and an unlimited sample. A reported extra
adjacency is therefore a possible consequence of the search restriction,
rather than evidence for a direct offer-to-revenue mechanism
\citep{exec14_discovery_design_oracles}.

This example also distinguishes a discovery query from an adjustment rule.
Conditioning on the mediators helps discover that a direct adjacency is
unnecessary. Adjusting them away when estimating the total effect of $A$
would remove the very paths that define that effect. An agent should not copy
the variables from a separating set into an outcome-regression adjustment
list without reviewing the target and graph. The same named set can be
appropriate for one computational question and inappropriate for another.

The proposed comparison should therefore bind the maximum conditioning order,
the admissible variable sets, the CI test and its missing-result behavior.
Record the number of attempted queries as well as completed ones and preserve
which search bounds were active. A faster configuration with smaller search
depth is a useful alternative when the user accepts its scope, but it is not
automatically a faster implementation of the same inferential procedure.
For banking data, higher conditioning order can also reduce usable observations
and aggravate singularity or mixed-type test limitations. Relaxing one
computational constraint can create another statistical problem.

This small model is a conformance reference for a future integration. It can
check that the adapter preserves query variables, reports the configured
depth and distinguishes a population zero from an unavailable test. The
calculation does not claim that any inspected package has been run or that
the entire search will visit a particular query in a particular order.
Whole-engine qualification still needs a pinned executable path and an
observed trace. The source-inspected Tetrad branch discussed earlier and this
mathematical example answer complementary questions.

\section{Several environments need a causal interpretation of their changes}
\label{exec14:l4-environments}

An assistant may propose using changes across branches, countries or policy
periods to orient a relation. Invariant prediction provides a substantive
alternative to single-environment graph search. The inspected formulation
assumes that, for a suitable predictor subset, the response mechanism and
independent error law remain invariant across environments. In the linear
case, the same coefficients and residual distribution apply in each setting.
Under its stated test and candidate-set assumptions, the intersection of
nonrejected subsets controls false inclusions in that invariant subset; it
does not promise recovery of every cause. Connecting the invariant subset to
causal parents requires the corresponding structural assumptions. In
particular, arbitrary changes to the response or its hidden causes are not
licensed by the presence of an environment label
\citep{petersicp2016}.

The distinction can be exhibited with another exact project law. Index
environments by a known number $m$. Let independent variables satisfy
$U\sim N(m,1)$, $E_A\sim N(0,1)$ and $E_Y\sim N(0,1/2)$, and write
\[
 A=U+E_A,\qquad Y=U-\frac m2+E_Y.
\]
Here $A$ has no causal effect on $Y$. Yet in every environment
\[
 E(A)=m,\quad E(Y)=m/2,\quad
 \operatorname{Var}(A)=2,\quad \operatorname{Cov}(A,Y)=1,
 \quad\operatorname{Var}(Y)=3/2.
\]
Gaussian conditioning gives $Y\mid A,m\sim N(A/2,1)$, independent of $m$.
A stable conditional predictor has emerged from changing a hidden common
cause and compensating the outcome intercept. The different model
$A\sim N(m,2)$, $Y=A/2+\eta$ with independent $\eta\sim N(0,1)$ has
the same observed law in every environment and an intervention derivative
of $1/2$. Thus the environment-indexed observations do not distinguish these
two causal explanations \citep{exec14_discovery_design_oracles}.

This construction does not contradict the invariant-prediction theorem. It
violates the structural interpretation needed to treat the stable observed
regression as the unchanged causal response mechanism: the hidden cause and
the outcome equation change together. It demonstrates why a test of one
observed conditional distribution cannot by itself certify the provenance of
the environment changes. An apparently successful comparison across many
branches can still be uninformative if their changes share this structure.
The number of environments is not a count of independent randomized
interventions on the proposed action.

For a deposit application, a market stress period may change demand, the
bank's pricing rule and measurement practices at the same time. A dated
policy document can establish that pricing changed, but it may not establish
that all other relevant mechanisms remained fixed. A useful language task is
to construct an environment dossier: what changed, when it became effective,
which variables it could affect, and what the sources leave unknown. If the
dossier supports a narrower intervention family, a numerical method can be
assessed under that family. Otherwise the environment comparison remains a
predictive or sensitivity analysis with explicit limits.

The public prototype can test this reasoning affordably. Generate environments
that alter only an admissible predictor mechanism and compare them with the
compensating-change construction above. Keep the target and permitted
documents fixed within each regime, and hide the generator from the assistant.
A successful workflow should distinguish failure to reject invariance from
evidence that the structural assumptions were supplied. A correct refusal in
the second regime and a useful result in the first are different outcomes
that should both appear in the evaluation. This is a proposed study design;
the exact algebra supplies reference cases without supplying performance
estimates for any model.

\section{Allowing mistakes in extracted restrictions without hiding them}
\label{exec14:l4-restriction-errors}

Hard restrictions become especially fragile when an agent extracts many of
them. Let $\mathcal G_D$ be the graph class retained from numerical evidence,
and let $R_1,\ldots,R_J$ be proposed predicates on a graph. The ordinary
intersection retains only graphs satisfying every predicate. A more explicit
sensitivity construction allows at most $k$ predicates to be wrong:
\[
 \mathcal G_{D,k}=\left\{G\in\mathcal G_D:
       \sum_{j=1}^J\{1-R_j(G)\}\leq k\right\}.
\]
This is a proposed deterministic construction, not an implemented general
graph solver. If the true graph lies in $\mathcal G_D$ and violates at most
$k$ proposed restrictions, it lies in $\mathcal G_{D,k}$. Enlarging $k$ cannot
shrink the retained class. Mapping that class into effect sets makes the
decision's dependence on the trusted restrictions visible.

For an exact finite illustration, label four otherwise admissible candidate
models $G_{ab}$ by two binary properties $a,b$. Suppose the extracted claims
assert $a=1$, $b=1$, and that exactly one property is one. The joint
intersection is empty. Allowing one erroneous claim retains $G_{01}$,
$G_{10}$ and $G_{11}$, since each violates exactly one claim; $G_{00}$
violates two and is excluded. A nonempty relaxed class has therefore not
proved that the true model was retained. If the truth is $G_{00}$, the
one-error sensitivity analysis still misses it. The remedy requires either
more trusted evidence or a larger declared error allowance, with the
consequent loss of precision \citep{exec14_discovery_design_oracles}.

A probabilistic guarantee requires actual error control in addition to that
logical construction. Write $N=\sum_j\mathbf1\{R_j(G_\star)=0\}$ for
the number of wrong restrictions at the true graph. If a justified evaluation
or observation law supplies marginal bounds
$\Pr\{R_j(G_\star)=0\}\leq\delta_j$, then
$E(N)\leq\sum_j\delta_j$. Since
$\mathbf1\{N>k\}\leq N/(k+1)$, it follows that
\[
 \Pr\{G_\star\notin\mathcal G_{D,k}\}
 \leq \alpha_D+\frac{\sum_j\delta_j}{k+1},
 \qquad
 \Pr\{G_\star\notin\mathcal G_D\}\leq\alpha_D,
\]
with the right side capped at one. Independence between the restrictions,
or between text and numerical evidence, is unnecessary for this inequality.
Its premises must concern the actual joint workflow, including adaptive
extraction. A model's confidence statement or vote frequency is not a
supplied $\delta_j$. The result is an elementary project deduction exposing
what a reliable probabilistic claim would need.

For three genuine marginal bounds of $.05$ each, the restriction component
is at most $.15$ when $k=0$ and $.075$ when $k=1$. These values are conditional
illustrations; no such language-model error bounds have been established in
this proposal. A larger retained class can make the effect or decision
unresolved even if its retention bound improves. This separates the value of
error tolerance from a promise of useful precision.

The unit of an error allowance needs substantive care. Three copied passages
can express the same mistaken policy statement. Treating them as three
independent confirmations is unjustified, while counting them as three
different restrictions can also distort an at-most-$k$ analysis. The proposed
system should retain a mapping from predicates to source assertions and
source families. An analyst can then define sensitivity to a failed source
family rather than to a number of duplicated sentences. A newly discovered
contradiction should create a revised analysis, preserving the earlier
restrictions and showing which effects or decisions changed.

\section{An intervention's value can be calculated before conducting it}
\label{exec14:l4-finite-information}

The perfect-information calculation above is an upper bound. A prospective
experiment normally supplies noisy information, so a prototype should also
evaluate the attainable value under a declared experiment. Return to the
chapter's two binary worlds. In the forward world, $A$ is a fair binary action
and $Y$ equals $A$ with probability $.9$. In the reverse world, $Y$ is fair
and $A$ equals $Y$ with probability $.9$. Both have the same observational
table. Under a controlled action $A=1$, however, the outcome success
probabilities are $.9$ and $.5$, respectively. This difference comes from
replacing the action equation, not from resampling the old observational
table.

Assume, solely for this planning calculation, that independently intervened
units can be drawn from the same structural world, that outcomes are observed
without loss, and that neither interference nor experimentation changes the
mechanism. If $K$ successes are observed in $n$ probes, let
\[
 f_H(k)=\binom nk(.9)^k(.1)^{n-k},\qquad
 f_L(k)=\binom nk(.5)^n.
\]
For an elicited prior probability $q$ on the forward world, Bayes' rule gives
\[
 q_k=\frac{qf_H(k)}{qf_H(k)+(1-q)f_L(k)}.
\]
The common binomial coefficient cancels in the likelihood ratio, so the
posterior odds equal the prior odds multiplied by
$1.8^k\, .2^{n-k}$. Unlike a prompt-vote fraction, this update has an
explicit observation law and an explicit relationship to the two worlds.
Its validity remains conditional on those two worlds and the stipulated
intervention protocol.

Connect it to the earlier offer decision by assigning gross incremental value
ten times the intervention effect and an offer cost of three. The effects
are $.8$ and zero, giving net offer values five and minus three; no offer
has value zero. For a future comparable account, the optimal action after
observing $k$ is to offer when $q_k>3/8$. Before the probe results are known,
the expected optimal value is
\begin{equation}
 V_n(q)=\sum_{k=0}^n
       \max\{0,\,5qf_H(k)-3(1-q)f_L(k)\}.
 \label{exec14:l4-sample-information-value}
\end{equation}
To obtain this expression, multiply the conditional optimum
$\max(0,8q_k-3)$ by the predictive mass
$qf_H(k)+(1-q)f_L(k)$ and simplify. The resulting sum includes every possible
outcome, including evidence that leads to no offer. There is no favorable
outcome selection.

At $q=1/2$, the no-information value is one. With one probe, a success leads
to an offer and a failure leads to no offer; the expected value becomes
$3/2$, so the expected value of sample information is $1/2$. With two probes,
only two successes justify an offer. The expected value is $33/20$, giving
sample-information value $13/20$. Exact enumeration for larger $n$ appears
in Table~\ref{exec14:l4-evsi-table}. These are arithmetic consequences of a
stipulated finite model, not observed experimental results or measured bank
benefits \citep{exec14_discovery_design_oracles}.



% SOURCE docs/release/2026-10-29/source/docs/writing-repair/2026-10-29/src/fddbb1903e836b29.tex CHAR 0
\begin{table}[htbp]
\centering
\caption{Exact information value under the stipulated two-world intervention law.}
\label{exec14:l4-evsi-table}
\begin{tabular}{rrr}
\toprule
Probes $n$ & Success count triggering an offer & Value of sample information\\
\midrule
0 & $0$ or more & $0.000000$\\
1 & $1$ or more & $0.500000$\\
2 & $2$ or more & $0.650000$\\
5 & $4$ or more & $1.015100$\\
10 & $8$ or more & $1.242492$\\
20 & $15$ or more & $1.440825$\\
50 & $37$ or more & $1.498585$\\
\bottomrule
\end{tabular}
\end{table}


% SOURCE docs/release/2026-10-29/source/docs/writing-repair/2026-10-29/src/9235cdc16df4fd1e.tex CHAR 25715


The value is bounded above by the perfect-information value of $3/2$ per
future account. It cannot decrease when another free, comparable probe is
made available, because a decision rule can ignore that probe. This is a
property of the decision problem and information nesting; it does not imply
that a paid study has nonnegative net value. If $N_f$ future accounts share
the relevant world and reward model, the gross planning value is
$N_f\{V_n(q)-V_0(q)\}$. The complete study cost, including probe treatment,
observation, review, delay and any required harm constraints, must be assessed
against that total. Multiplying by future accounts assumes that the same
learned information transfers to those accounts; heterogeneous worlds would
require a different calculation.

This calculation supplies a clear reason to prefer a generated prototype now.
It can verify the likelihood update, target preservation, action threshold
and complete outcome accounting. It can introduce dropout, shared shocks or
worlds outside the two-model family and test whether the workflow reports
the changed assumptions. An actual bank experiment is justified only when
its feasible action, observation law, transfer population and cost could
resolve a consequential uncertainty. The present calculation cannot supply
those inputs and does not authorize or perform an experiment.

A retrieval-only query has zero value for discriminating these two worlds
if its response law is identical in both and it conveys no new applicable
restriction. Its likelihood ratio is one and the posterior remains $q$.
Retrieval may still improve variable definitions, reduce clerical effort or
reveal a genuine distinguishing fact. Those are separate benefits to model
and measure. Repeating a prompt until it favors the preferred graph does
not reproduce the experiment analyzed in Equation~\eqref{exec14:l4-sample-information-value}.

The comparison also explains why full graph recovery is sometimes unnecessary.
When every retained model recommends the same feasible action over the
relevant cost range, further orientation may have little decision value.
When two models differ on a consequential action, the useful query is the
one that separates them under an admissible information mechanism. An agent
can help formulate that query and show which assumptions drive its value.
The proposal's initial objective is this reviewable connection from discovery
to a decision, with numerical and semantic limits preserved.
\label{exec14:l4-completion-end}


% SOURCE docs/release/2026-10-29/source/docs/writing-repair/2026-10-29/src/5bea1c25666be783.tex CHAR 0
\section{A correlation screen can discard a nonlinear cause}
\label{exec14:l4-nonlinear-screen}

The proposal should also evaluate what happens before the formal discovery
engine is called. A marginal correlation screen may reduce the number of
candidate pairs supplied to an agent, but its omissions constrain everything
downstream. The inspected Augur procedure uses a Spearman screen before
language judgments. Its reported forecasting performance does not establish
that this screen retains all intervention-relevant mechanisms
\citep{exec13_census_augur_arxiv}. An exact project example shows the
distinction without making any claim about an executed Augur benchmark.

Let $X\sim N(0,1)$ and $E\sim N(0,1)$ be independent, and set
\[
 Y=X^2-1+E.
\]
Then $E(Y\mid X=x)=x^2-1$, so $X$ and $Y$ are dependent, but
\[
 \operatorname{Cov}(X,Y)=E(X^3)-E(X)=0,
 \qquad \operatorname{Var}(Y)=E(X^4)-E(X^2)^2+1=3.
\]
The population Spearman correlation is also zero. To see this, the joint law
of $(X,Y)$ is unchanged when $X$ is replaced by $-X$. Since the normal
distribution function obeys $\Phi(-X)=1-\Phi(X)$, the expectation of
$(\Phi(X)-1/2)(F_Y(Y)-1/2)$ is equal to its negative and hence is zero.
The independent Gaussian error makes the outcome continuous, so the usual
continuous-rank interpretation applies. Neither a Pearson nor a Spearman
population threshold bounded away from zero retains this pair on the basis
of its marginal correlation.

The intervention response is not zero: replacing $X$ by $x$ gives
$E\{Y\mid\operatorname{do}(X=x)\}=x^2-1$. Moving from zero to one changes
the mean by one; moving from minus one to one changes it by zero. Thus even
the phrase ``the effect of $X$'' is incomplete without the intervention
contrast. In a banking scenario, $X$ might be a standardized deviation from
a reference pricing point and the response might be curved. The example
does not establish that such a curve applies to deposits. It identifies a
specific failure regime that a discovery proposal should cover
\citep{exec14_discovery_screening_value}.

A deterministic feature $H=X^2-1$ would reveal the association in this law:
$\operatorname{Var}(H)=\operatorname{Cov}(H,Y)=2$. That feature is a useful
model representation, but it does not make $H$ an independently manipulable
business action. It maps $X$ and $-X$ to the same value. The interface must
retain the original action, the transformation and the intervention contrast.
Discovering relations among engineered features can change the graph's
variables and its feasible interventions. An agent should not relabel those
features as independently adjustable levers.

Several remedies have distinct costs. Retaining all candidate pairs avoids
this particular screen but increases search and testing burden. A nonlinear
dependence screen can detect other features of the joint law, subject to its
sampling assumptions and power. A scientifically motivated feature family
can make a restricted functional model plausible, but choosing features after
examining the final evaluation sample changes the procedure being assessed.
A language proposal of a nonlinear relation can be retained as a hypothesis
without being treated as verified evidence. The public prototype should
compare these alternatives with the same initial information and report
which candidate relations were removed before graph search began.

\section{Information gain and decision value can choose different queries}
\label{exec14:l4-information-objective}

Model disagreement or expected entropy reduction can be useful ways to select
a new experiment. The inspected Model Discovery Agent uses a disagreement
surrogate for information gain within its candidate mechanism pool
\citep{v12_la_2608_09696}. The banking decision adds another requirement:
the acquired information must concern something that can change a feasible
action. A high information score need not meet that requirement. The following
project example makes the difference exact.

Let a candidate world have two independent binary properties, $R$ and $S$.
Assume $\Pr(R=1)=1/10$ and $\Pr(S=1)=1/2$. Only $R$ affects the offer
decision: offering gives five units when $R=1$ and minus three when $R=0$;
not offering gives zero. Property $S$ concerns another graph feature that
does not affect this target. The current expected offer value is
$5/10-27/10=-11/5$, so the current action is no offer. Perfect observation
of $R$ would allow offering only in its favorable state, with expected value
$1/2$. Perfect observation of $S$ cannot change the action and has decision
value zero.

Nevertheless, a query revealing $S$ removes one bit of world uncertainty,
while a query revealing $R$ removes only
$h_2(1/10)\simeq .468996$ bits, where
$h_2(p)=-p\log_2p-(1-p)\log_2(1-p)$. Maximizing entropy reduction over
the full candidate world therefore prefers the decision-irrelevant query.
This is a difference between objectives, not a computational failure. It can
arise even when the candidate probabilities and the query likelihoods are
exactly known.

Now consider an informative but imperfect query $Z$ about $R$, with
$\Pr(Z=1\mid R=1)=4/5$ and $\Pr(Z=1\mid R=0)=1/5$. A positive answer
gives posterior probability $4/13$, and a negative answer gives $1/37$.
Both are below the offer threshold $3/8$, so the optimal action remains no
offer for either answer. This query has positive mutual information about
$R$, namely $h_2(13/50)-h_2(1/5)\simeq .104818$ bits, but zero value
for this decision. Information that changes beliefs can still be too weak to
cross a relevant action threshold.

A more accurate query has true-positive probability $19/20$ and
false-positive probability $1/20$. Its positive posterior is $19/28$ and
its negative posterior is $1/172$. Offering only after a positive answer
now gives ex ante value
\[
 5\left(\frac1{10}\frac{19}{20}\right)
 -3\left(\frac9{10}\frac1{20}\right)
 =\frac{17}{50}=.34.
\]
This remains below the perfect-information value $.5$. A query cost of $.1$
per comparable future decision would leave planning net value $.24$ under
these stipulations. The same priced irrelevant query has negative net value.
The probability and cost assumptions are authored inputs, not measured
calibration or vendor prices \citep{exec14_discovery_screening_value}.

The lesson for an agent is precise. Retain a model-discrimination objective
when scientific understanding is the declared target. For a specified banking
action, propagate candidate worlds into utilities and ask whether the query
can change the preferred action. If several bank decisions share the graph,
the utility may include several uses; their populations and weights need to
be stated rather than silently inferred. Customer constraints and acquisition
costs also belong in this calculation. Optimizing information gain is not
wrong, but its connection to investment value must be demonstrated.

The prototype can implement these finite reference problems before conducting
any expensive experiment. It should reproduce the different query rankings,
preserve the posterior update and explain the dependence on the prior and
loss. If a language-only query is assigned a likelihood, that likelihood
requires its own evidence; model confidence is not a substitute. A correct
implementation of this arithmetic would demonstrate decision accounting in
a declared finite problem. It would not establish that the candidate-world
set includes the bank's true mechanism or that a real proposed query has the
stipulated reliability.
\label{exec14:l4-final-end}


% SOURCE docs/release/2026-10-29/source/docs/writing-repair/2026-10-29/src/1b3d2ff5062ac8b4.tex CHAR 0
\hypertarget{exec14:iid-discovery}{%
\section{From language guidance to a defensible discovery
decision}\label{exec14:iid-discovery}}

Language assistance can make a discovery workflow more useful by
clarifying variables, locating applicable knowledge, comparing
algorithms and exposing conflicting assumptions. Its strongest proposal
is a supervised route from meaning to an explicitly restricted causal
model. The decisive question is whether that route improves the
intervention decision while preserving what the observations cannot
distinguish. A graph that is more fluent, more stable under repeated
prompting, or easier to fit need not answer that question.

Consider a bank choosing a one-unit increase in a standardized deposit
offer \(A\) at the start of a month. Let \(Y\) be a standardized measure
of the following month's retained balances. The population consists of
accounts eligible when the offer is set; post-offer behavior cannot
become baseline information by renaming a column. For the example below,
account histories are independent, there is no interference, and the
decision compares \(do(A=1)\) with \(do(A=0)\). These are stipulated
conditions for a calculation, not properties inferred from a real bank.
Shared branches, calendar shocks, selection and equilibrium responses
would require a different observation model.

\hypertarget{a-clock-compatible-ambiguity-that-more-observational-data-cannot-remove}{%
\subsection{A clock-compatible ambiguity that more observational data
cannot
remove}\label{a-clock-compatible-ambiguity-that-more-observational-data-cannot-remove}}

Two centered Gaussian structural models illustrate the information
problem. In the first, the offer changes the outcome:

\[
 A=U_A,\qquad Y=\tfrac12 A+U_Y,\qquad
 U_A\perp U_Y,\quad
 \operatorname{Var}(U_A)=2,\quad\operatorname{Var}(U_Y)=\tfrac32.
\]

In the second, a pre-offer latent condition \(H\) affects both
variables:

\[
 A=H+V,\qquad Y=H+W,\qquad
 H,V,W\ \text{independent }N(0,1).
\]

Both mechanisms respect the recorded order of offer and outcome. Both
produce the same nonsingular observed distribution,

\[
 (A,Y)\sim N\!\left(0,
 \begin{pmatrix}2&1\\1&2\end{pmatrix}\right),
 \qquad E(Y\mid A=a)=\tfrac12a.
\]

The covariance determinant is three, so this ambiguity does not depend
on a degenerate observed table. Nevertheless, the intervention contrast
is \(1/2\) in the first model and zero in the second. Causal sufficiency
would exclude the second model by assumption. The clock alone does not
exclude it. A language system that interprets the phrase ``offer
response'' as proof of a direct effect has supplied a scientific
restriction that was absent from the numerical data.

There is also an exact limit on a forced choice based only on these
observations. If an algorithm selects the first model with probability
\(q\), its success probabilities in the two worlds are \(q\) and
\(1-q\), because the input distribution is identical. Their sum is one
and the smaller is at most one half, for any observational sample size.
This statement includes randomized algorithms. It does not say that all
causal discovery is impossible; it identifies the information missing in
this particular comparison. An applicable assignment restriction or an
intervention could distinguish the worlds. A repeated description of the
same observed association cannot.

This distinction determines the numerical baseline. A PC analysis under
causal sufficiency and the stipulated order answers a conditional
question. A latent-aware analysis must retain the possibility of common
causes. Neither method should be scored as wrong merely for declining to
select a single DAG when its declared model class permits both
mechanisms. For these two stipulated worlds the compatible effect values
are \(\{0,1/2\}\). That pair is not an exhaustive identified set over
arbitrary latent models, and it is not a finite-sample confidence set. A
prototype should retain those qualifications in the returned object.
\citep{exec14_iid_oracle}

\hypertarget{what-the-reported-gains-do-and-do-not-support}{%
\subsection{What the reported gains do and do not
support}\label{what-the-reported-gains-do-and-do-not-support}}

The available studies motivate several distinct product options.
Retrieval can supply missing definitions or restrictions; soft priors
can guide a limited numerical search; statistical prompting can make
existing numerical evidence intelligible; and debate can organize
competing explanations. The appropriate comparator changes with the
role. A retrieved definition is compared with a dictionary and a
documented expert rule, whereas a search prior is compared with the same
numerical algorithm and resource budget.

The primary results also preclude a single claim of uniform improvement.
LACR's retrieved-document condition reduces ASIA skeleton F1 from 1 to
0.73, while additional knowledge improves CORONARY. Its reported perfect
orientation score is conditional on recovered true-positive edges.
Missing and false edges therefore remain outside that particular
denominator. \citep{v12_la_2402_15301} Darvariu and colleagues report
useful combined mutual-information and language priors, especially with
limited search, but mutual information alone performs best on Child in
the reported comparison. Keeping positive search support is useful; it
does not calibrate the probability of a causal effect.
\citep{v12_la_2405_13551}

Statistical prompting is similarly task dependent. In Takayama and
colleagues' Table 3, the Sachs PC baseline has F1 0.44, the language
condition without the statistical prompt has 0.52, and the two
statistical-prompt conditions have 0.29 and 0.27. Other
algorithm/dataset combinations improve. These are comparisons within the
authors' experiment, not local replications or evidence that one
algorithm always dominates. Their lower fit criterion and their
closer-to-reference graph are different outcomes.
\citep{v12_la_2402_01454}

Cohrs and colleagues report query-order agreement of 0.510 and 0.870 for
their two model configurations, and a data-based method has higher SACHS
F1 than either language method. Their answer-proportion null concerns
the response process. Even a correctly calibrated test of that null
would not automatically test conditional independence in the bank
population. \citep{v12_la_2406_07378} The response-count calculation
developed earlier in this chapter supplies a separate finite-binomial
counterexample; it should be reused as a diagnostic, not represented as
a new model experiment.

The newer multi-agent comparison combines algorithm selection and code
execution with graph debate. Its inspected version reports that
additional debate rounds can degrade Survey and Auto results and that
the complete workflow consumes 16--25 times the tokens of standard
baselines. Those findings support a bounded single-round comparator and
full cost accounting, without establishing bank economics.
\citep{v12_la_2407_15073} The 2026 argumentation approach makes
conflicts explicit but retains hard semantic constraints after an
irreversible preliminary skeleton reduction. Its synthetic descriptions
are generated with access to the true graph and its runtime comparison
omits external model latency. These limits motivate neutral
descriptions, misleading-description controls and complete timing.
\citep{exec14_expert_argumentation}

Exposure and semantics require separate tests. The accessible-corpus
study by Feng and colleagues links relation frequency to performance and
examines contradictory statements and context changes. Its scope does
not establish the behavior of every closed model. Moreover, a changed
context may genuinely change the causal claim. A fixed-world paraphrase
test must preserve the population, timing and intervention; a
changed-world test must change its reference answer accordingly.
\citep{v12_la_2407_19638} Contradiction reduction and graph agreement
are useful diagnostics, but neither is a substitute for effect truth.
The inspected hallucination study explicitly retains that distinction.
\citep{exec13_hallucination}

Barakati and colleagues provide a domain application that usefully
exposes the translation problem. Their materials study maps
microscopy-derived unit-cell descriptors to literature terminology,
queries relationships using LangChain and a language model, and reruns
gCastle PC with required and forbidden relations. The stated response
schema includes insufficient information. This is a concrete
architecture for connecting measurements, terminology and candidate
restrictions. Its comparison chiefly discusses agreement between
inferred pathways and prior LiNGAM findings. Agreement does not
establish an intervention oracle, and spatial unit cells are not
automatically independent samples. The abstract uses the term
fine-tuning, whereas the inspected methods describe querying literature
and give no reproducible training recipe or checkpoint. The proposal
should therefore retain the documented workflow without inferring a
trained artifact or calibrated effect uncertainty. The linked repository
is a separate reuse lead requiring its own revision, rights and
execution review. \citep{exec14_iid_materials}

\hypertarget{missing-answers-and-recorded-identity-are-part-of-the-causal-interface}{%
\subsection{Missing answers and recorded identity are part of the causal
interface}\label{missing-answers-and-recorded-identity-are-part-of-the-causal-interface}}

Tracing the pinned research code reveals obligations that are less
visible in a method name. In the selected \texttt{causal\_GPT}
interface, the pure PC wrapper reads variable descriptions and stored
language answers. Its numerical \texttt{pvalue} method pools the two
query orders, then converts a decision to zero or one. A missing stored
row returns zero. The advertised weighted-vote branch also returns zero
beside an unfinished implementation marker. Thus unavailable evidence,
an unsupported aggregation method and a dependence decision can share
one numerical representation. These observations concern the inspected
source, not measured frequencies in an installed package.
\citep{exec14_iid_gpt}

The hybrid path contains a genuine numerical fallback for queries
outside the selected language-variable/conditioning scope. Within that
scope it collects unanswered queries, uses temporary decisions while
assembling them, and requests their answers in the surrounding driver.
That driver also saves result metadata with fixed model and temperature
strings despite accepting different values as arguments. For example, a
configured model B at temperature 0.2 would still receive the inspected
record's fixed model label and temperature 0.6. This is a literal
assignment trace, not an API call. A comparison cannot be reproduced
reliably if result identity and actual configuration can diverge.
\citep{exec14_iid_gpt}

The proposal should therefore adopt the useful separation of language
and numerical query channels while replacing the ambiguous interchange.
A language-answer object needs the target statement, operational
definitions, source evidence, model revision, settings, parsing status
and an explicit unknown state. A statistical test object needs its
observations, null, sampling law and calibration. A pending query must
not be published as a completed dependence judgment. Unsupported
weighted voting should refuse until implemented and checked. Result
identity should be populated from the actual request and response, with
missing metadata retained as missing.

\hypertarget{a-token-score-does-not-have-the-same-meaning-in-three-discovery-engines}{%
\subsection{A token score does not have the same meaning in three
discovery
engines}\label{a-token-score-does-not-have-the-same-meaning-in-three-discovery-engines}}

The five pinned LLM-and-SCD notebooks expose a complete research route
from an initial graph and bootstrap summaries through two prompts to
constrained rediscovery. In a selected collection cell, the code sums
the first token's top-five probability mass only when the token is
exactly \texttt{Yes} or \texttt{yes}. It separately accumulates
\texttt{No} or \texttt{no}. It does not identify all semantically
equivalent strings or turn the unreported mass into an explicit
uncertainty object. The resulting score is a property of the prompted
token distribution, not a probability that an edge exists in the bank.
\citep{exec14_iid_scd}

An exact authored example makes the parsing boundary visible. Suppose
the five tokens have masses 0.96, 0.02, 0.01, 0.005 and 0.005, assigned
to \texttt{Yes}, \texttt{No}, a newline, \texttt{I} and \texttt{It}. The
literal affirmative score is 0.96. Change only the first token to a
leading-space \texttt{Yes}; the exact-string collector assigns
affirmative score zero, although a whitespace-normalized semantic class
still has mass 0.96. These are constructed input distributions, not
observed endpoint outputs. They show why unrecognized formats and
truncated token lists need explicit treatment before any scientific
constraint is created. Simple normalization would fix this example but
would not establish causal calibration or solve multi-token answer
semantics. \citep{exec14_iid_oracle}

The conversion then magnifies a parsing choice. The selected notebook
marks a score below 0.05 as forbidden, above 0.95 as required, and the
interval between them as unknown. Consequently the two constructed
representations above become opposite hard constraints. At exactly 0.95
the notebook returns unknown, while the inspected paper's Algorithm 2
uses an inclusive upper threshold. The 0.05 boundary is unknown in both.
The notebook commit predates the inspected paper revision, so the
discrepancy is a version-alignment obligation, not evidence about the
authors' later experimental execution.
\citep{exec14_iid_scd,v12_la_2402_01454}

The target backend introduces another semantic choice. The PC conversion
creates required and forbidden directed-edge objects; exact search
receives an allowed supergraph, in which a high score permits rather
than forces an edge. The proposed pinned DirectLiNGAM target documents
directed-path prior knowledge. Its meaning must be checked against the
notebook's direct-impact question and the chosen matrix convention. For
a chain \(A\to M\to Y\), no direct \(A\to Y\) edge exists, but an
\(A\)-to-\(Y\) directed path does. Converting ``no direct impact'' to
``no directed path'' rules out this chain and its mediated effect. A
common matrix with numbers \(-1,0,1\) does not make these three
contracts interchangeable. \citep{exec14_iid_scd,v13_sw_lingam}

Notebook dependency installation is not locked to a reproducible
environment. The separate cycle-conversion notebook also passes the data
positionally into a constructor whose first parameter is
\texttt{random\_state} in the later inspected target, before calling
\texttt{fit} with the data. Compatibility must be established at an
actual chosen dependency revision; this source comparison does not
demonstrate what ran in the historical notebook. Acyclic conversion and
choosing among repaired constraints by fit on the same observations are
further parts of the adaptive procedure, and must accompany any later
uncertainty claim. \citep{exec14_iid_scd,v13_sw_lingam}

LLM-CD offers a useful alternative representation. The inspected
evaluation checks generated relations against directed reachability. Its
backend wrappers serialize ancestral constraints separately from absent
edges and, for one backend, order constraints. This is consistent with
the primary paper's distinction between ancestral priors and local
adjacency. The paper also states that the factorization of path priors
is approximate and that their common strength is problem dependent.
\citep{exec14_iid_cd,v12_la_2306_16902} Adoption still needs a bounded
adapter: the selected wrappers use shared working paths and shell calls
whose exit status is not checked before parsing, then remove
intermediate files. A nested backend license does not resolve missing
root-project terms. The proposed adapter should retain typed
constraints, use isolated run directories, check process results and
preserve all scientific inputs and outputs. Its terms and dependency
boundaries must be resolved before enabling it.

\hypertarget{a-public-and-generated-prototype-with-a-decision-criterion}{%
\subsection{A public and generated prototype with a decision
criterion}\label{a-public-and-generated-prototype-with-a-decision-criterion}}

A useful first prototype has two stages. Deterministic semantic fixtures
check missing answers, unsupported methods, token variants, exact
threshold boundaries, graph endpoint preservation, cycles,
edge-versus-path conversion and request/result identity. These tests are
cheap and have known expected answers. They establish an adapter
contract; they do not estimate live model quality. Only then should the
team profile a matched discovery experiment and register its final
design.

The experiment should split by mechanism family. Repeated numerical
samples, variable names, retrieved packets and prompts belong inside a
mechanism's cluster. The Gaussian pair above is a compulsory ambiguity
case, accompanied by causally sufficient, non-Gaussian and latent cases
with explicit clocks. Generator truth and effect oracles remain outside
the prompts. Natural descriptions, neutral renamed descriptions and
deliberately misleading descriptions are crossed with no retrieval,
applicable original evidence, wrong-domain evidence and contradiction.
Removing meaningful definitions changes the information supplied, so the
comparison must say when it tests robustness and when it tests the value
of information.

Each language arm needs the same numerical inputs and a declared
comparator: a dictionary/clock baseline; PC or FCI under the appropriate
assumptions; score or non-Gaussian methods within their qualified
regimes; language-only proposals; and soft or hard guidance with
declared strength and refusal rules. Final graphs alone are
insufficient. Record class and endpoint retention, invalid outputs,
valid abstentions, conditional effect-set retention and width, and
downstream action regret under known truth. Include every requested case
in the denominator, including parse failures, timeouts, empty outputs
and conflicting constraints. Measure total calls, tokens, retries,
numerical resources, latency and adjudication effort.

Before final runs, freeze versions, rights, splits, targets, sample
sizes, replication counts, tolerances, primary paired comparisons and
multiplicity treatment. Development profiling may set resource caps; it
cannot become an unannounced final experiment. Prompt variants and prior
strengths must not be selected on the final mechanisms. An optional live
model experiment needs the chosen endpoint's actual revision and
measured costs, but none of these design steps requires bank data or a
paid experiment now.

The banking example supplies a transparent economic threshold. Stipulate
USD 200 per outcome unit and USD 60 for the unit offer increase. Net
values are USD 40 in the direct-effect world and minus USD 60 in the
latent world. If a planner assigns probability \(p\) to the former, the
expected net action value is \(100p-60\). Because the observed
likelihood is the same, additional observational records alone cannot
update those odds within this two-world comparison. At \(p=0.9\) the
expected value is USD 30, while perfect knowledge would produce USD 36
by taking the action only in the favorable world. Thus USD 6 per
decision is an upper bound on the expected value of any information
experiment under those assumptions. At \(p=0.5\) it is USD 20. A real
noisy experiment supplies less information and incurs costs; neither
number is bank ROI or an authorized spending budget.
\citep{exec14_iid_oracle}

The resulting recommendation is to reuse typed evidence, explicit
conflict handling and bounded search guidance as design references; to
defer enabled package integration until the specified contracts are
verified; and to reject automatic conversion of response agreement into
causal confidence. The unresolved scientific task is to justify useful
graph/effect uncertainty for the complete adaptive procedure in declared
regimes. The small-language-model discovery lead currently has metadata
but no acquired original; its methods cannot support a technical
comparison yet. Broader literature coverage and that original remain
explicit survey work. They do not prevent the proposed public/generated
prototype from testing the concrete semantic and decision boundaries
established here.

\label{exec14:iid-discovery-end}


% SOURCE docs/release/2026-10-29/source/docs/writing-repair/2026-10-29/src/5e1bb6d4d921349d.tex CHAR 0
\section{Expert-guided graph ensembles and their intervention limits}
\label{exec20:expert-scope}

A bank building a diagram from several discovery algorithms will often face
agreement about a relationship but disagreement about its direction. The
methods may also miss different relationships because they use different
scores, tests and structural restrictions. Asking a domain expert to resolve
selected disagreements is a plausible use of an assistant. The design question
is what the expert is being asked, how its answer changes the graph, and whether
the resulting graph supports the bank's intervention query.

Tench and Demeester propose a dynamic expert-guided ensemble that distinguishes
adjacency from orientation. It counts how many component graphs contain any
edge between a pair, considers pairs in descending frequency, and adds directed
edges subject to an acyclicity check. A lower adjacency threshold determines
which pairs are eligible. For eligible pairs without a strict majority, it asks
the expert whether a plausible causal relationship exists. A separate threshold
controls when direction votes suffice and when expert orientation is requested.
The procedure uses equal component weights and a finite query budget. This is
greedy graph construction; the term model averaging does not establish a
Bayesian posterior over causal models. \citep{exec20_expert_averaging}

\subsection{An expert's causal ordering is not a direct edge}
\label{exec20:expert-semantics}

The source deliberately permits an expert to answer about a direct or indirect
causal path. This is weaker than knowing whether two nodes are adjacent in the
full graph. In a stipulated offer--uptake--balance chain
$O\rightarrow U\rightarrow Y$, changing the offer may affect balance through
uptake even though $O\rightarrow Y$ is not a direct edge. An expert can answer
the path question correctly while accepting a spurious direct connection
proposed by a component method. The paper recognizes this distinction in its
simulated expert and compares it with an expert that knows direct edges.
\citep{exec20_expert_averaging,exec20_expert_averaging_reference}

The proposed bank interface should therefore carry the proposition being
elicited: adjacency, ancestry, orientation conditional on adjacency, or a
temporal exclusion. It should also permit insufficient knowledge. A reply about
plausible ancestry cannot silently become proof of a direct edge, and a forced
binary answer cannot represent every legitimate uncertainty state. The set of
variables matters: marginalizing uptake changes the observed-variable graph
even when the underlying mechanism is unchanged.

The method's output is a DAG. Its component comparison includes FCI, while its
noisy benchmarks can contain latent confounding. A directed output cannot by
itself retain every bidirected or unresolved endpoint in a latent graph. Before
an adapter is proposed, the conversion from each component's graph type must be
explicit, with unsupported endpoints preserved or refused. Acyclicity checks
alone do not establish compatibility with an input equivalence class or identify
an intervention in the presence of hidden common causes.
\citep{exec20_expert_averaging}

\subsection{Frequency, diversity and structural compatibility}
\label{exec20:expert-votes}

For one binary question, suppose eleven voters each have probability $3/4$ of
being correct. Under independent errors their majority is correct with
probability
\begin{equation}
 \sum_{k=6}^{11}\binom{11}{k}
       \left(\frac34\right)^k\left(\frac14\right)^{11-k}
 =\frac{1012581}{1048576}\simeq0.966.
 \label{exec20:expert-majority}
\end{equation}
If all eleven instead copy the same $3/4$-correct signal, majority correctness
remains $3/4$. These are two authored probability laws with the same individual
accuracy. They show why individual competence alone cannot justify improvement
toward certainty. The paper itself acknowledges that algorithms share data and
often share statistical principles. Its discussion motivates diversity, but
does not estimate a calibrated graph probability from vote frequency.
\citep{exec20_expert_averaging,exec20_expert_averaging_reference}

Even individually valid DAGs can produce incompatible edge majorities. Consider
the three graphs
\begin{equation}
 \begin{split}
 G_1&=\{A\rightarrow B,\ B\rightarrow C\},\\
 G_2&=\{B\rightarrow C,\ C\rightarrow A\},\\
 G_3&=\{C\rightarrow A,\ A\rightarrow B\}.
 \end{split}
 \label{exec20:expert-cycle}
\end{equation}
Each is acyclic, yet every edge of $A\rightarrow B\rightarrow C\rightarrow A$
has frequency $2/3$. A greedy acyclicity rule must discard or reconsider an
edge. Its result can depend on tie handling and earlier choices. Returning one
DAG makes the output structurally admissible under that rule, but it neither
preserves every component claim nor proves which disputed edge is false.
\citep{exec20_expert_averaging_reference}

\subsection{The reported evidence is informative and conditional}
\label{exec20:expert-evidence}

The main experiments use seven benchmark network structures with synthetic
categorical samples, clean and corrupted versions, four sample sizes from 100
to 100,000, and twenty column orderings. The principal simulated expert is
correct on 70\% of existence questions and 90\% of orientation questions.
Thresholds of zero for adjacency and 0.9 for orientation, with a budget of twice
the node count, were selected using separate randomly generated networks. The
paper reports broad gains in its graph scores, alongside a precision--recall
tradeoff and varying comparison denominators when algorithms fail or return
undefined precision. These are reported experiments; they are not newly
reproduced results in this proposal. \citep{exec20_expert_averaging}

The appendix tests a narrower real-model setting on five clean networks,
using \texttt{gpt-5-nano-}\allowbreak\texttt{2025-08-07} and excluding Pathfinder
because its variable names lack the needed semantics and Formed because many
edges represent diagnostic rather than causal relations. Expert responses are
cached across seeds and sample sizes. Both argument orders are queried, with
disagreement causing a fallback to component voting. The reported per-response
orientation and existence accuracies are not the accuracy of that combined
fallback policy. \citep{exec20_expert_averaging}

\begin{center}
\begin{tabular}{lrr}
\hline
Reported case & Orientation accuracy & Existence accuracy\\
\hline
Mean across five networks & 93.3\% & 66.4\%\\
Property & 83.9\% & 43.9\%\\
Alarm & 82.8\% & 52.2\%\\
\hline
\end{tabular}
\end{center}
Property's existence accuracy falls below the paper's better-than-random expert
assumption. A useful mean therefore cannot replace domain-specific evidence,
and accuracy on all possible pairs need not equal accuracy on the difficult
pairs selected adaptively by an ensemble. Cached responses also do not become
twenty independent expert realizations because the numerical algorithms are
rerun twenty times. The appendix's near reconstruction of the Asia network
raises a memorization concern; unsuccessful recall of another full network
cannot establish that its concepts or fragments were absent from training.
\citep{exec20_expert_averaging}

The source states that it supplies no formal graph-recovery guarantee and has
not empirically qualified continuous or mixed observations. Its cost discussion
includes component discovery, expert requests and structural-distance scoring.
The proposed comparison should likewise retain timeout, memory and invalid-
output denominators, distinguishing a conditional score on successful outputs
from all-request performance. The source's proposed imputation sensitivity
analysis does not supply the missing runtime results or a bank operating cost.

\subsection{Resolve the printed query contract before reuse}
\label{exec20:expert-budget}

The prose says that a fixed budget bounds queries and exhausted decisions fall
back to majority voting. The printed Algorithm~1 leaves two details unresolved.
For an existence query, the counter increment lies inside the branch reached
only when the expert accepts the connection. With three eligible minority
connections, budget one, and three rejections, a literal trace makes three
expert calls while the printed counter remains zero. Separately, when an
accepted adjacency has two legal directions, neither reaches the orientation
threshold, and no budget remains, the printed conditional chain has no final
branch adding the majority direction. An example with direction votes
$3/4$ and $1/4$, threshold $9/10$ and budget zero therefore leaves the edge
unassigned under the printed chain. \citep{exec20_expert_averaging_reference}

These are specification discrepancies, not verified defects in the released
implementation. Resolving them requires an immutable source version and its
tests and call logs. A proposed adapter should count every attempted external
call, distinguish logical questions from the two argument-order requests,
account for retries and invalid answers, and define the behavior at exhaustion.
Caching needs a key that binds the variable meanings, full request, model
snapshot and relevant settings. Otherwise a cost-saving cache can silently
reuse evidence for a different scientific question. The method and data lead
in the paper are useful acquisition routes; package names and reported license
labels are not sufficient permission or provenance for direct adoption.

\subsection{A selected graph still has to answer the intervention query}
\label{exec20:expert-identification}

For a separate exact example, consider two standardized observational summaries
$X,Y$ whose temporal ordering has not been established. Let
$\rho=3/5$. The following two zero-mean Gaussian structural models use
independent errors within each model:
\begin{equation}
 \begin{aligned}
 \mathcal M_1:\quad &X=\epsilon_X,\quad Y=\rho X+\epsilon_Y,\\
 &\operatorname{Var}(\epsilon_X)=1,\quad
 \operatorname{Var}(\epsilon_Y)=1-\rho^2,\\[2pt]
 \mathcal M_2:\quad &Y=\eta_Y,\quad X=\rho Y+\eta_X,\\
 &\operatorname{Var}(\eta_Y)=1,\quad
 \operatorname{Var}(\eta_X)=1-\rho^2.
 \end{aligned}
 \label{exec20:expert-worlds}
\end{equation}
Both imply covariance
$\left(\begin{smallmatrix}1&3/5\\3/5&1\end{smallmatrix}\right)$ and hence the
same observational distribution. Their effects on the mean of $Y$ of increasing
the intervention level of $X$ by one unit differ:
\begin{equation}
 \tau(\mathcal M_1)=\frac35,\qquad \tau(\mathcal M_2)=0.
 \label{exec20:expert-effects}
\end{equation}
More algorithms fitted to the same distribution do not identify the missing
direction. Truthful temporal restrictions, suitable experiments or independently
justified structural knowledge can eliminate a candidate; that extra information
must be stated. The example does not override a known banking information clock.
\citep{exec20_expert_averaging_reference}

Suppose a one-unit increase in $Y$ is provisionally valued at USD~100 and the
intervention costs USD~40. The two compatible margins are USD~20 and USD~$-40$.
Treating a 90\% graph vote as probability for $\mathcal M_1$ gives effect
$27/50$ and margin USD~14, but the vote has supplied no probabilistic calibration
that warrants that calculation as expected value. A decision under unresolved
model ambiguity must retain the compatible effects or adopt an explicitly
justified probability model. This is an authored planning illustration, not
evidence of model performance or bank returns. \citep{exec20_expert_averaging_reference}

The proposal should therefore retain expert-guided averaging as a candidate
for generating and reviewing diagrams. A public prototype comparison can use
known generated mechanisms, withheld domains, correlated component failures,
wrong expert assertions, budget exhaustion and incompatible graph endpoints.
It should compare static aggregation, selective expert queries, a numerical
method without expert input and a workflow that preserves unresolved graph
classes. The outcomes should include the correctness of the intervention
answer, supported refusal, retained ambiguity, all physical requests and total
review cost. Graph scores remain useful diagnostics, but qualification of an
intervention service requires its own evidence.



% SOURCE docs/release/2026-10-29/source/docs/writing-repair/2026-10-29/src/4f95c57067806217.tex CHAR 0
\section{Domain knowledge as evidence for a causal graph}
\label{exec22:domain-scope}

A collections manager considering a payment holiday needs to distinguish a
plausible story from an intervention that improves customer outcomes. A language
model may suggest that financial strain affects delinquency, that servicing
changes financial strain, and that earlier delinquency affects the servicing
offer. Each statement needs a population, a time index, a variable definition
and an evidential basis. Joining the three statements without those qualifications
can create a cycle whose meaning is unclear. A retrieval system can help locate
the underlying evidence, but a retrieved sentence does not by itself establish
an invariant structural mechanism. Two recent clinical studies make the promise
and the limitations of this assistance concrete.

\subsection{Graph elicitation and the meaning of completeness}
\label{exec22:domain-elicitation}

Aziz and Brookhart study a hypothetical comparison of statins with no treatment
for preventing cardiovascular disease. They run four prompting strategies on
GPT-4o and a zero-shot strategy on GPT-o1, with three separate conversations per
configuration. Two author reviewers assess consistency, acyclicity, source
accuracy, completeness and prompt adherence. This is a study of fifteen outputs
in one familiar clinical scenario, rather than an estimate of performance over
a population of banking questions. The original was published on 30 October
2025. \cite{exec22_domain_aziz}

Its completeness criterion asks whether the graph contains seven groups of
variables drawn from an ASCVD risk score. This is a useful check for conspicuous
omissions, but it does not validate the graph's arrows or establish that its
adjustment set identifies the statin effect. Predictors can have different causal
roles: a pretreatment common cause, a mediator and a selection variable cannot
be treated interchangeably simply because each predicts the outcome. In a bank,
a strong arrears predictor measured after a servicing offer may be an outcome
of that offer. Including it in a list is different from correctly deciding
whether to adjust for it.

Table 1 and the article's accompanying text report fourteen acyclic graphs and
seven graphs containing six of the seven criteria. The seven-of-fifteen count
does not support the abstract's ``more than half'' wording. Race is absent from
all fifteen outputs. That is a reproducible observation about the study's
criterion; it does not establish a universal requirement to insert a particular
attribute into every causal model. A social attribute's role needs an explicit
account of the processes it represents and the target population. Blindly
adding a variable to improve a checklist score could conceal the same scientific
problem as omitting it. \cite{exec22_domain_aziz,exec22_domain_reference}

The study also separates producing citations from producing reliable evidence.
One-shot outputs supply sources but include nonexistent references and incorrect
titles, dates or author details. Chain-of-thought outputs give explanations
without sources in all three runs. The GPT-o1 zero-shot configuration consistently
supplies references and code, yet one code output needs manual repair. Its third
run is described as supplying nine sources and then assessed using denominators
of six; the original outputs are needed before aggregating that run's citation
accuracy. The source assessments described in the paper use title and abstract
review, which is weaker than checking that primary methods and results support
a particular causal arrow. \cite{exec22_domain_aziz}

One instruction-based output contains arrows in both directions between statin
use and cholesterol. The authors correctly identify the need for temporally
indexed nodes if feedback is intended. The banking counterpart is
\(D_t\to A_t\to D_{t+1}\), where \(D_t\) is delinquency information available
before action \(A_t\). Collapsing \(D_t\) and \(D_{t+1}\) into a single
``delinquency'' node loses the information clock. A graph assistant should request
that missing distinction rather than silently reverse an arrow to satisfy a
syntactic check. The study supports a brainstorming and source-checking role;
it supplies no validation of a bank effect estimator.

\subsection{Retrieval helps under a particular reference and query}
\label{exec22:domain-retrieval}

Lin and colleagues combine language models, retrieval and observational discovery
for chronic low back pain. Their February 2026 medRxiv original is a preprint,
not a peer-reviewed result. They use baseline observations from 1,868 participants
and a fourteen-variable expert reference. Crucially, the authors explicitly state
that this reference contains cycles and is not a DAG. The reference was developed
from clinical expertise and a literature synthesis; it is a declared comparison
target, not observed intervention truth. \cite{exec22_domain_lin}

The three knowledge systems share Mistral 7B-Instruct v0.3. Ordinary retrieval
supplies document chunks selected by vector similarity. Knowledge-graph retrieval
extracts entities and relations, merges entities, forms communities and retrieves
entities, attached relations, community summaries and supporting text. The paper
describes 600-token chunks with 100-token overlap, an 8,000-token retrieval budget,
and a pipeline involving constrained extraction, a graph database and community
clustering. These are descriptions of the authors' system; they do not establish
that its dependencies are pinned, licensed for the proposed distribution or
operationally integrated here. \cite{exec22_domain_lin}

Both retrieval systems use the same source articles that the experts used to
construct the comparison graph. This is an informative test of reconstructing
an aligned body of knowledge. It does not measure performance when the literature
is incomplete, contradictory, selected by a different team, or independent of the
reference adjudicators. The authors discuss this alignment explicitly. Their
argument that the model predates publication of the expert graph also does not
by itself authenticate the model's training corpus or exclude related earlier
content. The appropriate extension is a deliberately held-out reference and
document design, not a declaration that contamination has been ruled out.

For the data component, the paper describes hill climbing with local generalized
linear models, BIC scores and 500 bootstrap samples. A thresholded consensus graph
forms the initial model. The system then combines its edges with the knowledge
system's proposals and asks a further direction question when both directions
are suggested. Plausibility, association, temporality and a direct causal query
are separate prompt conditions. Their relative performance matters as much as
the strongest result. \cite{exec22_domain_lin}

\begin{center}
\begin{tabular}{p{0.48\linewidth}rr}
\hline
Reported configuration & F1 & Difference \\
\hline
Observational initial model & 0.396 & --- \\
Initial model plus LLM, plausibility & 0.636 & --- \\
Initial model plus ordinary retrieval, plausibility & 0.714 & --- \\
Initial model plus graph retrieval, plausibility & 0.745 & 0.031 over 0.714 \\
Graph retrieval alone after direction resolution & 0.737 & 0.008 below 0.745 \\
\hline
\end{tabular}
\end{center}

These are displayed scores against the single expert graph. The improvement
from 0.396 to 0.745 is 0.349, but the extra gain over ordinary retrieval is 0.031,
and the gain from adding observational discovery to the strongest graph-retrieval
configuration is 0.008. None is a paired confidence interval or a measured return
on engineering effort. Before direction resolution, graph retrieval's plausibility
F1 is 0.602, below the LLM-only value of 0.627. After resolution the values are
0.737 and 0.636. In the final combined systems, ordinary retrieval exceeds graph
retrieval for association by 0.060 and for temporality by 0.130. A blanket ranking
of retrieval architectures would conceal these reversals. \cite{exec22_domain_lin,exec22_domain_reference}

The strongest combined configuration still reports FDR 0.360, meaning disagreement
among its predicted directed edges against this reference under the paper's
definition. It is not a certified probability that a proposed bank causal claim
is false. Small differences also deserve care: subtracting the printed 0.396
from 0.636 gives 0.240, whereas Table 5 prints a change of 0.241. Rounding of
unreported underlying values can explain that discrepancy. It supplies a reason
to recover the original edge records, not a reason to invent corrected results.

\subsection{Direct edges, indirect effects and feedback need different objects}
\label{exec22:domain-semantics}

The chronic-pain plausibility and temporality questions explicitly allow a
relationship ``through one or more intermediate variables.'' Their answers
therefore need not describe direct adjacency among the measured variables.
For the chain
\begin{equation}
 A\longrightarrow M\longrightarrow Y,
 \qquad E=\{(A,M),(M,Y)\},
 \label{exec22:domain-chain}
\end{equation}
\(A\) is an ancestor of \(Y\) even though \((A,Y)\notin E\). A positive
indirect-effect answer cannot be inserted as a direct edge without changing the
graph and possibly its conditional independences. Conversely, when a variable
is unmeasured, a projected graph may legitimately have different edge semantics.
The interface must state which object it is constructing. The paper's source
text is sufficient to identify the mismatch that needs checking; it is not
sufficient to reconstruct every edge's intended semantics. \cite{exec22_domain_lin,exec22_domain_reference}

Taking a union of proposed edges also has no universal accuracy guarantee.
Starting from the correct two-edge set in equation~\eqref{exec22:domain-chain},
adding the false edge \(Y\to A\) produces two true positives, one false positive
and no false negatives. Thus
\begin{equation}
 F_1=\frac{2\,\mathrm{TP}}{2\,\mathrm{TP}+\mathrm{FP}+\mathrm{FN}}
     =\frac45<1.
 \label{exec22:domain-f1}
\end{equation}
It also creates a directed three-cycle without any reciprocal pair. Resolving
only two-way pairs would not remove it. This exact example does not assert that
the paper's code necessarily produces such a failure. It shows why the observed
nonnegative changes in its result table cannot be promoted to a theorem that
combining knowledge and data always improves the graph. \cite{exec22_domain_reference}

There is a further distinction in the observational comparison. Moralization
drops arrowheads and connects parents of a common child. For
\(A\to M\leftarrow B\), it adds the \(A\)--\(B\) link even though that link
is absent from the original skeleton. Consequently, the paper's moralized-graph
comparison is not simply a comparison of adjacencies after ignoring direction.
Good agreement on that object cannot identify a unique causal ordering. Similarly,
a cyclic expert reference and a DAG learner require an explicit account of the
representation mismatch before their distance is interpreted as scientific error.

A useful banking representation retains the proposition type, time horizon,
population, supporting span and evidence strength for each proposed relation.
It distinguishes a direct cause from an ancestor, an association, an ontological
relationship and a modeler's assumption. A type hierarchy or a semantic link in
a knowledge graph is not automatically a causal edge. Hard constraints can narrow
the graph search when justified by an information clock or design rule. Uncertain
literature claims should remain revisable evidence; treating them as inviolable
constraints can exclude the true graph. When sources conflict, returning the
alternatives preserves information that forced orientation destroys.

\subsection{Identical records can support opposite investment decisions}
\label{exec22:domain-value}

An exact illustration explains why better graph agreement alone is insufficient
for a product decision. Let \(A\) denote a centered, standardized servicing
intensity and \(L\) a centered loss index, where one loss unit is stipulated to
equal USD~100. Consider two authored Gaussian structural models. In the first,
all exogenous variables are independent and
\begin{equation}
 A=U_A,\qquad L=-\tfrac12 A+\varepsilon,
 \qquad \operatorname{Var}(U_A)=1,\quad
 \operatorname{Var}(\varepsilon)=\tfrac34.
 \label{exec22:domain-causal}
\end{equation}
In the second, \(U,\eta_A,\eta_L\) are mutually independent centered Gaussians,
each with variance \(1/2\), and
\begin{equation}
 A=U+\eta_A,\qquad L=-U+\eta_L.
 \label{exec22:domain-latent}
\end{equation}
Both models imply exactly the same centered bivariate Gaussian observed law:
\begin{equation}
 \operatorname{Var}\!\begin{pmatrix}A\\L\end{pmatrix}=
 \begin{pmatrix}1&-1/2\\-1/2&1\end{pmatrix}.
 \label{exec22:domain-covariance}
\end{equation}
In the first model, replacing the action equation
by a one-unit higher intervention lowers expected loss by one half unit. In the
second, replacing the action equation leaves \(U\) and \(L\) unchanged, so the
effect is zero. The observed regression slope is \(-1/2\) in both. No amount
of replication of this observed joint law distinguishes the two intervention
effects. The Gaussian loss index is an illustrative local model, not a claim
that realized bank losses can have an unrestricted negative support.
\cite{exec22_domain_reference}

Suppose the higher action costs USD~20 per eligible customer over the declared
horizon, with no other monetized effects. Its modeled incremental margin is
\begin{equation}
 V=-100\{\mathbb E[L\mid\operatorname{do}(A=a+1)]
             -\mathbb E[L\mid\operatorname{do}(A=a)]\}-20,
 \label{exec22:domain-margin}
\end{equation}
which equals USD~30 in the first model and USD~\(-20\) in the second. Retaining
both models therefore retains the range \([-20,30]\). Under a separately declared
subjective weight \(w\) on the first model, expected margin is \(50w-20\),
positive only when \(w>2/5\). A language-model vote fraction, retrieval score or
graph F1 is not that probability. The practical task is to find evidence that
can distinguish the models: an admissible randomization, a justified additional
measurement, or an independently supported restriction. More fluent explanations
of the same observed correlation cannot do so. \cite{exec22_domain_reference}

\subsection{A public evaluation and a concrete integration decision}
\label{exec22:domain-validation}

The proposal should first test graph assistance on public documents and generated
cases with declared causal laws. The evaluation needs two distinct targets.
One is reconstruction of a specified expert account, for which alternative
adjudicators and documented disagreements are legitimate evidence. The other is
recovery of causal implications under a known data-generating law. Conflating
these targets would reward agreement with an expert even when the expert's graph
is incomplete. The Gaussian pair above provides an inexpensive refusal case:
a system given only the shared observed law should retain both effects unless
the task supplies additional identifying information.

Cases should also include indirect-only relations, long cycles, feedback requiring
time expansion, unmeasured common causes, contradictory papers and evidence that
is silent about direction. Keep the adjudication documents separate from the
retrieval corpus for the primary test; report same-corpus reconstruction as a
different condition. Hold out domains as well as individual documents. Match
source access, context allowance and reviewer effort across data-only,
expert-only, ontology-assisted, LLM-only, ordinary-retrieval, graph-retrieval and
combined systems. This tests whether constructing a knowledge graph is worth its
additional extraction, maintenance and review cost.

Measure source identity and span entailment, variable/time alignment, direct-edge
and ancestral accuracy, equivalence-class retention, abstention, valid
identification and decision consequences separately. Repeat prompts and seeds;
do not treat correlated edge votes from a shared corpus as independent experts.
Record graph-building tokens, query tokens, latency, correction effort and
maintenance alongside numerical scores. An 8,000-token context cap does not
equalize the cost of building the two retrieval systems. The public study should
predeclare thresholds for useful assistance and costs before inspecting the
held-out results. There is no measured financial benefit in the two clinical
papers that can set those thresholds for a bank.

For implementation planning, adopt the evidence-provenance pattern and the
separation of variable definition, relation proposal and graph validation.
Integrate retrieval behind a typed interface only after pinning the underlying
corpus, software, model and rights. Retain the ordinary-retrieval comparator;
the results do not justify porting an entire graph-retrieval stack in advance
of a matched benefit test. Reject unconditional union, compulsory orientation
and silent conversion of an indirect relation into a direct edge. The chronic-pain
appendix's Boolean answer schema and fallback to general knowledge lack an
explicit abstention channel; the proposed interface should supply one.
\cite{exec22_domain_lin}

The same reviewed request can pass through a library call, CLI or MCP tool. It
should carry the target, variable clocks, corpus identities, proposed relation
types and permitted operations, and return evidence-backed alternatives and
identification status. Retrieved prose is evidence to inspect, not an instruction
that authorizes another tool call. A cycle or unsupported edge should produce a
specific unresolved result that downstream estimation preserves. These are
proposed contracts, not demonstrated cross-repository integrations. Public
prototype merit must precede an expensive experiment with private bank data.


% SOURCE docs/release/2026-10-29/source/docs/writing-repair/2026-10-29/src/db604379b4b12e69.tex CHAR 0
\section{Factual support in a causal assistant}
\label{exec22:fact-scope}

A retrieved explanation of a lending policy can be well supported by its cited
document and still leave the causal question unanswered. The document may
describe association, concern a different customer population, or rely on an
unverified assignment assumption. Wang and colleagues' survey of factuality
organizes evaluation and improvement around knowledge, retrieval and domain
specificity \cite{exec22_factuality_wang}. Its selected arXiv edition is useful
background for designing evidence assistance. Its taxonomy and surveyed
examples should not be read as independent validations of each underlying
method or as a theorem about causal diagrams.

The survey scopes factuality around determinable, undisputed facts. A causal
proposal frequently includes a different type of claim: whether a servicing
offer would change default, whether selection creates a spurious association,
or whether a time-indexed edge is needed. These questions can remain disputed
even when dates, document quotations and variable definitions are correct.
The assistant therefore needs an explicit distinction between what a source
says, whether that statement is warranted outside the source, and whether it
supports the proposed causal role for the population and time in question.

\subsection{Match the evaluation to the claim}

The survey distinguishes rule-based, neural, human and language-model
evaluation. Lexical overlap and semantic similarity can help locate likely
support, but neither is sufficient for entailment. ``The offer causes
repayment'' and ``the offer does not cause repayment'' share most words while
disagreeing on the decision-relevant proposition. Even a citation that
accurately supports the first sentence may describe a randomized intervention
in another setting or an author's unsupported interpretation of a regression.
Checking attribution and checking causal identification answer different
questions.

Atomic claim checking offers a useful design discipline: split the output into
claims whose support can be examined. For example, an assistant might assert
that treatment preceded the outcome, an instrument predicts treatment, its
exclusion restriction holds, and the estimated effect exceeds cost. The first
two may have direct documentation or statistical evidence. Exclusion usually
requires a structural argument; an effect-above-cost claim additionally needs
the target population, units, uncertainty and cost assumptions. One global
``factual'' score cannot make these evidence types interchangeable.

A proposed evidence record should therefore retain the exact source span,
edition, scope and claim role; distinguish quotation from project inference;
and preserve a conflict or unsupported status. Human reviewers can assess
scientific assumptions, but the survey also discusses subjectivity and
inconsistency in human evaluation. Model judges introduce their own bias and
domain limitations. Agreement among judges that share the same document,
training exposure or prompting pattern is not an independent scientific
replication. Calibrating a judge for citation entailment would still leave
causal identification and monetary decision quality to be assessed separately.

\subsection{Retrieval and structured knowledge help under stated conditions}

The survey identifies failure routes involving missing or conflicting
information, misleading evidence, distraction and misinterpretation. This
suggests concrete proposal requirements. Evidence search should preserve
contradictory passages and their populations, intervention definitions and
dates. Retrieval should allow ``insufficient evidence'' rather than making
general knowledge completion an obligatory final step. An assistant should
explain when a source settles a definition, offers a mechanism hypothesis or
reports an identified effect under assumptions. Increasing retrieved volume
does not by itself resolve disagreement among those roles.

Structured knowledge can improve entity and query mapping. For banking, this
could distinguish an application score from a post-decision monitoring score,
or connect two names for the same branch. A knowledge-graph relation such as
``associated with'', ``measured after'' or ``mentioned together'' is still not
a direct causal edge. An indirect mechanism should not silently become a
direct edge among observed variables. Temporal feedback requires explicit
time indexing or an appropriate cyclic model, rather than an acyclic drawing
that hides the feedback. The assistant's contribution is to expose candidate
interpretations and supporting evidence so that graph assumptions can be
reviewed.

The survey discusses debate and interactive retrieval as strategies described
in the literature. Those accounts motivate candidate workflows, not verified
superiority of a multi-agent implementation in this proposal. Several agents
repeating an unsupported premise may amplify confidence without adding
information. A useful design would require an independently sourced challenge,
record the unresolved disagreement and retain an abstention route. Its merit
must be measured against simpler retrieval and human-assisted baselines at
matched evidence and review budgets.

The survey's financial examples mainly concern language tasks. In that context
``causal language model'' can denote autoregressive next-token modeling.
That architectural term does not identify an intervention effect. Financial
text accuracy, sentiment prediction and financial question answering may help
prepare evidence, while a causal bank decision still requires an intervention,
target population, information clock, identification argument and uncertainty.
Domain fluency is useful only when those roles remain visible.

\subsection{An evaluation that can justify a prototype}

Public evaluation can test these distinctions before private bank experiments.
Freeze known-law causal cases and separate documentary corpora, including
negated statements, obsolete policies, contradictory populations, indirect
relations and observationally equivalent causal explanations. Compare simple
retrieval, structured retrieval and proposed agent workflows with matched
documents and budgets. Hold out documents and domains, and give adjudicators
the actual claim definitions rather than asking for an undifferentiated
plausibility score. These are proposed studies, not completed live-model
evaluations.

Report evidence retrieval and entailment, entity/time mapping, graph-role
errors, identification/refusal, target-weight preservation and decision loss
separately. Known-law cases establish ground truth by construction; public
observational data can test reproducibility and practical relevance but need
their own identifying assumptions. Neither should be labeled authentic bank
qualification. Include adversarial cases where every quoted sentence is
accurate but the causal conclusion or target population is wrong. Count corpus
curation, retrieval, repeated model calls, adjudication and maintenance when
comparing costs.

In the proposed CLI and MCP workflow, evidence status should accompany the
claim into the causal analysis request. Retrieved text supplies evidence, not
tool authorization. The result should distinguish supported premises,
unverified assumptions, identified or unidentified effects and the population
to which an interval refers. This would allow the assistant to explain why a
proposal needs more information or a narrower claim. The survey supports the
need to examine these evidence failures; actual reliability and economic merit
remain empirical questions for the bounded public prototype.


% SOURCE docs/release/2026-10-29/source/docs/writing-repair/2026-10-29/src/30f827c6014a5a93.tex CHAR 0
\section{An autonomous discovery workflow and the limits of its authority}
\label{exec17:copilot-development}

An analyst who can run a discovery algorithm still faces a substantial sequence
of choices: expressing the question, checking the observation clock, choosing
an admissible method, configuring it, interpreting its graph and communicating
what the graph permits. Causal-Copilot addresses much of this sequence in a
single agent workflow. Wang and colleagues describe natural-language input,
data diagnostics, algorithm selection and configuration, execution revision,
graph refinement and report generation. This is a useful architectural
comparison for the proposed banking assistant. Its value lies partly in
organizing work that a numerical discovery routine leaves to the analyst.
The extent of that automation and the causal authority of its output require
separate examination.~\citep{exec17_copilot_paper}

The evidence considered here has two distinct versions. The paper is
arXiv:2504.13263v2, dated 21 April 2025. The inspected repository snapshot is
commit \texttt{a31aa7c}, dated 15 May 2026. Both precede the literature cutoff,
but the later source has not been established as the code that produced the
paper's experiments. Paper results below are author-reported results;
source findings describe particular static paths in the later snapshot.
Neither supplies a reproduction of the other. This distinction permits a
useful component assessment without attributing a later implementation choice
to an earlier experiment.~\citep{exec17_copilot_paper,exec17_copilot_code}

\subsection{Translating a retention question into a causal request}

Suppose a manager asks, ``Which retention offers prevent withdrawals?''
Causal-Copilot's input and preprocessing modules offer a concrete starting
point: parse a natural-language objective and constraints, inspect the supplied
data, and allow the user to revise the request. Its subsequent modules filter
and rank algorithms, configure execution and revise errors. The paper also
describes an interactive graph revision loop and a generated PDF report;
its web interface is an author-described interface rather than a deployment
tested here. These capabilities show why the proposal needs a complete
analyst journey, including intelligible reasons for method selection, rather
than a collection of numerical commands.~\citep{exec17_copilot_paper}

For the retention request, a proposed causal contract must go further than
extracting the words ``offer'' and ``withdrawal.'' Let $X$ denote assignment of
a specified offer before a decision deadline, and let $Y$ indicate retention
over a fixed subsequent horizon. The population must specify eligibility at
that deadline. A receipt date, an acceptance date and an assignment date can
produce different treatments; an account balance observed after an offer can
be a consequence of treatment rather than an admissible adjustment variable.
Multiple offers require either a specified version or an explicit policy.
The assistant should expose these alternatives and ask for the information
that determines the intended contrast. A fluent question-to-graph translation
does not establish the timing or truth of the underlying variables.

The selected source makes part of the selection process concrete. A user-chosen
algorithm bypasses the filtering step. Otherwise, the filter sends the query,
data statistics, supplied knowledge, resource information and the requested
DAG or CPDAG form to a temperature-zero language-model selection step. Such a
prompt can organize alternatives and enforce an output preference. It cannot
make an observationally undetermined direction identifiable merely because a
DAG was requested. The proposal should retain the candidate class, rejected
alternatives and assumptions that justify the selection. User authority to
choose an algorithm is likewise distinct from evidence that its assumptions
hold.~\citep{exec17_copilot_code}

The proposed contract would therefore distinguish a documented business fact,
a modeling assumption, a statistical diagnostic and an agent suggestion.
For example, an audited assignment rule may justify an independence premise;
an apparent lack of linear association cannot certify it. A suggested latent
``relationship strength'' variable may explain selection, but naming it does
not measure it. In the eventual CLI or MCP request, these distinctions need
structured fields that survive the call to a numerical engine and its returned
report. They are proposed additions to this library, with semantic and runtime
qualification still required.

\subsection{Stable association can support incompatible decisions}

The distinction between workflow assistance and identification has a small
exact demonstration. Let $U$ and $V$ be independent, with
$U\sim\operatorname{Bernoulli}(1/2)$ and
$V\sim\operatorname{Bernoulli}(1/10)$. Write $\oplus$ for binary exclusive-or.
Consider two stipulated structural worlds:
\begin{equation}
 \begin{array}{lll}
 A:& X=U,&Y=X\oplus V,\\
 B:& X=U\oplus V,&Y=U.
 \end{array}
 \label{exec17:copilot-worlds}
\end{equation}
In world $A$, the offer changes retention. In world $B$, the hidden customer
state $U$ precedes both assignment and the later outcome, producing selection
without an offer effect. These are mathematical examples, not observed banking
mechanisms or fitted Causal-Copilot models. Both induce the same strictly
positive observed distribution:
\begin{equation}
 P(0,0)=P(1,1)=\frac9{20},\qquad
 P(0,1)=P(1,0)=\frac1{20}.
 \label{exec17:copilot-observed}
\end{equation}
Consequently, $E(Y\mid X=1)-E(Y\mid X=0)=4/5$ in either world.
An intervention replacing the equation for $X$ leaves $Y=U$ unchanged in
world $B$. In world $A$, it gives retention probabilities $9/10$ and $1/10$.
The interventional contrast is therefore $4/5$ in $A$ and zero in $B$.
Enumerating the four $(U,V)$ atoms reproduces both the observed law and these
effects. The pair demonstrates nonidentification in a class admitting both
worlds; it is neither a general sharp effect bound nor a statement about an
actual discovery run.~\citep{exec17_copilot_exact}

Now condition on six independent observations comprising three $(0,0)$ pairs
and three $(1,1)$ pairs. In an ordinary iid empirical bootstrap of size six,
the covariance of $X$ and $Y$ is positive whenever both types appear.
It vanishes only for either of the two constant-type resamples. Hence
\begin{equation}
 P_*\{\widehat{\operatorname{Cov}}(X,Y)>0\}
 =1-2(1/2)^6=\frac{31}{32}=0.96875.
 \label{exec17:copilot-stability}
\end{equation}
Here $P_*$ is conditional resampling probability, not probability that a causal
direction is true. Exhausting all $6^6=46{,}656$ index resamples gives 45,198
positive covariances. The conditioning event has probability
$\binom63(9/20)^6=531441/3200000$ under either observed law, so it is a possible
sample with an explicit probability. The stability calculation concerns a
covariance sign; it does not claim that any specific upstream algorithm returns
an edge with that frequency.~\citep{exec17_copilot_exact}

With a stipulated USD\,100 contribution per incremental retained customer and
a USD\,30 expense per assigned offer, the contrast against withholding is
\begin{equation}
 \Delta_A=100(4/5)-30=50,\qquad
 \Delta_B=100(0)-30=-30\quad\text{USD per offer}.
 \label{exec17:copilot-value}
\end{equation}
The same observed data can thus favor opposite actions in the two admitted
worlds. More confident repetition of the association does not select between
them. A useful assistant should identify what would change this conclusion:
a credible assignment experiment, defensible additional restrictions, or new
information that excludes a world. Which information is feasible remains a
banking design question. Bootstrap stability is still useful for describing
sampling sensitivity under a declared resampling law; its interpretation must
stop at that law.~\citep{exec17_copilot_exact}

\subsection{What the reported discovery comparisons establish}

The paper's synthetic experiments supply known graphs and vary the generating
conditions. They provide evidence about selection and discovery within those
conditions, including favorable settings and substantive alternatives. The
default tabular row in Table 2 reports Causal-Copilot F1 of 0.900, compared with
0.920 for PC, 0.910 for FCI and 0.920 for GES. In the sparse row, its 0.820 is
below GPT-4o's 0.830 and FCI's 0.840. These countercomparisons matter when
evaluating an adaptive selector: automation can be useful without dominating
each fixed method. They do not support a universal superiority claim.
\citep{exec17_copilot_paper}

The financial scenario is generated, with 50 nodes across three financial
domains; the reported Copilot F1 is $0.650\pm0.130$. Financial vocabulary does
not turn this known-truth simulation into an authentic bank validation. The
temporal Table 3 is explicitly linear and stationary. Its default row gives
0.673 for Copilot, 0.695 for PCMCI and 0.733 for DYNOTEARS. The table's
$n=5000$ row reports 0.668 for Copilot, whereas the adjacent prose says 0.649.
The very-large row labels $p=100$, whereas part of the prose says $p=1000$.
These two disagreements remain unresolved; no raw output or experiment-bound
configuration is supplied by this assessment to reconcile them.
\citep{exec17_copilot_paper}

The evaluation description on physical page 47 imposes a twenty-minute runtime
cap, describes a best-DAG conversion for tabular predictions and evaluates
temporal summary graphs while excluding diagonal entries. Such conventions
change the evaluated object. A summary adjacency can record that some lagged
relationship exists while losing which lag matters. Excluding a diagonal
also excludes a variable's own autoregressive relationships from that graph
score. Completion under a runtime cap is relevant operational evidence, but
an unavailable run and an incorrect completed graph are different outcomes.
A banking comparison should retain both, specify the scored graph object,
and examine the intended intervention query rather than infer effect quality
from a summary F1 alone.~\citep{exec17_copilot_paper}

The later source supplies a precise example of why the conversion convention
needs an explicit name. In \texttt{evaluator.py}, entries $[i,j]=1$ mean
$j\to i$, and the default is \texttt{sample=False}. The default metric method
calls a helper that consults \texttt{true\_graph} when orienting ambiguous
predicted edges. Its alternative sampling branch also selects using minimum
SHD against truth. The integrated tabular benchmark constructs the default
evaluator and reaches this helper; its temporal branch instead calls direct
binary scoring. This is truth-assisted benchmark evaluation in the inspected
paths. It is not evidence that a deployed discovery procedure receives the
unknown true graph, and it has not been bound to the 2025 experiments.
\citep{exec17_copilot_code}

For two nodes, hold fixed a single legal asymmetric undirected prediction and
consider opposite one-edge true DAGs. Both directions are compatible with this
unoriented prediction. With ordered off-diagonal binary F1, a truth-assisted
conversion can attain one in both worlds. A conversion fixed to one direction
before seeing truth attains one in one world and zero in the other. Under a
stipulated equal weighting of the worlds,
\begin{equation}
 \overline{F1}_{\mathrm{assisted}}=1,\qquad
 \overline{F1}_{\mathrm{fixed}}=\frac12.
 \label{exec17:copilot-score}
\end{equation}
This authored enumeration isolates the information used by the score; it is
not an estimate of the paper's score inflation. An optimistic compatibility
score can answer a legitimate question about whether an ambiguous output
contains the truth. It should be reported as that question, alongside ambiguity
and a declared prediction rule, rather than interpreted as a direction the
algorithm selected without truth.~\citep{exec17_copilot_exact}

\subsection{A graph-refinement invariant needs an executable guard}

Physical page 7 of the paper states that LLM recommendations are soft
suggestions and may not override high-confidence bootstrap decisions. That
is an intelligible division of authority. A source-level assessment must ask
whether the reachable editing path enforces it, including the representation
of opposite directions of the same pair. In the later snapshot,
\texttt{bootstrap\_recommend} identifies high directed frequencies at 0.9 or
above and low frequencies at 0.1 or below. Low ordered entries also enter its
``middle'' collection. The LLM evaluation combines that collection and
deduplicates it as unordered pairs. Its prompt distinguishes two directions,
answer C for no relationship and answer D for unknown.
\citep{exec17_copilot_paper,exec17_copilot_code}

Consider a conditional static witness using the source's receiver/source
matrix convention. Stipulate a raw $0\to1$ edge and twenty bootstrap results,
nineteen in that direction and one in the reverse direction:
\begin{equation}
 G=\begin{pmatrix}0&0\\1&0\end{pmatrix},\qquad
 B=\begin{pmatrix}0&0.05\\0.95&0\end{pmatrix}.
 \label{exec17:copilot-edit-input}
\end{equation}
Set uncertain-edge frequencies to zero. The forward ordered entry is high;
the reverse entry is low and therefore also enters the middle collection.
Unordered deduplication makes the pair eligible for LLM examination. In the
inspected \texttt{quality\_judge} path, \texttt{fixed\_pairs} is empty. If an
answer C is accepted for this pair, the resulting forbid entry reaches an edit
that sets both entries of $G$ to zero. A separately authored finite translation
of these expressions reproduces the mutation.
\citep{exec17_copilot_code,exec17_copilot_exact}

The conditions in this witness are essential. The bootstrap collection and
accepted answer are stipulated, not observed model outputs. The main program
reaches the judge and its quality path, but no upstream process or background
knowledge implementation was executed here. The witness stops immediately
after the LLM mutation; later KCI pruning and cycle handling are not simulated.
It therefore identifies an unguarded intermediate edit under stated inputs,
rather than claiming a reproduced final graph failure. It also does not
establish that the 2025 implementation violated the paper's rule.
\citep{exec17_copilot_code,exec17_copilot_exact}

A proposed repair would encode protection on unordered pairs before any edit:
\begin{equation}
 \begin{gathered}
 \mathcal P=\{\{i,j\}:i\ne j,\ B_{ij}\ge0.9\},\\
 \{i,j\}\in\mathcal P\ \Longrightarrow\
 G^{\mathrm{after\ LLM}}_{\{i,j\}}=G^{\mathrm{before\ LLM}}_{\{i,j\}}.
 \end{gathered}
 \label{exec17:copilot-guard}
\end{equation}
The equality refers to both ordered entries for the protected pair. Applying
this rule blocks the same conditional deletion in the authored reference.
A full integration would have to define the precedence of user changes,
statistical constraints and subsequent pruning, then verify the invariant at
each relevant boundary. Protecting an edge implements an authority rule; it
does not prove that the protected edge is causally correct. The two-world
example already shows why such a proof cannot follow from frequency alone.
\citep{exec17_copilot_exact}

\subsection{Temporal semantics, runtime scope and component adoption}

Temporal support requires following the path that is actually called. The
inspected bootstrap helper has an iid row branch and a temporal branch with
block size ten and block starts sampled without replacement. That is a
particular resampling scheme, not automatically a justified time-series
bootstrap. More directly, the inspected judge calls the helper with
\texttt{ts=False}, even though the helper offers a temporal branch. The
program wrapper can preserve returned lag matrices, and the state declares
raw, revised and lagged fields. Storage capacity does not establish that the
judge resamples, edits and reports the same temporal object.
\citep{exec17_copilot_code}

For a servicing query, an edge $X_{t-1}\to Y_t$ and an edge
$X_{t-5}\to Y_t$ imply different feasible intervention times. Summarizing
both as $X\to Y$ can be useful for an overview while being insufficient for
the requested policy. A proposal integration should preserve lag, endpoint
type, observation clock and the declared summary operation in its result.
It should select a resampling law justified for that target and explicitly
refuse a temporal confidence claim when only an iid route is available.
Whether a fixed block size and without-replacement scheme have suitable
behavior depends on the assumed dependence and asymptotics; the inspected
branch alone supplies neither a theorem nor a qualification study.

Runtime claims also need a boundary. The inspected benchmark timer encloses
\texttt{Programming.forward}, rather than the whole natural-language,
diagnostic, selection, refinement and reporting journey. Its source also
replaces missing values with zero, a preprocessing choice whose causal
consequences depend on the missingness mechanism. Default pytest settings
exclude integration, LLM and LaTeX markers; this is evidence of declared tests
and selected exclusions, not absence of tests. The requirements contain pinned
entries alongside an unpinned \texttt{huggingface-hub}. An environment lock,
end-to-end runtime measurement and semantic acceptance suite therefore remain
necessary before a reproducible runtime comparison.
\citep{exec17_copilot_code}

The bounded adoption decision is to use the workflow decomposition as an
architectural reference and assess components behind explicit contracts.
Question parsing, method explanations and report drafting are candidates for
adaptation with evidence-linked inputs and human revision. A statistical
wrapper is a candidate for reuse after validating its graph and temporal
semantics against an independent reference. The truth-assisted evaluator is
unsuitable as the sole operational accuracy score; it can be retained only as
a separately named compatibility diagnostic. The unguarded refinement path
requires a repair and invariant checks before adopting the paper's authority
promise. None of these static decisions qualifies an installed package.
\citep{exec17_copilot_code,exec17_copilot_exact}

The root license grants MIT permissions with a 2026 contributor notice; the
inspected vendored causal-learn notice grants MIT permissions with a 2022
notice. Both require preservation of the notice and disclaim warranty. These
two originals support a bounded licensing assessment of the inspected files.
They do not settle every transitive dependency, model provider, data asset,
deployment or redistribution right. Direct adoption remains conditional on
that complete inventory and on a tested environment; a port would create an
additional numerical and semantic equivalence obligation rather than removing
the rights question.~\citep{exec17_copilot_rights}

The comparison also exposes deficits in this proposal. It has not demonstrated
an integrated Copilot-like business conversation, algorithm selector and
interactive graph revision experience. Nor has it qualified the proposed
ambiguity-preserving engine boundary, temporal guard, evidence-linked report
or complete workflow costs. These are concrete prototype objectives. A first
public-data assessment can combine known-mechanism simulations with paired
requests that change assignment knowledge, lag information or the authority
to edit a graph. It should retain refusals, failures, ambiguity and total
workflow expense in the denominator, and use the exact examples above as
small reference cases. Such an assessment can justify whether an expensive
bank experiment is worth conducting. It cannot substitute public observational
labels for unknown causal truth or declare a banking product validated.
\label{exec17:copilot-development-end}


% SOURCE docs/release/2026-10-29/source/docs/writing-repair/2026-10-29/src/a4161d5b111a17a6.tex CHAR 0
\section{Choosing what to reuse from Causal-Copilot}
\label{sec:exec25-copilot-reuse}

A bank evaluating Causal-Copilot needs a component decision as well as the
scientific assessment in Section~\ref{exec17:copilot-development}. The useful
question is which responsibilities the proposed library should retain and
which it could delegate after a bounded trial. The inspected README links the
paper, repository and demonstration under one project identity and describes
V1 and V2 releases. Those names establish a documented project lineage; they
do not identify three independently validated engines. This assessment binds
the Lancelot39 repository at commit \texttt{a31aa7c}, dated 15 May 2026, and the
21 April 2025 paper separately. No claim that the later source produced the
earlier experiments follows from their shared name. Other similarly named
repositories require their own identity and revision checks.
\citep{exec17_copilot_paper,exec17_copilot_code}

The present choice is architectural reference use. The workflow decomposition
is a concrete comparator for business dialogue, method selection, execution
and reporting. Direct adoption would additionally require a reproducible
environment, a complete dependency and rights inventory, and demonstrated
semantic behavior. Selective adaptation of a parser or report template may be
less extensive, but still requires checking what information it loses. An
isolated integration could retain the upstream workflow while placing an
independent contract around it. A language port would assume responsibility
for every numerical, graph and failure conversion; it would not by itself
repair causal authority or reduce maintenance obligations. The proposed
library should reject automatic promotion of a suggested direction to an
identified intervention claim, while retaining useful components whose
contracts can be qualified. These are component choices, not a declaration
that the entire project is either safe or unsuitable.
\citep{exec17_copilot_code,exec25_copilot_choice}

\subsection{The boundary between assistance and scientific authority}

Consider again a retention offer assigned at a stated deadline. A proposed
request would carry a stable request identifier, a data and schema digest,
the eligible population, assignment and outcome clocks, treatment versions,
the requested contrast and a ledger of assumptions with their provenance.
The agent may suggest revisions, but a revision changes a versioned request
rather than silently changing the estimand. The selected upstream filter
bypasses algorithm filtering when the user already selected an algorithm;
that preference therefore needs a separate assumption check at the proposed
library boundary. A successful method call must not certify the assumptions
that motivated it.~\citep{exec17_copilot_code}

The corresponding result would distinguish an observational graph, a causal
claim conditional on assumptions, and an unavailable answer. It would preserve
variable identities, the convention for matrix entries, endpoint marks, lags
and raw versus revised graphs. The selected evaluator uses the convention
that entry $(i,j)$ refers to the arrow $j\to i$; silently adopting the
opposite convention would reverse a reported intervention direction. The
earlier source review also found that the inspected judge calls the bootstrap
helper with \texttt{ts=False}. A declared temporal request therefore needs
evidence that its actual path preserves the intended dependence and lag
semantics before receiving a temporal uncertainty claim. Neither a stored
lag array nor a successful JSON response supplies that evidence.
\citep{exec17_copilot_code}

For a minimal proposed representation, let $\mathcal G(r)$ denote the set of
directed possibilities represented by result $r$. With just two variables,
an unresolved adjacency can represent
\[
 \mathcal G(r_u)=\{X\to Y,Y\to X\},\qquad
 \mathcal G(r_d)=\{X\to Y\},\qquad
 \mathcal G(r_0)=\{\hbox{no edge}\}.
\]
The first two have the same skeleton but different directional information;
the first and third distinguish unknown direction from an absent relation.
This small authored representation is not a full CPDAG or PAG definition. It
demonstrates why a conversion test should compare represented possibilities,
not merely dimensions or adjacency counts. The complete adapter would need
the appropriate graph-class semantics, including latent-confounding marks
where relevant.~\citep{exec25_copilot_choice}

An edit proposal should carry the preceding graph digest, a reason, its
evidence, its author and the exact endpoints it changes. Scientific authority
and runtime permission are separate. A reviewer may permit an edit without
certifying the new direction. Conversely, a well-motivated edit must still
fail if it targets an outdated graph. The inspected refinement path and its
conditional witness motivate a symmetric protected-pair rule: both orderings
of a protected pair must be checked. An unknown answer should preserve
uncertainty rather than become a no-edge decision. The protected-pair rule
implements a declared authority policy; it does not prove that the protected
edge is true. Later statistical pruning needs its own explicit precedence
rule and audit, rather than an implicit exemption.
\citep{exec17_copilot_code,exec17_copilot_exact,exec25_copilot_choice}

\subsection{A public trial with distinct scientific and operational oracles}

The next useful experiment is a bounded prototype on generated mechanisms
and permitted public data. Before execution, the study should freeze task
families, graph conventions, provider and package versions, prompts, model
settings, seeds, splits and scoring rules. Generated structural models supply
known intervention effects for scientific checks. Public observational data
can test missingness, schema changes and workflow behavior, but unknown
causal truth must remain unknown. The earlier observational-equivalence
construction supplies a refusal test: without additional distinguishing
information, the two admitted worlds cannot justify one unconditional effect.
\citep{exec17_copilot_exact,exec25_copilot_choice}

Three comparisons answer different questions. A fixed method with a
prespecified configuration tests the cost and benefit of adaptive selection.
The source-pinned upstream workflow, once permitted and reproducible, tests
the documented agent journey. A proposed contract-wrapped workflow tests
whether preserved ambiguity, edit authority and typed refusals improve the
decision under the same tasks and information. A human-assisted baseline
can measure the review burden, with effort and expertise recorded. These
arms must receive matched data and information; allowing only one arm to use
the true graph would answer a different question. The inspected evaluator's
truth-assisted orientation may be retained as a named compatibility score,
but it cannot be the sole measure of an operational direction prediction.
\citep{exec17_copilot_code,exec25_copilot_choice}

The proposed qualification suite separates several failure modes. A variable
permutation fixture checks names and matrix conventions. Opposite admissible
orientations test ambiguity preservation. Paired lag-one and lag-five
requests test whether the result retains the information clock. Protected
edits in both identifier orders and a stale graph digest test edit authority.
Malformed answers, provider refusal, timeout, cancellation, duplicate delivery
and a crash between execution and report generation test the runtime
contract. Repeating a request after an uncertain timeout must not invent a
completion or double-count an accepted result. These are proposed tests; no
upstream package or provider was run to pass them in this assessment.
\citep{exec25_copilot_choice}

Scientific acceptance should be prespecified separately from those exact
interface invariants. A structurally valid graph can still be wrong, and
well-calibrated uncertainty in one stipulated regime need not transfer to
another. The trial should report graph ambiguity, target identification,
effect error and coverage where ground truth is available, together with
correct refusals, inappropriate refusals and failed runs. Concrete error
tolerances, power and sample size require the selected decision and precision
requirements; they are not supplied here by an arbitrary universal threshold.
This design can establish whether a larger experiment is worth its cost
without claiming authentic bank validation.~\citep{exec25_copilot_choice}

\subsection{Rights, ownership and the denominator of value}

The inspected root and vendored causal-learn notices are MIT notices. They
support reference and component assessment subject to their conditions, but
do not replace an inventory of all dependencies, model services, assets and
data. The selected requirements mix pinned packages with an unpinned
\texttt{huggingface-hub}; default pytest settings exclude integration, LLM
and LaTeX markers. Thus a named historical commit is not yet a complete
reproducible or maintained deployment profile. Present upstream support and
compatibility are unresolved rather than inferred from a historical README.
\citep{exec17_copilot_code,exec17_copilot_rights}

The scientific owner would sign the target, assumptions, graph class and
scoring rule. The adapter owner would maintain conversion, edit and failure
contracts. Platform and procurement owners would establish the environment,
service terms and retention policy. An evaluation owner would maintain
independent oracles and all-request accounting. A port additionally needs a
named maintainer for semantic parity across upstream changes. These are
proposed responsibilities; no institution has accepted them merely because
they appear in this proposal.~\citep{exec25_copilot_choice}

For an illustrative trial, suppose preparation costs USD\,1,200, all attempts
USD\,1,800, review USD\,2,400 and allocated maintenance USD\,600. Total cost
is USD\,6,000. With 120 requests, 100 completions and 80 scientifically
admissible results, the costs per request, completion and admissible result
are respectively USD\,50, USD\,60 and USD\,75. Completion is $5/6$, whereas
the admissible fraction is $2/3$. These exact authored quantities illustrate
why reporting only successful calls hides the cost of failures and refusals.
They are neither a vendor quote nor an estimated return on bank investment.
The useful denominator must follow the actual decision, with every attempt
retained and admissibility criteria fixed before outcomes are observed.
\citep{exec25_copilot_choice}

Reference use can therefore be decided now without declaring integration
ready. Reconsider adoption only after a selected artifact and complete
environment have passed the stated scientific and runtime contracts, rights
are established for the intended deployment, and matched full-cost evidence
supports the decision. Experiment-to-code binding remains necessary for any
reproduction claim. Generic agent repositories and diagram assistants still
require their own concrete candidate decisions.
\label{sec:exec25-copilot-reuse-end}



% SOURCE docs/release/2026-10-29/source/docs/writing-repair/2026-10-29/src/d237e95e325e91c3.tex CHAR 0
\section{Choosing agent components without changing the causal question}
\label{sec:exec25-agents-reuse}
Six concrete components clarify what a causal assistant can contribute to a
banking decision. LLM-CD, causal\_GPT and LLM-and-SCD supply different routes
from language-derived judgments to discovery constraints. CausalAgent exposes
selected analysis tools; oci-agent organizes notebook execution; OntoLearner
supplies an ontology-learning and loading component. These roles are useful
references for a proposed assistant, but they are not interchangeable complete
products. Their inspected revisions precede the 7 October 2026 literature cutoff;
the later acquisition dates and immutable source identities remain in the
source audit. The choices below concern those retained revisions and selected
paths, not current support or an installed service. Causal-Copilot retains its
separate lineage-specific decision.~\citep{exec14_iid_cd,exec14_iid_gpt,exec14_iid_scd,exec13_data_causalagent_code,exec13_data_oci_code,exec13_data_ontolearner_code}

Consider a bank asking whether an earlier servicing call reduces delinquency.
The initial sentence omits a population, a treatment version, an information
clock and an outcome horizon. A useful assistant should expose these omissions:
is the intervention an offer, a completed call or a scripted conversation;
which accounts are eligible; what was known before assignment; and when is
missed payment measured? The proposed output is a versioned question contract,
with alternatives and unresolved assumptions, before it becomes a graph or a
number. A domain dictionary can reconcile two names for an account status, but
that reconciliation alone does not prove temporal precedence or an adjustment
set. A learned graph can inform the next investigation while the intervention
query remains unresolved. This division of responsibility motivates the
component choices and the refusal conditions below.~\citep{exec25_agents_choice}

\subsection{Language-derived restrictions require a typed graph interface}
The inspected LLM-CD revision is
\texttt{8ea28948affdab88bee91}\allowbreak\texttt{a8b7858e89dc9424ffe}
(15 February 2025). Its CaMML wrapper serializes ancestral restrictions with
\texttt{=>} and absent direct-edge restrictions with \texttt{->}; the MINOBSx
wrapper has distinct blocks for absent edges, ordering and ancestry. These are
scientific distinctions, not merely alternate encodings. In the authored graph
$A\to B\to C$, $A$ is an ancestor of $C$ although the direct edge $A\to C$ is
absent. A statement that a servicing policy can eventually affect arrears must
not silently become a direct causal arrow after conversion. Conversely,
excluding a direct arrow leaves mediated paths possible. The proposed contract
therefore carries the relation type, ordered variable identifiers, source
statement, authority and confidence interpretation separately.~\citep{exec14_iid_cd,exec25_agents_choice}

The same retained wrappers change the working directory, use shared output
paths, invoke shell commands and then parse outputs without a checked success
branch. The optional shell timeout in the MINOBSx path does not itself certify
that the subsequent file is a fresh complete result. These are static findings,
not observed native failures. Conditional integration would require a bounded
adapter with per-request workspaces, explicit process status, output identity,
cleanup policy and cancellation semantics. It would also require resolved root
rights and backend entitlements: the nested BI-CaMML BSD-style notice does not
establish a license for the whole repository, and the inspected dependency list
includes \texttt{gurobipy==10.0.2}. The present choice is to retain the typed
constraint design as a reference and defer direct adoption and source porting
until those rights and the target environment are established. Independent
implementation of a published idea has its own source and qualification
obligations.~\citep{exec14_iid_cd,exec25_agents_choice}

At causal\_GPT revision
\texttt{2001fb5660d8e5d9a6c3}\allowbreak\texttt{a159b41d739316b8f91a}
(12 June 2024), the selected PC interface aggregates saved language-model
answers. Its statistical branch converts the answer-proportion decision into
numeric zero or one; the unfinished weighted-vote path and the missing-row
fallback also return zero. Such a number cannot be accepted as a calibrated
conditional-independence $p$-value under the bank's observation law. At a
threshold $\alpha=1/20$, a missing-result zero falls below the threshold just as
an apparently strong rejection would. The contract must instead return a
separate unavailable status and preserve whether evidence came from a
language-model judgment, a data test or domain knowledge. Repeated model
answers characterize a response-generating process; they do not by themselves
supply the sampling law of account outcomes.~\citep{exec14_iid_gpt,exec25_agents_choice}

The retained BSD-3-Clause notice makes causal\_GPT a candidate for a separately
reviewed adaptation, subject to its dependencies and provider terms. The
proposed adaptation separates answer collection from statistical testing,
retains raw answer provenance and missingness, and prevents a numeric adapter
from erasing that distinction. A historical OpenAI dependency pin and several
unlocked dependencies are declarations to reconcile, not a reproduced
environment. Reference use is selected now; adoption of the binary interface
as a bank test is rejected. A new wrapper could be assessed against generated
laws with known conditional independences, but this is a proposed experiment,
not a result of the source inspection.~\citep{exec14_iid_gpt,exec25_agents_choice}

LLM-and-SCD revision
\texttt{5aaeaea89e71e2cff443}\allowbreak\texttt{289e473c06d2f79172c5}
(31 January 2024) illustrates a third route. Code-only extraction of its
retained notebooks shows first-token top-five probability mass for the strings
\texttt{Yes} and \texttt{yes}, and strict thresholds below 0.05 or above 0.95
in selected constraint constructions. Backend constraints differ in meaning.
Neither token probability nor crossing one of these thresholds establishes a
probability that the causal relation is true. The proposal can retain this
elicitation design as a reference while preserving the exact prompt, model,
token treatment, thresholds and conversion meaning as experimental inputs.
The retained complete tree has no root license, and notebook dependencies do
not freeze the historical environment. Direct adoption and source porting are
therefore deferred; a conditional route first resolves rights, distinguishes
ancestry from direct edges, and measures scientific error and abstention on
held-out tasks. Saved notebook outputs are not treated as newly reproduced
evidence.~\citep{exec14_iid_scd,exec25_agents_choice}

\subsection{Tool transport, notebook targets and ontology identity}
CausalAgent revision
\texttt{87a4f6a45e9d77f9b0e8}\allowbreak\texttt{b7796489c47e50af41b8}
(28 January 2026) supplies a concrete reference for exposing causal tools through
MCP. In the selected module, the PC routine is imported whereas the OLC routine
is called without a definition or import in that module; the call is enclosed
in an exception handler. A separate output parser returns an empty dictionary
on malformed JSON. The review does not claim an observed service failure or
that every project path has this behavior. It identifies the exact paths that
must be resolved before the declared tool can be qualified. An empty valid
object, an invalid response and a scientifically unidentified request must have
different statuses. The root MIT notice and the selected MCP dependency pin do
not supply a running service or complete transitive rights. Reference use of
the tool organization is selected; conditional integration needs callable
binding, schema, graph and failure tests.~\citep{exec13_data_causalagent_code,exec25_agents_choice}

For oci-agent revision
\texttt{a978030e2c69cddab235}\allowbreak\texttt{e8737ee0ca527759ac49}
(27 July 2026), the selected runner executes a specification-selected notebook
with \texttt{NotebookClient(timeout=-1)}, writes the executed notebook, then
reads a result file. Its package metadata declares version 0.2.0 and a CLI
entry point. These facts motivate bounded execution, approved notebook
identity, resource accounting and a result digest in the proposed interface.
They do not establish an MCP service or qualified execution here. The retained
Apache-2.0 notice supports a concrete rights review for adaptation, while
notebook content, dependencies, data and provider use remain separately
scoped.~\citep{exec13_data_oci_code,exec25_agents_choice}

The target deserves equal attention. In the retained notebook, the selected
ATO path trims by propensity and takes a uniform AIPW summary of the retained
rows. Trimming and conventional overlap weighting generally define different
target populations. To see the distinction without running the notebook,
consider two equally prevalent strata, with supplied propensities $e_1=1/2$
and $e_2=1/10$, and supplied effects $\tau_1=0$ and $\tau_2=10$. Both strata
survive a stipulated interval $[0.05,0.95]$. Their uniformly averaged effect is
\begin{equation}
 \tau_{\mathrm{trim}}=(\tau_1+\tau_2)/2=5.
\end{equation}
For the explicitly defined overlap target with weights $e_s(1-e_s)$,
\begin{equation}
 \tau_{\mathrm{overlap}}=
 \frac{(1/4)\,0+(9/100)\,10}{1/4+9/100}=\frac{45}{17}.
\end{equation}
These are exact population calculations under stipulated stratum effects and
propensities, not finite-sample claims about an estimator. They show why a
label cannot authorize target substitution. Identification from banking data
would additionally require the appropriate consistency, exchangeability,
positivity and information-clock conditions. Conditional reuse must expose
the actual trimming rule, target weighting, nuisance fits and uncertainty
assumptions; a changed target requires a new question version. A business
owner seeking the effect among borderline accounts should not receive a
uniform retained-population average under an overlap label.~\citep{exec13_data_oci_code,exec25_agents_choice}

OntoLearner revision
\texttt{074a41ba23a85f5b72d4}\allowbreak\texttt{e28da042eab3d50bd90a}
(3 September 2026) is a reference for semantic resources rather than a complete
causal question-and-diagram assistant. The inspected loader requests a
Hugging Face dataset without a revision argument. Fixing the code commit
therefore does not fix the ontology bytes. The root MIT notice also does not
establish rights for each downloaded ontology. The proposed interface records
dataset repository, immutable revision, file hash, rights, concept identifiers
and mapping provenance. It keeps an ontology relation such as membership or
subclass distinct from causal influence. A mapping from two database fields
to the same concept is a reviewable semantic hypothesis until the banking
meaning and time of observation are established. Conditional component reuse
is plausible after these controls and task-specific evaluation; generic
approval of causal semantics is not.~\citep{exec13_data_ontolearner_code,exec25_agents_choice}

\subsection{A common qualification contract with component-specific decisions}
The two umbrella leads are resolved into the six named component decisions,
not adopted as families. All six are retained as bounded references. LLM-CD
and LLM-and-SCD require unresolved root rights before source copying; the
other four permit a more concrete adaptation assessment under their retained
notices, without assuming that dependencies, data or providers are covered.
For each candidate, integration remains conditional on its scientific and
operational contract. A port is a separate implementation with the same
semantic tests, not an escape from source rights or statistical assumptions.
The proposal rejects treating language judgments as calibrated data tests,
transport declarations as running services, ontology relations as causal
edges, or one estimand label as proof of another target. It does not reject
all uses of the underlying projects.~\citep{exec25_agents_choice}

The proposed common request contains a logical request identifier, a question
version, data and schema digests, population, treatment and outcome clocks,
graph class, typed restrictions, model and prompt identities, environment,
rights profile and a budget deadline. A response separates scientific status
from execution status: a computation may finish with a conditional answer,
an identified set or a refusal; a timeout may leave execution completion
unknown. Raw and edited graphs retain variable order, endpoint marks and edit
authority. Notebook outputs expose target weights; ontology outputs expose
source identities. A CLI or MCP envelope must preserve these fields through
replay, rather than letting a caller infer meaning from a successful process
exit. These are proposed interface obligations, not an implemented common
service.~\citep{exec25_agents_choice}

A public qualification study should compare a fixed conventional workflow,
a permitted frozen upstream component, a proposed wrapped version and a
human-assisted comparator with recorded effort. Generated structural models
supply known graph or effect truth; permitted public observational data can
stress schema, missingness and workflow without inventing its assignment
truth. The study must freeze versions, tasks and held-out splits, prompts,
provider settings, environments, scientific tolerances and scoring before
measurement. An ancestry-only relation, an absent cache row, threshold-boundary
answers, malformed JSON, stale results, missing callable bindings, deadline
expiry, changed notebook targets and a changed ontology revision each need a
specific oracle. Rephrased questions must preserve their estimand or explicitly
open a new version. The same protocol records successful refusals and
unknown-completion states rather than dropping them from the denominator.
No such native comparison is claimed here.~\citep{exec25_agents_choice}

The ownership and economic case follows the same boundary. Scientific owners
approve the population, graph restrictions, target and uncertainty; adapter
owners maintain conversion and failure semantics; platform and rights owners
freeze permitted artifacts; evaluators control held-out oracles; and maintainers
review model, prompt and dataset drift. These are proposed responsibilities,
not signed institutional assignments. For a purely hypothetical batch, USD 900
of preparation, USD 1,100 of all attempts, USD 2,500 of review and USD 500 of
maintenance allocation sum to USD 5,000. With 100 attempted requests, 80
completed requests and 50 scientifically admissible results, costs are USD 50,
USD 62.50 and USD 100 respectively. A tool-call timer cannot decide which
component improves the banking workflow. The next investment decision needs
matched information, accepted scientific outputs and complete measured costs,
with explicit limits on rights, targets and native qualification.~\citep{exec25_agents_choice}


% SOURCE docs/release/2026-10-29/source/docs/writing-repair/2026-10-29/src/9aa6322883ab0b20.tex CHAR 0
\section{Temporal discovery with hidden causes: what an agent can add}
\label{exec15:latent-temporal}

\subsection{Start with the intervention, the clock and the hidden alternative}

A servicing team observes that a borrower receiving more contact this month
has a different repayment outcome next month. It wants to change contact
intensity, not merely forecast the outcome. Let $X_{it}$ denote the intensity
for borrower $i$ at decision time $t$, and let $Y_{i,t+1}$ denote the subsequent
net recovery before contact cost. A useful first target is
\[
 \tau=E\{Y_{i,t+1}(\operatorname{do}(X_{it}=1))
             -Y_{i,t+1}(\operatorname{do}(X_{it}=0))\}.
\]
Here the expectation is over a declared eligible borrower population and its
predecision histories. A multistage policy requires a separate target that
also specifies subsequent actions. Recording a message after a payment does
not make the message a prepayment intervention. Event time, availability time,
decision time and follow-up must therefore be represented separately. An agent
can elicit these distinctions and propose candidate mechanisms. Neither the
ordering of the records nor a persuasive mechanism description establishes
exchangeability, the absence of selection, or a manipulable treatment.
\citep{v12_la_2606_23608,exec15_latent_design}

The distinction can be seen without fitting a model. Consider a centered
two-wave Gaussian example, with independent standard normal $U,e_x,e_y$:
\[
 X=U+e_x,\qquad Y=2U+e_y.
\]
The hidden state $U$ can represent persistent financial distress. There is no
effect of intervening on $X$. A second world has $X\sim N(0,2)$ and
$Y=X+e$, where $e\sim N(0,3)$ is independent of $X$. Its intervention effect
is one. Both worlds have exactly the same observed bivariate Gaussian law,
with covariance matrix
\[
 \Sigma=\begin{pmatrix}2&2\\2&5\end{pmatrix}.
\]
The regression slope is one in both worlds; the time order is also the same.
More observations of this law cannot choose the causal effect. If one outcome
unit is stipulated to be USD~20 and the added contact costs USD~8, the two net
effects are minus USD~8 and plus USD~12. A language model that labels $U$
``distress'' has supplied an interpretation, not an observation of $U$ or a
reason to discard either world. These exact constructions explain why a
proposal must preserve ambiguity until a defensible restriction, proxy,
assignment mechanism or experiment supplies additional information.
\citep{exec15_latent_exact}

Time indexing also changes the graph class. The arrows
$A_t\to B_{t+1}$ and $B_t\to A_{t+1}$ are compatible with an acyclic
time-unrolled model. Collapsing time produces the cycle $A\to B\to A$.
Likewise, $A_t\to A_{t+1}$ is ordinary autoregressive persistence, not a
contemporaneous self-cycle. An instruction to ``repair the DAG'' by deleting
these arrows would remove valid candidate mechanisms. The product needs a
typed graph carrying lag units, contemporaneous edges, endpoint marks and
latent-variable semantics before it asks an agent to explain or modify it.
\citep{exec15_latent_exact,exec15_latent_design}

\subsection{Five frontier roles that should be evaluated separately}

RealTCD uses language-derived meta-initialization within temporal discovery
from observational and interventional regimes. Its unknown-target result is
valuable because intervention targets need not all be supplied in advance.
It nevertheless assumes causal sufficiency. Theorem~3.1 concerns a global
score maximizer under sufficient conditional-density capacity, intervention
faithfulness, strict positivity, finite entropy and suitably small penalties.
These are substantive restrictions on the data and optimization problem.
LLM initialization does not prove that a numerical optimizer reaches that
maximizer. A month with unusually many arrears is not automatically an
intervention regime: an unrecorded macroeconomic shock can change both
contact and repayment. For banking, the proposed use is to suggest initial
lagged structures and possible regime targets while retaining a separate
register of which regime assignments are observed, assumed or conjectured.
Its reported benchmark advantages are author results under their designs,
not evidence that hidden banking causes have been recovered.
\citep{v12_la_2404_14786}

Sequential Causal Discovery with Noisy Language Model Priors addresses a
different problem. Its sequential axis consists of arriving observation
batches; it should not be silently equated with a fully specified temporal
structural model. The inspected formulation allows latent confounding under
a restriction that each confounder affects two observed variables. Language
judgments update edge-level beliefs in a Bayesian-inspired scheme. The
authors explicitly distinguish this from full probabilistic inference over
graphs, and leave broader confounding, calibration and interventional
extensions as future work. Their discussion also warns against transferring
standard independent, stationary bandit regret guarantees to the dependent
language-response channel. This is a useful approach to noisy suggestions,
but an edge belief cannot be presented as a calibrated posterior over all
banking mechanisms or as confidence in a policy value.
\citep{v12_la_2506_16234}

Large Causal Models provide an essential numerical comparator. They are
pretrained numerical transformers for temporal discovery, not large language
models interpreting business text. The inspected paper assumes additive
noise, causal stationarity, no hidden confounders and no contemporaneous
effects. Its limitations section says that the model does not detect or
correct violations of these assumptions; predictions can instead reflect
training-distribution biases. Reported speed and accuracy under its test
design motivate a challenger in that declared domain. They do not justify
using the model as a hidden-confounder detector. A bank comparison should
therefore separate the gain from numerical pretraining from any gain due to
language descriptions, expert restrictions or additional information.
\citep{v12_la_2602_18662}

Model Discovery Agent assigns language another role: proposing additions to
a mechanistic hypothesis library. Candidate models are fitted and compared
using evidence approximations, including BIC or tempered sequential Monte
Carlo in different settings. The inspected experiments use open-loop inputs;
the method discussion distinguishes these from closed-loop actions that
depend on observed history. The finite-pool interpretation is particularly
relevant to banking. Even exact evidence would give a posterior restricted
to the enumerated candidates. In the implemented search, evidence is
approximate, top-pool pruning is deterministic hypothesis search, and
unproposed model mass is not represented. Small residuals can delay further
expansion, and a discarded model can later become relevant. An agent can
thus make a model library more useful without establishing that it contains
the true policy-response mechanism. A servicing application needs an archive
of rejected alternatives and challenges that distinguish their intervention
predictions, not only their fit to logged trajectories.
\citep{v12_la_2608_09696}

The protected-core architecture advocated in Causal Discovery in the Era of
Agents supplies a way to combine these roles. The formal discovery method
continues to determine the meaning of its output; the agent proposes context,
retrieval, assumptions and analyses around it. A CPDAG represents a class of
DAGs, and a PAG preserves ambiguity arising with latent or selection
variables. Neither should become a single confident DAG because the report
needs a diagram. This is an architectural argument and prospective
evaluation standard. It is not an empirical demonstration of safe banking
performance, and its causal-learn+ identity is distinct from the CausalAgent
repository inspected below.
\citep{v12_la_2606_23608}

Evidence against unqualified language assistance belongs alongside these
favorable possibilities. The inspected study on integrating language priors
reports degradation when priors are inaccurate, particularly for the
temporal Water and Mildew examples. A model's familiarity with a variable
name can encode the wrong temporal scale, population or intervention.
Consequently, a fair comparison must include deliberately reversed or
misleading priors, opaque variable names, and a no-language baseline. Tuning
prior strength on the final evaluation mechanisms would conceal the very
failure the comparison is meant to measure.
\citep{v12_la_2306_16902,exec15_latent_design}

\subsection{Additional non-Gaussian information: a useful but restricted route}

Shi and colleagues develop a two-wave method that can distinguish some
causal alternatives using higher-order cumulants. Their model is
\[
 X=\beta_1U+e_x,\qquad Y=\beta_2U+\eta X+e_y,
\]
with temporal ordering, linearity, independent errors and latent component,
and specified non-Gaussian information. Define the fourth mixed cumulants
$C_{31}=\operatorname{cum}(X,X,X,Y)$,
$C_{22}=\operatorname{cum}(X,X,Y,Y)$ and
$C_{13}=\operatorname{cum}(X,Y,Y,Y)$. The first statistic is
\[
 \Delta_1=C_{31}C_{13}-C_{22}^2.
\]
For this scalar-component model, writing
$\kappa_U=\operatorname{cum}_4(U)$ and
$\kappa_x=\operatorname{cum}_4(e_x)$ gives
\[
 \Delta_1=\kappa_U\kappa_x\beta_1\beta_2^2
                 \eta(\beta_2+\eta\beta_1).
\]
The proof's nonzero cumulants and loadings, and the excluded cancellation
$\eta=-\beta_2/\beta_1$, matter. Under those conditions the statistic supplies
information absent from the Gaussian example. It does not create information
in that example by changing the estimator.
\citep{exec15_latent_shi,exec15_latent_exact}

Two exact checks make the boundary concrete. First, let $U$ take values
$-1,0,1$ with probabilities $1/6,2/3,1/6$. It is nonnormal, has variance
$1/3$, and has fourth cumulant $1/3-3(1/3)^2=0$. With independent
Rademacher errors, $X=U+e_x$ and $Y=U+X+e_y$ have a direct effect of one.
Nevertheless, $C_{31}=C_{22}=C_{13}=-2$, so $\Delta_1=0$; the relevant
third cumulants are also zero. This does not contradict the result with its
full nonzero-cumulant premises. It shows why ``nonnormal'' alone is an
insufficient eligibility test and why nondetection must remain inconclusive.
\citep{exec15_latent_exact}

Second, consider independent Rademacher $U,V,e_x,e_y$ with
\[
 X=U+V+e_x,\qquad Y=U+2V+e_y.
\]
There is no direct $X\to Y$ path. Direct enumeration of all sixteen equally
likely atoms, independently checked by adding factor cumulants, gives
$C_{31}=-6$, $C_{22}=-10$ and $C_{13}=-18$. Thus $\Delta_1=8$.
After standardizing $X$ and $Y$, whose variances are three and six,
the determinant is $8/(3^2 6^2)=2/81$, still nonzero. The two hidden
loading vectors $(1,1)$ and $(1,2)$ are not proportional. Although the
paper discusses aggregating hidden factors into a global component, one
common scalar component is not automatically an adequate representation of
arbitrary loadings with the required independent errors. This example lies
outside that scalar restriction; it supplies a diagnostic requirement for a
proposed banking application, not a counterexample under all theorem
premises. Multiple forms of persistent borrower distress should not be
collapsed into a single latent score merely because that makes the test
available.
\citep{exec15_latent_shi,exec15_latent_exact}

The proposed multiple-testing algorithm standardizes the observations and
uses percentile bootstrap intervals for three statistics with a Bonferroni
allocation of $\alpha/3$. The paper explicitly treats failure to detect as
inconclusive. In one specified weak-signal simulation, reported power rises
from 4.8\% at $N=2{,}000$ to 99.5\% at $N=200{,}000$. These are useful
warnings about information requirements, not a general sample-size formula
for a bank. Calendar dependence, selection, heavy tails and estimated
high-order moments require their own inferential argument. A bootstrap
implementation cannot supply that argument by itself. Moreover, detecting
the presence of a path does not consistently estimate its magnitude. The
contact decision still needs an effect or effect set in recovery units.
The appropriate reuse is a narrowly qualified statistical challenger and
assumption diagnostic, with stronger restrictions declared explicitly.
\citep{exec15_latent_shi}

\subsection{Screening, hierarchy and prediction are complementary evidence}

Yu and colleagues illustrate a more modest and practically useful language
role in dementia research: screening thousands of candidate variables using
word and document embeddings, followed by manual review and statistical
discovery. This is not evidence that a modern agent independently discovers
a true temporal mechanism. Their pipeline uses a mixed graphical model
skeleton followed by FCI-MAX, temporal restrictions, fifty imputed datasets,
and edge-frequency thresholds of 30\%, 50\% or all datasets. Manual review
restores risk factors missed by the information-retrieval models. The lesson
for a bank is to test the recall of the variable-screening stage before
celebrating a smaller graph: a missed common cause can undermine the
downstream analysis. Agreement across imputations is a stability statistic,
not automatically a calibrated graph posterior or a causal confidence set.
The source's cohort prose reports 406,837 participants whereas its table
reports 406,781; that discrepancy is unresolved, and no exact cohort-size
claim is used to justify a banking comparison. UK Biobank access requires a
researcher application, and the supplementary code has not been inspected
here. Publicly described data should therefore be distinguished from data
available for unrestricted immediate reuse.
\citep{exec15_latent_yu}

Spatio-Temporal Hierarchical Causal Models offer another adjacent route.
The hierarchy distinguishes units, such as regions, from repeated subunits,
such as sensors. Its collapse arguments take the number of subunits toward
infinity while the numbers of units and time points remain fixed. The
stated conditions include convergence of mechanisms and empirical measures,
static unit confounders, neighbor restrictions, a known global ordering of
contemporaneous spatial influence, and common functional mechanisms. The
adjustability result additionally requires observable history to block the
relevant backdoor paths. More observations within a branch cannot establish
that premise or make a time-varying unobserved shock static. A proposed bank
use must specify which objects are units and subunits and why the asymptotic
sequence represents the intended decision.
\citep{exec15_latent_sthcm}

The paper also states instrument-based identification claims that must be
read with a target-specific argument. A valid instrument alone does not
identify an arbitrary population intervention distribution. To see the
problem, let an independent binary instrument $Z$ encourage treatment among
a population split equally between compliers and never-takers. Compliers
have $Y(0)=0,Y(1)=1$; never-takers have $Y(0)=0$ and either $Y(1)=0$ or
$Y(1)=1$. Both worlds satisfy exclusion, relevance and monotonicity and
produce the same observed $(Z,A,Y)$ law. The complier effect is one, but
the population average effect is respectively one half or one. Additional
restrictions are necessary before translating an instrument into the
population recovery target. This construction does not purport to refute
every specialized hierarchical result; it prevents a generic IV sentence
from being treated as a universal identification theorem.
\citep{exec15_latent_sthcm,exec15_latent_exact}

The hierarchical paper reports favorable simulation comparisons and traffic
applications, including stress settings with drift or feedback. Such results
are useful leads for a prototype. A finite simulation mean of 4.91 near a
truth of five is not itself a proof of unbiasedness, and observed robustness
does not remove the theorem's assumptions. Its public code lead is a
separate item for version, rights and numerical review. It is also a
statistical modeling contribution rather than a demonstrated language-agent
capability.
\citep{exec15_latent_sthcm}

Augur uses language-generated causal explanations and teacher--student
distillation for discrete trend prediction. Its graph-overlap consensus and
summary-to-graph consistency scores help select training explanations.
Agreement among sampled graphs, however, can reproduce one shared
misconception. A narrative can match every edge of an incorrect graph.
The inspected forecasting and explanation results therefore motivate a
prediction or explanation baseline, not a claim of hidden-confounder
recovery or identified intervention value.
\citep{exec13_census_augur_arxiv}

CausalPulse illustrates the difference between a successful workflow and
successful causal diagnosis. In its manufacturing evaluation, reported
aggregate operational success reaches 98.73\% for PSE, while root-cause
Hits@1 is 0.44 and F1 is 0.39. Its reference labels are manually identified
``immediate variables changed'' at an anomaly timestamp. Those labels and
process-order rules may be useful for its diagnostic task, but changes can
be downstream consequences, and neither operational success nor those
annotations establish a bank's intervention truth. The appropriate banking
comparison reports workflow completion, diagnostic ranking, identification
and decision performance separately. No one metric can stand in for the
others.
\citep{exec13_pulse}

\subsection{Inspect the executable boundary before proposing adoption}

The pinned LCM source declares package version 0.1.0, Python at least 3.10,
an Apache-2.0 root license and six direct dependencies, including NumPy below
two and PyTorch Lightning 2.0.8. These declarations do not settle all
dependency or checkpoint rights. The advertised checkpoints support up to
twelve variables, lag three and length five hundred; their external bytes
have not been qualified here. The inspected forward path returns a tensor
with axes batch, target variable, source variable and lag index. The example
applies a sigmoid, and its convention maps index $l$ to lag
$\ell_{\max}-l$. For example, index two at maximum lag three means lag one.
A wrapper that reverses source and target or treats the last axis as elapsed
calendar time silently changes the causal claim.
\citep{exec15_latent_lcm}

The README's example contains a concrete masking issue. It loops over the
three lag positions but indexes the two variable dimensions:
\begin{verbatim}
for l in range(pred.shape[-1]):
    pred[:, l, l] = 0
\end{verbatim}
For shape $(1,12,12,3)$ this clears nine cells, all lag entries on variable
diagonals zero through two; diagonals three through eleven remain. More
fundamentally, valid lagged self-edges should survive unless the declared
model excludes them. The remedy is therefore to decide which temporal
edges are admissible before correcting the index loop. The example's global
min--max normalization also has a zero denominator for a constant input,
and its random padding needs explicit seeds, validity semantics and
evaluation controls. The inspected loading helper returns no model on a
missing or failed checkpoint and uses a loader path with unrestricted
weight deserialization. Checkpoint identity, loading behavior and failure
classification belong in a prototype contract. These are static-source and
exact-coordinate findings, not measurements from running LCM.
\citep{exec15_latent_lcm,exec15_latent_exact}

For DMIRLAB's CausalAgent, the inspected route spans tool selection, MCP
dispatch, the statistical backend, graph extraction, cycle processing and
report preparation. Its root license is MIT; requirements include specified
MCP and LangGraph versions as well as lower-bounded statistical dependencies,
so a resolved compatible environment remains necessary. The MCP module
imports the PC function but invokes an OLC function without importing it,
even though that function exists in the backend. A missing binding and an
absent optional CDMIR dependency are different failures. Neither should be
described as a scientific conclusion that hidden causes are absent. The PC
path uses Fisher's $Z$ at level 0.05; a generic tool description cannot
extend that test's validity to arbitrary data-generating laws.
\citep{exec15_latent_causalagent}

Graph conversion introduces another distinction. The PC path retains raw
endpoint matrices, while the OLC formatter documents its own row/column and
coefficient conventions. The downstream cycle detector keeps entries equal
to one and transposes them into a directed graph. If a bidirected endpoint
pair is admitted, this projection produces two directed arrows and a false
directed two-cycle. This is a conditional converter example, not an assertion
that ordinary PC outputs bidirected graphs. Undirected and circle endpoints
also need explicit treatment; deleting them is information loss, not
resolution of uncertainty. The detector catches exceptions and returns the
same pair as ``no cycle.'' A validation error must instead be an explicit
unavailable state.
\citep{exec15_latent_causalagent,exec15_latent_exact}

Tracing the caller prevents an incomplete repair recommendation. It does
run the cycle detector again after proposed edits, but only warns if cycles
remain and still reports the original cycle count as fixed. The editing
function deletes one matrix entry without a paired-endpoint contract. In
the normal dictionary-result path, a reversed type guard leaves the critical
edge list empty; the emitted revised graph is then an empty list, and the
edited working matrix is not used for that output. The alternate string
branch attempts dictionary access on a string. Independently, the edge
evaluator's prompt placeholders differ from the keys supplied at invocation
and its exception path returns a low-confidence retain default. Correcting
only that prompt would not repair the dictionary branch or the lost edited
graph. A bounded repair must address the entire typed result path and verify
the resulting report, while keeping original and proposed graphs separately
available. No live application or language model was invoked for these
observations.
\citep{exec15_latent_causalagent}

These findings support reference use before integration. Direct adoption is
deferred because target semantics, native behavior, dependencies and weight
identities remain unqualified. Porting entire neural or agent stacks would
add maintenance and equivalence obligations without an established benefit.
A narrow adapter is a plausible next step: preserve upstream numerical
methods while owning the clock, graph, failure and evidence contracts.
Rejecting all research use would discard useful challengers, but accepting
an MCP response or an acyclic projection as causal certification would be
equally unjustified. The proposed adapter must round-trip permitted endpoint
pairs and lagged edges, distinguish validation errors from negative findings,
and retain every edit's origin and justification.
\citep{exec15_latent_design}

\subsection{A public and generated comparison that can resolve useful questions}

The first investment should be a bounded prototype using generated truth
and suitable public observations. Public time series can supply realistic
clocks, missingness, vocabulary and forecasting challenges. They supply
causal ground truth only when assignment or interventions are independently
documented. Generated structural models supply known intervention responses
and controlled violations. Combining those roles is more informative than
declaring a public observational graph to be the true bank mechanism.
It also avoids making expensive bank experiments a prerequisite for
investigating basic conversion errors.
\citep{exec15_latent_design}

The comparison should have distinct arms: a numerical method without
language; the same method with expert-supplied restrictions; language
proposals accepted only through the formal core; deliberately corrupted
priors; and the numerical LCM challenger within its declared sufficient,
stationary, lag-only domain. Mechanism families should vary scalar versus
multiple latent loadings, weak or zero higher-order signal, temporal
feedback and contemporaneous effects, irregular or aggregated clocks, and
drift and selection. Identical observations and information budgets should
be used wherever the scientific target permits. An expert arm given true
hidden structure is an information oracle, not a fair practical competitor.
\citep{exec15_latent_design}

Evaluation units must be independent mechanisms or complete trajectories
under a declared sampling design, not repeated prompts or neighboring rows.
Development, calibration and final evaluation need separate mechanisms.
Freeze source and weight identities, variable descriptions, prompts,
prior-strength rules and stopping budgets before the final set. A pilot can
estimate runtime and variance to choose a defensible replication budget;
no final qualification dataset or success rate is claimed here. Report
endpoint accuracy and retained-class validity separately from identification
correctness, effect-set coverage, width, availability, decision loss and
total review/computation cost. Unavailable tools, failed fits, timeouts and
refusals remain in their declared denominators. A repaired graph conversion
can pass exact native fixtures while the scientific model remains
inapplicable to a particular hidden-confounder regime.
\citep{exec15_latent_design}

Several remedies are already concrete. Exact endpoint and lag fixtures can
detect conversion loss. Constant-input and padding fixtures can enforce
preprocessing eligibility. A missing OLC binding or backend can produce a
typed failure. The dictionary, prompt, structured-response and output
branches can be corrected together and tested with generated results.
Public metadata-screening exercises can measure omitted-variable recall
against an independently curated reference. None of these steps requires
pretending that a language-generated latent label is an observed variable.
They should proceed while the broader statistical questions remain open.
\citep{exec15_latent_design}

\subsection{The unresolved bridge from graph uncertainty to monetary uncertainty}

Let $M_\star$ denote the true model in a stated model universe and let
$\widehat{\mathcal M}$ be a data-dependent retained set. Suppose a justified
procedure provides
\[
 P(M_\star\notin\widehat{\mathcal M})\leq\alpha_G.
\]
For each retained model, let $C_M$ contain the target effects compatible with
its assumptions and an appropriate observed-law confidence set. Suppose
the construction also justifies a bound $\alpha_E$ for failure of the
true-model effect set under the actual selection and sampling procedure.
Then the union
\[
 C=\bigcup_{M\in\widehat{\mathcal M}}C_M
\]
misses the true effect with probability at most $\alpha_G+\alpha_E$, by the
union bound. Independence between these two failures is unnecessary once
both event bounds are valid. Obtaining those bounds is the difficult part;
using the same data for selection and effect estimation cannot be ignored.
A model universe omitting the true mechanism also invalidates the premise.
\citep{exec15_latent_exact,exec15_latent_design}

For illustration, budgets of 0.03 and 0.02 allow a 0.05 overall failure bound.
An exact hundred-atom construction with three graph failures and two disjoint
effect failures attains this bound. No graph-frequency histogram, low
entropy over LLM outputs, or successful DAG check establishes either
premise. Nor does a narrow interval conditional on one selected graph
account for excluded alternatives. In the Gaussian banking example, even
infinite observational data leave both effects zero and twenty dollars
available unless extra assumptions are justified. Conservative uncertainty
can therefore remain wide for an informational reason, not because a
numerical solver needs improvement.
\citep{exec15_latent_exact}

The known research route is to declare a restricted temporal model family,
derive or verify its retention and observed-law confidence procedure,
propagate all compatible effects, and account for numerical approximation.
Exact small models should precede dependent-data studies and larger
learners. General nonlinear latent systems with selection, changing
mechanisms and potentially false language priors are not covered by that
route merely because a restricted example succeeds. Useful finite-sample
precision and robustness in those broader regimes remain open questions.
When equivalent models imply opposite decisions, a valid outcome is an
effect set or a request for additional information, rather than a forced
point estimate.
\citep{exec15_latent_design}

An illustrative value-of-information calculation helps decide whether that
additional information merits a later experiment. Suppose, explicitly as a
decision-maker's stipulated prior, the direct-effect world has probability
$3/5$. Acting immediately has expected net value
$(3/5)12+(2/5)(-8)=4$ dollars per eligible case. Perfectly distinguishing the
worlds before acting would yield $(3/5)12=7.2$ dollars, so perfect information
is worth at most USD~3.2 before its cost. This is an upper bound under the
stipulated prior and two-world decision problem, not an observed bank price
or the value of an imperfect trial. A robust decision-maker who declines to
assign that prior may instead retain the net-effect set $\{-8,12\}$ and
withhold action under a worst-case criterion. The proposal can explain both
decision rules without inventing institutional preferences.
\citep{exec15_latent_exact}

The resulting program is specific about what can be done now: improve
semantic contracts, repair selected adapters in a prototype, test exact
counterexamples, and run predeclared public/generated comparisons. It is
equally specific about what is unresolved: broad graph/effect-set
calibration, model misspecification, useful precision and the information
needed to distinguish observationally equivalent mechanisms. This review
recovers three additional primary originals, but the diffusion-model lead
still lacks a securely acquired original and the Earth-science lead returned
HTML rather than a genuine technical PDF. Specialist search saturation,
uninspected supplementary code and checkpoint qualification remain open.
These coverage limits do not prevent the concrete prototype work, and they
do prevent a claim of an exhaustive or generally qualified temporal causal
agent.
\citep{exec15_latent_design}
The remaining step is to preserve these distinctions when a proposed graph
becomes a formal causal query. The next study shows why a valid calculation
can still answer the wrong population question, and why a failed criterion
must not become an unsupported claim of nonidentification.
\label{exec15:latent-temporal-end}


% SOURCE docs/release/2026-10-29/source/docs/writing-repair/2026-10-29/src/84c8b19e3681956e.tex CHAR 0
\section{From language suggestions to an identification certificate}
\label{exec14:llm-identification}

An analyst asks whether servicing a delinquent account will increase repayment.
An assistant may propose a graph, an adjustment set or an instrument, yet these
are different objects from an identified effect. The product must preserve the
question through each transformation: which customers, which action, at which
time, compared with what alternative, and under which observation law? A formal
engine can answer the resulting mathematical question conditional on its model.
It cannot establish the institutional truth of an assumption supplied by the
language model. This division of responsibility creates a useful role for an
assistant without treating its fluency as a certificate. \citep{exec14_ident_contract}

\subsection{Evidence from different tasks should remain distinguishable}

\emph{CLadder} provides formal causal question generation and different
verbalizations, helping separate calculation from familiar real-world
associations. Its authors report improvements from structured CausalCoT
prompting and retain the difficulty of turning language into the right formal
inputs. The inspected paper and repository are separate versioned objects; a
result reported for one dataset version is not automatically a qualification of
every released generator template. In particular, a selected repository template
places an instrumental-variable ratio inside an average-treatment-effect query.
That label needs its own structural and population assumptions.
\citep{v12_la_2312_04350,exec14_ident_packages}

\emph{CounterBench} makes another contribution: 1,000 balanced binary questions
about deterministic structural mechanisms, spanning basic, joint, nested and
conditional questions. In the inspected first version, Chen and colleagues report
GPT-4o accuracy of 52.5 percent with standard prompting, 78.8 percent with
CausalCoT and 92.0 percent with CoIn. The same table reports a decrease from
51.8 to 49.8 percent for Davinci. CoIn organizes extraction, iterative reasoning
and backtracking through demonstrations; the description of an algorithm does
not by itself make the model's generated steps a machine-checked proof. The
authors identify probabilistic mechanisms and hidden confounding as extensions.
These favorable and adverse results are historical paper observations, not
measurements of this product or of current services. The inspected v1 title is
\emph{A Benchmark for Counterfactuals Reasoning in Large Language Models}; it
should not be silently conflated with later title or release variants.
\citep{exec14_ident_counterbench}

\emph{METER} instead aligns three linguistic question families over 4,145 shared
contexts. Its first level concerns causal relations expressed in text, which the
paper calls discovery. That task differs from recovering a graph or equivalence
class from an observational law. Intervention and counterfactual options are
generated by a model and then edited and filtered by human annotators. For
GPT-4o the authors report zero-shot accuracies of 87.92, 77.96 and 67.42 percent,
with corresponding CoT accuracies of 86.46, 75.14 and 66.18 percent. Its useful
lessons concern shared context, distraction by irrelevant facts, causal reversal
and unsupported statements. Curated semantic labels do not supply a complete
formal identification oracle. Comparing its headline percentage directly with
CounterBench would confound the question, answer space, information and target.
\citep{exec14_ident_meter}

Gendron and colleagues separate graph extraction from evaluating counterfactual
mechanisms. Their experiments show that providing the correct graph need not
remove mechanism-evaluation errors. They also acknowledge using model-derived
counterfactual values as proxies for some real-world ground truth. This is a
reason to retain independent numerical or logical references when evaluating a
banking workflow. A coherent graph and a confident answer can fail at different
stages. \citep{v12_la_2410_06392}

\emph{CausalDS} moves toward an agent workflow with public scene files, separate
private truth, observation noise and identification-gated questions. It
distinguishes population-average intervention targets from counterfactual targets
and conceptual variables from their noisy measurements. Those distinctions are
directly relevant to a bank's data dictionary. However, benchmark architecture
and the semantics of every implementation branch require separate inspection.
The selected source retains unknown values for missing engines and exceptions,
while returning false with error metadata for invalid target nodes. A negative
criterion result can also survive an unavailable or failed fallback. A consumer
that reads only the Boolean can turn a malformed or unresolved request into a
claim of scientific impossibility. \citep{v12_la_2607_08093,exec14_ident_packages}

Formal proof assistance introduces a further boundary. \emph{CausalSmith}
distinguishes a mechanically checked proof from a faithful formalization of the
intended statement. A changed hypothesis, a vacuous definition or an assumed
hard step can yield an artifact that fails the original scientific purpose.
The authors' internal catalog reports 14 accepted, 59 downgraded and 71 failed
runs among 144 recorded runs; the identification and structural-model clusters
contain one acceptance among 94 runs. Reruns count separately. These are not
independent estimates of autonomous research quality: the paper acknowledges
the absence of independent human quality validation and an estimated error rate
for statement matching. A valid proof of the wrong banking target remains the
wrong deliverable. \citep{exec13_census_causalsmith}

\subsection{A valid instrument need not answer the population question}

Consider a stipulated servicing experiment. At time zero an eligible population
receives a randomized encouragement $Z$. Actual servicing $A$ occurs by day
seven, and binary repayment $Y$ is measured by day thirty. Let $A(z)$ denote
take-up under encouragement $z$, and $Y(a)$ repayment under servicing $a$.
Assume encouragement is independent of the take-up type and potential outcomes,
that it affects repayment only through servicing, and that consistency and no
interference hold. Take-up is monotone: there are no borrowers who accept only
when they are not encouraged. These assumptions define the example; an
assistant's proposed instrument does not establish them. \citep{exec14_ident_exact}

Suppose one quarter of borrowers are always-takers with $(A(0),A(1))=(1,1)$,
one quarter are never-takers with $(0,0)$, and one half are compliers with
$(0,1)$. The two candidate worlds have the same take-up and differ only in
potential repayments:
\begin{center}
\begin{tabular}{lccc}
\hline
Stratum & Probability & World A: $(Y(0),Y(1))$ & World B: $(Y(0),Y(1))$ \\
\hline
Always-taker & $1/4$ & $(0,1)$ & $(1,1)$ \\
Never-taker & $1/4$ & $(0,1)$ & $(0,0)$ \\
Complier & $1/2$ & $(0,1)$ & $(0,1)$ \\
\hline
\end{tabular}
\end{center}
Under either world the complete observed joint law of $(Z,A,Y)$ is the same.
For $Z=0$, a quarter have $(A,Y)=(1,1)$ and three quarters have $(0,0)$;
for $Z=1$, these probabilities are reversed. Consequently,
\[
 \frac{E[Y\mid Z=1]-E[Y\mid Z=0]}{E[A\mid Z=1]-E[A\mid Z=0]}
 =\frac{3/4-1/4}{3/4-1/4}=1.
\]
The population average effect is nevertheless 1 in world A and $1/2$ in world B.
Exact enumeration reproduces every cell and contrast; neither world is fitted
to bank data. \citep{exec14_ident_exact}

The reason is visible algebraically. Let $C$ denote the complier stratum.
Under the stated independence and exclusion restrictions, the first-stage
difference is $P(C)$: always-takers and never-takers do not change treatment.
For the same reason the reduced-form outcome difference is
$P(C)E[Y(1)-Y(0)\mid C]$. Dividing therefore identifies the complier effect
when $P(C)>0$. The ratio supplies no new observation of the always-takers'
untreated outcome or the never-takers' treated outcome. Here their observed
outcomes and the binary restriction permit a population-effect range
$[1/2,1]$. Varying the unobserved outcomes, or mixing compatible subtypes,
attains the endpoints and intermediate values in this stipulated class.
It would be incorrect to attach the ratio's value to the unrestricted
population target. \citep{exec14_ident_exact}

The distinction can reverse an investment conclusion. Stipulate a USD 100 value
for repayment and a USD 75 servicing cost per treated borrower. A hypothetical
policy treating everyone rather than nobody then has net value between minus
USD 25 and plus USD 25 per eligible borrower. The complier calculation alone
gives plus USD 25 and cannot settle that population policy. This comparison is
neither the value of the encouragement policy nor observed bank ROI; feasibility,
actual costs and transport beyond the stipulated population remain separate
questions. The exact example explains why the query's population must survive
every language, tool and reporting conversion. \citep{exec14_ident_exact}

\subsection{Suggestion packages and formal engines have different obligations}

The selected pinned PyWhyLLM paths illustrate a useful but limited component.
Its public initialization imports suggestion classes. The main identification
suggester delegates backdoor proposals to confounder suggestions, leaves a
frontdoor method unimplemented, and retains a commented estimand path. A
separate simple suggester prompts for IV, backdoor or frontdoor names and
extracts response tags. Protocol classes declare interfaces; the validation
suggester adds prompted critique. None of these selected paths supplies a
formal identification certificate. Its selected unit tests mock model responses
and compare extracted lists, establishing intended parsing behavior rather than
live scientific validity. \citep{exec14_ident_packages}

The proposed reuse is therefore a small optional elicitation component. Its
outputs should be candidate variables or assumptions with provenance, subject
to schema checks, variable allowlists, finite retries and review. A package name
or a method called validation should not grant scientific authority. The
root MIT notice has been inspected separately from complete runtime and
dependency rights. CLadder and CausalDS likewise need version-specific
generator, observation-model and scorer audits, with their code and data rights
kept distinct. CounterBench and METER currently support the literature argument;
adopting their released datasets requires item-level semantics and rights review.
No selected external integration is enabled by this analysis.
\citep{exec14_ident_packages,exec14_ident_contract}

The formal engine needs a stricter output contract. An identified result retains
the exact query, graph, observation law, assumptions, supported model class,
expression and checkable certificate. A not-identified result needs a compatible
counterexample or an applicable complete-method certificate for the same target
and class. Unknown covers insufficient assumptions, unsupported methods, invalid
input, absent engines, timeouts and unverified negative criterion results.
Execution status and scientific status should be distinct fields, with reasons
and fallback history. For example, failure to find a backdoor set may justify
trying another identification method; it does not by itself prove that no
observational formula exists. \citep{exec14_ident_contract}

This matters in the inspected CausalDS path. Filtering an IV candidate from a
population-ATE answer avoids one target substitution, but an empty remaining
method dictionary becomes false. When a more general fallback is absent or
fails, that result still needs its criterion scope. Invalid graph-node cases
also return false alongside error metadata. The proposed wrapper should inspect
all of those fields and preserve unknown or invalid-request semantics as
appropriate. This is a static source finding and a proposed conversion rule,
not an observed failure of an installed CausalDS runtime.
\citep{exec14_ident_packages,exec14_ident_contract}

\subsection{A public prototype that tests the actual weakness}

The smallest useful prototype combines independently enumerated mechanisms with
audited formal-engine outputs. It asks the assistant to translate a banking
question, expose its assumptions, select a supported query and explain the
result without changing its scope. Same-observed-law worlds test whether the
workflow recognizes missing information. Paraphrases and renamed variables test
whether it preserves meaning. A changed population or observation mask should
trigger a changed query, while harmless wording should not. Deliberately invalid
graphs, unsupported classes and timeouts test whether uncertainty survives.
Private truth artifacts must remain separate from assistant inputs.
\citep{exec14_ident_contract}

Report false identification, target substitution, invalid certificates,
unavailable results and the reasons for abstention separately. A high answer
accuracy on identified questions cannot conceal confident answers to
unidentified ones. Measure complete costs: formalization review, tool calls,
failed executions, proof checking and human correction as well as model tokens.
These measurements can use public and generated cases before an expensive bank
experiment. They establish whether a prototype merits institutional evaluation;
they do not manufacture assignment truth or stakeholder acceptance.
\citep{exec14_ident_contract}

The engineering route is concrete. Representative rates of business-to-formal
error, reliable detection of missing assumptions and broad faithful proof
translation remain research questions. A proof checker can reduce one class of
error while leaving those questions unresolved. The product proposal should
fund a scoped prototype with explicit failure cases and retain the unresolved
claims, rather than promise that formal syntax alone solves causal reasoning.
\citep{exec14_ident_contract}
\subsection{What the complete discovery argument supports}

Language assistance can supply a candidate definition, restriction, mechanism
or analysis. Its value depends on the information that proposal adds and the
consequences of admitting it. The chapter's distinct examples show why graph
agreement, response stability, statistical evidence, identification and
monetary value require separate assessments. Retaining a compatible model
class can be the scientifically useful result even when a single answer would
be easier to report.

A defensible proposal therefore preserves the target and clock, compares
against numerical and expert alternatives, and carries semantic failures and
uncertainty through to the decision. Source-level repairs can make a tool's
behavior more faithful to its contract; they cannot create an identifying
observation. Where the remaining alternatives imply different actions, the
next question is which feasible experiment could distinguish them at worthwhile
cost. That is the subject of the next chapter.
\label{exec14:llm-identification-end}

